% Generated by scripts/manuscript_map.py. Edit docs/manuscript-map.json.
\chapter{Manuscript index}\label{sec:manuscript-index}
Read this chapter in manuscript order. Each entry links the draft source and any corresponding Lean declarations. The coverage assessment refers to the full manuscript item. The component links lead to checked statements and their scope qualifications in this blueprint. For numbered results, the proof note compares the manuscript argument with the implemented route. Lean names are in the GaussianTesting namespace. Source links open line-numbered copies on the documentation site.

\section{Manuscript section 1. Introduction}
\subsection{Theorem 1.1}\label{paper:r1-1}
\textbf{Checked.} Gaussian tests have perfect completeness and one-sided soundness with O(epsilon\textasciicircum{}(-2) log(1/delta)) input copies, where epsilon is the trace-distance threshold. \bpsource{material/manuscript.tex}{126}{Manuscript, line 126.}

The public fermionic and bosonic testing theorems use the manuscript’s pointwise squared-overlap promise and its equivalent trace-distance promise for the independent state classes. They have perfect completeness, rejection probability at least 1-delta, the existing explicit measurement protocols, and integer copy counts independent of the mode count. General bosons use (k,C)=(4,48), fermions (2,8); the bosonic interface also covers mean-zero states (2,8) and coherent states (2,4). The exact count is k*ceil((C/epsilon²)*log(1/delta)), bounded uniformly by kC/epsilon²*log(1/delta)+k and by 2kC/epsilon²*log(1/delta) for delta<=exp(-1), retaining the manuscript source’s ceiling qualification.

\textbf{Proof route.} Apply Corollaries 5.3 and 6.8 to the equivalent supremal-fidelity promise. The pure-state trace-norm identity converts distance at least epsilon from every Gaussian state to the pointwise overlap condition. The same independent repeated measurements have the previously proved accepting readouts from Figure 1 and exact ceiling copy counts; the explicit rate constants are independent of the number of modes.

\bpdecls{\textbf{Lean.} \bpsource{GaussianTesting/Testing.lean}{58}{\texttt{fermionic\_gaussian\_testing}}, \bpsource{GaussianTesting/Testing.lean}{78}{\texttt{bosonic\_gaussian\_testing}}, \bpsource{GaussianTesting/Testing.lean}{98}{\texttt{fermionic\_gaussian\_testing\_traceDistance}}, \bpsource{GaussianTesting/Testing.lean}{119}{\texttt{bosonic\_gaussian\_testing\_traceDistance}}, \bpsource{GaussianTesting/Fidelity.lean}{58}{\texttt{pureFidelity\_sq\_le\_iff}}, \bpsource{GaussianTesting/Fermions/Testing.lean}{97}{\texttt{fermionicTestCopies\_le\_rate}}, \bpsource{GaussianTesting/RepeatedTest.lean}{114}{\texttt{testingCopies\_le\_rate}}, \bpsource{GaussianTesting/Fermions/BellMeasurement.lean}{113}{\texttt{fermionicBellMeasurement\_eq\_PF}}, \bpsource{GaussianTesting/Bosons/MeanZeroMeasurement.lean}{96}{\texttt{meanZeroNumberMeasurement\_eq\_P0}}, \bpsource{GaussianTesting/Bosons/CoherentMeasurement.lean}{114}{\texttt{coherentNumberMeasurement\_eq\_Q}}, \bpsource{GaussianTesting/Bosons/FourCopyMeasurement.lean}{105}{\texttt{fourCopyNumberMeasurement\_acceptance}}.}

\textbf{Theorem node.} \bplink{target:measurement-endpoints}{Testing Gaussianity with constant sample complexity}.

\subsection{Pure-state trace-distance promise}\label{paper:trace-distance}
\textbf{Checked.} For unit vectors, epsilon-farness in trace distance is equivalent to squared overlap at most 1-epsilon\textasciicircum{}2 against every target state. \bpsource{material/manuscript.tex}{134}{Manuscript, line 134.}

For unit vectors in an arbitrary complete complex Hilbert space, half the trace norm of the difference of their rank-one projections is sqrt(1-|inner|²). The norm is defined as the trace of CFC.abs(A)=sqrt(A†A); positivity, summability of its diagonal, and independence of the Hilbert basis are proved for the pure-state difference. The nonnegative trace-distance threshold is equivalent to the squared-overlap promise. This includes infinite-dimensional bosonic Hilbert spaces and coincident rays.

\textbf{Proof route.} For A=P\_psi-P\_phi and d=sqrt(1-|inner|²), rank-one algebra gives A³=d²A and trace(A²)=2d². If d>0, d\textasciicircum{}(-1)A² is positive and squares to A², so it equals |A| by uniqueness of the positive square root. If d=0, A=0. The summable diagonal of this explicit modulus has trace 2d in every Hilbert basis. Divide by two and square nonnegative thresholds to obtain the promise equivalence.

\bpdecls{\textbf{Lean.} \bpsource{GaussianTesting/PureStateDifference.lean}{30}{\texttt{pureStateDifference}}, \bpsource{GaussianTesting/TraceDistance.lean}{35}{\texttt{traceNorm}}, \bpsource{GaussianTesting/TraceDistance.lean}{146}{\texttt{pureStateTraceDistance}}, \bpsource{GaussianTesting/PureStateDifference.lean}{50}{\texttt{pureStateDifference\_sq}}, \bpsource{GaussianTesting/PureStateDifference.lean}{62}{\texttt{pureStateDifference\_cube}}, \bpsource{GaussianTesting/PureStateDifference.lean}{92}{\texttt{pureStateDifference\_fourth}}, \bpsource{GaussianTesting/PureStateDifference.lean}{99}{\texttt{pureStateDifference\_eq\_zero\_of\_overlap\_sq\_eq\_one}}, \bpsource{GaussianTesting/TraceDistance.lean}{77}{\texttt{abs\_pureStateDifference}}, \bpsource{GaussianTesting/TraceDistance.lean}{121}{\texttt{isPositiveTraceClass\_abs\_pureStateDifference}}, \bpsource{GaussianTesting/TraceDistance.lean}{128}{\texttt{traceNorm\_pureStateDifference}}, \bpsource{GaussianTesting/TraceDistance.lean}{151}{\texttt{pureStateTraceDistance\_eq\_sqrt}}, \bpsource{GaussianTesting/TraceDistance.lean}{157}{\texttt{pureStateTraceDistance\_eq\_half\_traceNorm}}, \bpsource{GaussianTesting/TraceDistance.lean}{166}{\texttt{le\_pureStateTraceDistance\_iff}}.}

\textbf{Blueprint.} \bplink{lem:pure-state-trace-distance}{Pure-state trace distance}, \bplink{target:measurement-endpoints}{Testing Gaussianity with constant sample complexity}.

\subsection{Parity of the difference output}\label{paper:difference-parity}
\textbf{Checked.} The difference register obtained from identical bosonic inputs has even parity, so its mean vanishes when first moments exist. \bpsource{material/manuscript.tex}{143}{Manuscript, line 143.}

For every pure input, inverse balanced rotation turns input swap into parity on the difference register. Every reduced-state Gram column is even, and the reduced operator has range in the even subspace. Its countable-basis position and Fourier-coordinate Gram densities are integrable, have mass equal to the reduced trace, and integrate over each measurable event to the actual bounded-event operator trace. Unit inputs give normalized densities. Global integrability of coordinate times the relevant density is the first-moment existence hypothesis; both position and physical momentum means then vanish. Global moment integrability also implies the individual column hypotheses. The complete number basis supplies the required countable basis on L²(R\textasciicircum{}m).

\textbf{Proof route.} The balanced swap/parity identity fixes the two-identical-copy output under difference-register parity. Right-column transport gives even Gram vectors, and their convergent Gram sum gives even range of the reduced density operator. Fourier transformation preserves even parity. The nonnegative squared-amplitude series defines the global position and Fourier-coordinate densities; Tonelli and Parseval give their masses and nonnegative domination of each column. Bounded position-indicator operators and their Fourier conjugates identify the densities with the reduced state’s outcome probabilities. Reflection invariance makes every integrable odd density-weighted observable have zero integral, including position and momentum (2 pi times Mathlib’s Fourier coordinate).

\bpdecls{\textbf{Lean.} \bpsource{GaussianTesting/Bosons/DifferenceParity.lean}{36}{\texttt{schrodingerDifference\_even}}, \bpsource{GaussianTesting/Bosons/DifferenceParity.lean}{52}{\texttt{schrodingerDifference\_column\_even}}, \bpsource{GaussianTesting/Bosons/DifferenceParity.lean}{65}{\texttt{schrodingerReducedState\_even}}, \bpsource{GaussianTesting/Bosons/DifferenceParity.lean}{128}{\texttt{waveParity\_fourier\_of\_even}}, \bpsource{GaussianTesting/Bosons/DifferenceParity.lean}{186}{\texttt{waveGramDensity}}, \bpsource{GaussianTesting/Bosons/DifferenceParity.lean}{225}{\texttt{waveGramDensity\_integrable}}, \bpsource{GaussianTesting/Bosons/DifferenceParity.lean}{232}{\texttt{waveGramDensity\_integral}}, \bpsource{GaussianTesting/Bosons/DifferenceParity.lean}{252}{\texttt{waveGramDensity\_integrable\_component\_moment}}, \bpsource{GaussianTesting/Bosons/DifferenceParity.lean}{291}{\texttt{schrodingerDifferenceDensity}}, \bpsource{GaussianTesting/Bosons/DifferenceParity.lean}{298}{\texttt{schrodingerDifferenceFourierDensity}}, \bpsource{GaussianTesting/Bosons/DifferenceParity.lean}{304}{\texttt{schrodingerDifferenceDensity\_total}}, \bpsource{GaussianTesting/Bosons/DifferenceParity.lean}{476}{\texttt{schrodingerDifferenceDensity\_normalized}}, \bpsource{GaussianTesting/Bosons/DifferenceParity.lean}{338}{\texttt{wavePositionEvent}}, \bpsource{GaussianTesting/Bosons/DifferenceParity.lean}{445}{\texttt{waveFourierEvent}}, \bpsource{GaussianTesting/Bosons/DifferenceParity.lean}{436}{\texttt{schrodingerDifferenceDensity\_positionEvent}}, \bpsource{GaussianTesting/Bosons/DifferenceParity.lean}{468}{\texttt{schrodingerDifferenceFourierDensity\_positionEvent}}, \bpsource{GaussianTesting/Bosons/DifferenceParity.lean}{324}{\texttt{schrodingerReducedState\_globalFirstMoments}}, \bpsource{GaussianTesting/Bosons/DifferenceParity.lean}{486}{\texttt{schrodingerReducedState\_firstMoments\_of\_global}}.}

\textbf{Blueprint.} \bplink{target:fig1-general}{Figure 1(d), four-copy measurement}, \bplink{lem:difference-parity}{Parity and mean of the difference output}.

\subsection{Theorem 1.2}\label{paper:r1-2}
\textbf{Checked.} Informal robustness theorem. For a nonempty certainty set, a swap-commuting two-copy projection and three-copy gap at 0 <= alpha < 2/3 give the linear fidelity consequence for every tau >= 0. \bpsource{material/manuscript.tex}{159}{Manuscript, line 159.}

For any complex Hilbert space with the stated three-copy spectral inclusion and nonempty certainty set, every admissible error gives the claimed lower bound on squared fidelity. The proof uses the same certainty-state witness as Theorem 3.4.

\textbf{Proof route.} Apply the linear conclusion of Theorem 3.4 to the certainty set. Squared fidelity is at least the squared overlap of its witness.

\bpdecls{\textbf{Lean.} \bpsource{GaussianTesting/GradientFlow/HilbertSpectralGap.lean}{72}{\texttt{testing\_fidelity\_of\_threeCopy\_spectrum}}, \bpsource{GaussianTesting/GradientFlow/SharpRobustness.lean}{69}{\texttt{RobustnessBounds.linear\_fidelity}}, \bpsource{GaussianTesting/GradientFlow/SharpRobustness.lean}{95}{\texttt{testing\_fidelity\_of\_gap}}, \bpsource{GaussianTesting/GradientFlow/SharpRobustness.lean}{119}{\texttt{TensorSquare.testing\_fidelity\_abstract}}.}

\textbf{Theorem node.} \bplink{thm:gradient-flow-intro}{Linear fidelity bound from the spectral gap}.

\subsection{Figure 1(a)}\label{paper:f1a}
\textbf{Checked.} Apply a two-copy Bell measurement and accept precisely the outcomes compatible with Lambda=0. \bpsource{material/manuscript.tex}{187}{Manuscript, line 187.}

The explicit product Bell vectors form an orthonormal basis for every mode count, including zero. Their unitary decoder has classical phase and parity bitstrings a,b. The integer score is sum\_k (product\_\{l<k\} (-1)\textasciicircum{}b\_l) (-1)\textasciicircum{}a\_k (1-(-1)\textasciicircum{}b\_k), including the Jordan-Wigner prefix parity. Accept exactly score zero. Pulling back this diagonal outcome projection through the Bell decoder gives PF on the full two-copy Hilbert space. The accepted outcome probability equals the PF expectation and hence p(psi) on two identical copies.

\textbf{Proof route.} Implement the stated Bell measurement in the standard product Bell basis. Compute the eigenvalues of each c\_j tensor c\_j from the local Bell signs and the Jordan-Wigner strings, then sum to obtain Lambda’s integer score. The kernel projection fixes precisely the zero-score Bell vectors and annihilates the others. Extending the identity over the orthonormal basis proves U\_Bell† D\_A U\_Bell=PF; summing the accepted coordinate squared norms gives the Born probability.

\bpdecls{\textbf{Lean.} \bpsource{GaussianTesting/Fermions/BellBasis.lean}{40}{\texttt{bellLocalAmplitude}}, \bpsource{GaussianTesting/Fermions/BellBasis.lean}{56}{\texttt{fermionicBellVector}}, \bpsource{GaussianTesting/Fermions/BellBasis.lean}{108}{\texttt{fermionicBellVector\_orthonormal}}, \bpsource{GaussianTesting/Fermions/BellBasis.lean}{114}{\texttt{fermionicBellBasis}}, \bpsource{GaussianTesting/Fermions/BellBasis.lean}{124}{\texttt{fermionicBellReadout}}, \bpsource{GaussianTesting/Fermions/BellBasis.lean}{137}{\texttt{fermionicBellReadout\_apply}}, \bpsource{GaussianTesting/Fermions/BellEigenvalues.lean}{82}{\texttt{fermionicBellMajoranaEigenvalue}}, \bpsource{GaussianTesting/Fermions/BellEigenvalues.lean}{120}{\texttt{fermionicBellEigenvalue}}, \bpsource{GaussianTesting/Fermions/BellEigenvalues.lean}{125}{\texttt{fermionicBellAccepts}}, \bpsource{GaussianTesting/Fermions/BellEigenvalues.lean}{96}{\texttt{majoranaPair\_fermionicBellVector}}, \bpsource{GaussianTesting/Fermions/BellEigenvalues.lean}{145}{\texttt{lambdaF\_fermionicBellBasis}}, \bpsource{GaussianTesting/Fermions/BellMeasurement.lean}{78}{\texttt{fermionicBellOutcomeProjection}}, \bpsource{GaussianTesting/Fermions/BellMeasurement.lean}{89}{\texttt{fermionicBellMeasurement}}, \bpsource{GaussianTesting/Fermions/BellMeasurement.lean}{113}{\texttt{fermionicBellMeasurement\_eq\_PF}}, \bpsource{GaussianTesting/Fermions/BellMeasurement.lean}{122}{\texttt{fermionicBellReadout\_probability}}, \bpsource{GaussianTesting/Fermions/BellMeasurement.lean}{135}{\texttt{fermionicBellMeasurement\_acceptance}}.}

\textbf{Blueprint.} \bplink{target:fig1-fermions}{Figure 1(a), fermionic Bell measurement}, \bplink{target:measurement-endpoints}{Testing Gaussianity with constant sample complexity}.

\subsection{Figure 1(b)}\label{paper:f1b}
\textbf{Checked.} Phase shift every mode of register two by pi/2, apply the positive-angle balanced splitter, and accept equal total particle counts on the two registers. \bpsource{material/manuscript.tex}{211}{Manuscript, line 211.}

The physical circuit phases register two by (-i) to its total occupation, then applies the positive balanced splitter. Pulling back the equal-total-count outcome gives exactly schrodingerTensorP0 on the full physical two-copy space; its Born acceptance is p0. The number-difference/negative-Lambda identity is checked on the finite-particle core, and transformed number eigenvectors plus density identify the bounded projections. No separate unbounded number-operator closure is assumed.

\bpdecls{\textbf{Lean.} \bpsource{GaussianTesting/Bosons/MeanZeroMeasurement.lean}{40}{\texttt{waveSecondRegisterPhase\_numberCoordinates}}, \bpsource{GaussianTesting/Bosons/EqualNumberProjection.lean}{80}{\texttt{equalNumberInverseSub\_numberDifference}}, \bpsource{GaussianTesting/Bosons/EqualNumberProjection.lean}{226}{\texttt{equalNumberCircuit\_conj\_projection}}, \bpsource{GaussianTesting/Bosons/MeanZeroMeasurement.lean}{55}{\texttt{waveEqualTotalNumberProjection}}, \bpsource{GaussianTesting/Bosons/MeanZeroMeasurement.lean}{96}{\texttt{meanZeroNumberMeasurement\_eq\_P0}}, \bpsource{GaussianTesting/Bosons/MeanZeroMeasurement.lean}{107}{\texttt{meanZeroNumberMeasurement\_acceptance}}, \bpsource{GaussianTesting/Bosons/Main.lean}{359}{\texttt{P0B}}, \bpsource{GaussianTesting/Bosons/BeamSplitter.lean}{548}{\texttt{ULp}}.}

\textbf{Blueprint.} \bplink{target:fig1-mean-zero}{Figure 1(b), equal-number measurement}, \bplink{target:fig1-general}{Figure 1(d), four-copy measurement}.

\subsection{Figure 1(c)}\label{paper:f1c}
\textbf{Checked.} Apply the negative-angle balanced splitter and test register two for vacuum. \bpsource{material/manuscript.tex}{249}{Manuscript, line 249.}

The total-number readout is defined in the proved physical oscillator basis. Zero total occupation is exactly the vacuum outcome for every mode count. The circuit applying the inverse balanced splitter and reading zero total number in register two has accepting observable exactly schrodingerQ, and its Born probability on two identical inputs is schrodingerCoherentAcceptance.

\bpdecls{\textbf{Lean.} \bpsource{GaussianTesting/Bosons/CoherentMeasurement.lean}{66}{\texttt{waveTotalNumberProjection}}, \bpsource{GaussianTesting/Bosons/CoherentMeasurement.lean}{89}{\texttt{waveTotalNumberProjection\_zero}}, \bpsource{GaussianTesting/Bosons/CoherentMeasurement.lean}{107}{\texttt{coherentNumberMeasurement}}, \bpsource{GaussianTesting/Bosons/CoherentMeasurement.lean}{114}{\texttt{coherentNumberMeasurement\_eq\_Q}}, \bpsource{GaussianTesting/Bosons/CoherentMeasurement.lean}{121}{\texttt{coherentNumberMeasurement\_acceptance}}, \bpsource{GaussianTesting/Bosons/BeamSplitter.lean}{851}{\texttt{QB\_eq\_conj}}, \bpsource{GaussianTesting/Bosons/BeamSplitter.lean}{548}{\texttt{ULp}}, \bpsource{GaussianTesting/Bosons/BeamSplitter.lean}{556}{\texttt{VLp}}, \bpsource{GaussianTesting/Bosons/BeamSplitter.lean}{662}{\texttt{vacProj}}, \bpsource{GaussianTesting/Bosons/NumberCoherent.lean}{102}{\texttt{waveNumberCoordinates\_conj\_Q}}.}

\textbf{Blueprint.} \bplink{lem:coherent-projection}{Coherent accepting projection}, \bplink{target:fig1-coherent}{Figure 1(c), vacuum measurement}.

\subsection{Figure 1(d)}\label{paper:f1d}
\textbf{Checked.} Apply negative-angle balanced splitters to wires (1,2) and (3,4). Keep wires (2,4) and use the measurement from panel (b). \bpsource{material/manuscript.tex}{273}{Manuscript, line 273.}

The physical circuit applies U\_\{-pi/4\} to both pairs, discards outputs (1,3), and applies the phase/splitter/equal-count measurement from panel (b) to retained outputs (2,4). Its accepting operator equals schrodingerFourCopyProjection on the full four-copy Hilbert space. Consequently its Born probability on four identical inputs is schrodingerFourCopyAcceptance. The preparation equality and operator equality follow directly from the displayed wiring.

\bpdecls{\textbf{Lean.} \bpsource{GaussianTesting/Bosons/FourCopyMeasurement.lean}{37}{\texttt{fourCopyFigurePreparation\_explicit}}, \bpsource{GaussianTesting/Bosons/FourCopyMeasurement.lean}{44}{\texttt{fourCopyFigurePreparation\_eq}}, \bpsource{GaussianTesting/Bosons/FourCopyMeasurement.lean}{65}{\texttt{fourCopyFigureProjection\_eq}}, \bpsource{GaussianTesting/Bosons/FourCopyMeasurement.lean}{83}{\texttt{fourCopyNumberMeasurement}}, \bpsource{GaussianTesting/Bosons/FourCopyMeasurement.lean}{99}{\texttt{fourCopyNumberMeasurement\_eq\_projection}}, \bpsource{GaussianTesting/Bosons/FourCopyMeasurement.lean}{105}{\texttt{fourCopyNumberMeasurement\_acceptance}}.}

\textbf{Blueprint.} \bplink{target:fig1-general}{Figure 1(d), four-copy measurement}.

\subsection{Gaussian states and first and second moments}\label{paper:gaussian-moments}
\textbf{Open.} Within the Gaussian class, the first and second creation/annihilation moments determine the state. \bpsource{material/manuscript.tex}{323}{Manuscript, line 323.}

The concrete orbit and covariance exactness theorems do not provide an identification of all physical Gaussian states by these moments.

\textbf{Blueprint.} \bplink{prelim:bosons}{Bosonic Gaussian preliminaries}.

\subsection{Covariance characterization}\label{paper:covariance-criterion}
\textbf{Partial coverage.} Gaussianity corresponds to all bosonic symplectic covariance eigenvalues or all fermionic covariance singular values being one. The bosonic criterion uses finite second moments and the stated vacuum normalization. \bpsource{material/manuscript.tex}{326}{Manuscript, line 326.}

The fermionic covariance saturation calculation supports exactness. The full physical covariance criteria, including the finite-second-moment bosonic statement, require separate identifications.

\bpdecls{\textbf{Lean.} \bpsource{GaussianTesting/Fermions/CovarianceBasic.lean}{23}{\texttt{majMoment}}, \bpsource{GaussianTesting/Fermions/CovarianceBasic.lean}{69}{\texttt{norm\_lambdaF\_tp2}}, \bpsource{GaussianTesting/Fermions/GaussianCertainty.lean}{40}{\texttt{fermionicGaussian\_of\_lambdaF\_eq\_zero}}.}

\textbf{Blueprint.} \bplink{lem:exact-fermions}{Fermionic states accepted with certainty}.

\subsection{Orthogonal copy symmetries}\label{paper:copy-symmetry}
\textbf{Open.} The manuscript describes the orthogonal-group commutant for the restricted Gaussian actions and the restriction to transformations fixing the all-ones vector when displacements are allowed. \bpsource{material/manuscript.tex}{347}{Manuscript, line 347.}

The introduction and the new Section 5 and Section 6 discussions state the full orthogonal-copy commutant dualities. The fermionic testing proofs use the checked SO(2) and SO(3) second-quantized copy actions, without the general O(t) commutant theorem. The bosonic O(t) inverse-pullback action is constructed on copy-first coordinates Fin t × Fin m, matching R tensor I\_m; its full commutant identification is not formalized. These general duality assertions remain unused background. In the infinite-dimensional bosonic case, generation is understood in the von Neumann algebra sense.

\textbf{Blueprint.} \bplink{prelim:fermions}{Fermionic Gaussian preliminaries}, \bplink{prelim:bosons}{Bosonic Gaussian preliminaries}.

\subsection{Closed span and two-copy optimality}\label{paper:span-optimality}
\textbf{Checked.} Gaussian tensor squares have closed linear span equal to the relevant kernel, intersected with the symmetric subspace in the bosonic case. This yields an optimal perfectly complete two-copy test on identical inputs. \bpsource{material/manuscript.tex}{353}{Manuscript, line 353.}

For every mode count including zero, the fermionic Gaussian-square span equals the full kernel and PF has symmetric support. The bosonic closed span of independent normalized centered Gaussian squares equals the symmetric part of the actual closed generator kernel on full L2, without a moment hypothesis on kernel vectors. On any complex Hilbert space, a positive contraction accepting every vector in a set with certainty dominates the projection onto its closed span. This proves fermionic minimal acceptance on arbitrary two-copy vectors. For bosons, the actual P0 commutes with swap and agrees with the span projection on identical inputs, giving the manuscript optimality statement against every perfectly complete positive contraction, with no swap assumption on the competing effect. The manuscript wording is retained in its identical two-copy-input context. No span hypothesis remains in either public optimality theorem.

\textbf{Proof route.} For fermions, both Gaussian parity orbits span a subspace invariant under Majorana pairs and simultaneous Gaussian rotations. The joint Bell eigenspaces are one-dimensional. A nonzero vacuum-square overlap and signed coordinate permutations supply every zero-eigenvalue Bell line. The manuscript specifies full-orthogonal-group twirling in the cited argument, including both parity orbits; Lean retains the agreed direct spanning proof. For bosons, identify the closed generator kernel with invariance under the actual copy-rotation group. Average polynomial Gaussian wavefunctions over a full rotation period and then over copy swap. Each rotation orbit has a finite-dimensional polynomial span, so its Bochner average remains polynomial Gaussian. Injectivity of polynomial Gaussian realization turns the fixed-vector statements into polynomial identities. The substitutions u=x+iy and v=x-iy identify the generator with the difference of the two Euler operators; physical swap together with a quarter turn becomes interchange of u and v. A finite Fischer-duality argument proves that balanced, swap-invariant polynomials are generated by u\_i v\_j+u\_j v\_i: an annihilating functional gives a factorial-weighted dual polynomial killed by every symmetric mixed derivative, and the checked joint-kernel calculation forces it to vanish. These generators become twice x\_i x\_j+y\_i y\_j. Identical quadratic phase multipliers preserve centered Gaussian states and hence their closed square span. Differentiation of matrix coefficients under the integral, dominated by an integrable product of two L² functions, makes that span stable under multiplication by each Gram quadratic. The vacuum square and algebra induction therefore put every Gram-polynomial Gaussian in the span. Density of polynomial Gaussians and continuity of the rotation/swap average give every symmetric kernel vector. The reverse inclusion follows from Proposition 6.1 and symmetry of identical squares. Finally, positivity and perfect completeness force an effect to fix the closed square span. On identical inputs its projection agrees with the actual P0, because P0 commutes with swap; this gives the bosonic acceptance comparison.

\bpdecls{\textbf{Lean.} \bpsource{GaussianTesting/Fermions/GaussianSquareSpan.lean}{426}{\texttt{fermionicGaussianSquareSpan\_eq\_kernel}}, \bpsource{GaussianTesting/Fermions/GaussianSquareSpan.lean}{458}{\texttt{swap2\_PF}}, \bpsource{GaussianTesting/PerfectCompleteness.lean}{32}{\texttt{effect\_fixes\_of\_acceptance\_eq}}, \bpsource{GaussianTesting/PerfectCompleteness.lean}{47}{\texttt{effect\_fixes\_closedSpan}}, \bpsource{GaussianTesting/PerfectCompleteness.lean}{64}{\texttt{starProjection\_le\_of\_fixes}}, \bpsource{GaussianTesting/PerfectCompleteness.lean}{86}{\texttt{closedSpan\_projection\_le\_effect}}, \bpsource{GaussianTesting/PerfectCompleteness.lean}{96}{\texttt{closedSpan\_acceptance\_le}}, \bpsource{GaussianTesting/Fermions/PerfectCompleteness.lean}{24}{\texttt{fermionicTest\_projection\_le\_effect}}, \bpsource{GaussianTesting/Fermions/PerfectCompleteness.lean}{41}{\texttt{fermionicTest\_acceptance\_le}}, \bpsource{GaussianTesting/Bosons/GaussianSquareDensity.lean}{35}{\texttt{meanZeroGaussianSquareSpan}}, \bpsource{GaussianTesting/Bosons/GaussianSquareSpan.lean}{64}{\texttt{schrodingerSymmetricSpace}}, \bpsource{GaussianTesting/Bosons/GaussianSquareSpan.lean}{148}{\texttt{meanZeroGaussianSquareSpan\_eq\_symmetric\_kernel}}, \bpsource{GaussianTesting/Bosons/GaussianSquareSpan.lean}{169}{\texttt{mem\_meanZeroGaussianSquareSpan\_iff}}, \bpsource{GaussianTesting/Bosons/GaussianSquareDensity.lean}{212}{\texttt{polynomialGaussian\_mem\_meanZeroGaussianSquareSpan}}, \bpsource{GaussianTesting/Bosons/CopyPolynomialInvariants.lean}{21}{\texttt{polynomial\_mem\_copyGramAlgebra\_of\_invariant}}, \bpsource{GaussianTesting/Bosons/BalancedPolynomialGeneration.lean}{316}{\texttt{mem\_symmetricMixedAlgebra\_of\_balanced\_swap}}, \bpsource{GaussianTesting/Bosons/BalancedQuadratic.lean}{244}{\texttt{symmetricMixedLowering\_jointKernel\_eq\_zero}}, \bpsource{GaussianTesting/Bosons/PolynomialGaussianInjective.lean}{24}{\texttt{polynomialGaussian\_injective}}, \bpsource{GaussianTesting/Bosons/PerfectCompleteness.lean}{138}{\texttt{meanZeroGaussianSquareSpan\_projection\_square}}, \bpsource{GaussianTesting/Bosons/PerfectCompleteness.lean}{147}{\texttt{bosonicMeanZeroGaussian\_optimality}}.}

\textbf{Blueprint.} \bplink{prelim:fermions}{Fermionic Gaussian preliminaries}, \bplink{target:span-optimality}{Gaussian-square span and optimality}.

\section{Manuscript section 2. Preliminaries}
\subsection{Density operators and pure states}\label{paper:density}
\textbf{Checked.} Density operators are positive trace-class operators of trace one. Pure states are rank-one projectors represented by unit vectors modulo phase. \bpsource{material/manuscript.tex}{398}{Manuscript, line 398.}

On any complete complex Hilbert space, positive trace class is defined by positivity and summability of the diagonal in a Hilbert basis. The criterion, positive trace, and complex trace against every bounded observable are proved independent of the basis. Every density operator has a normalized summable Gram representation, obtained from the positive square root. Tensor products and their physical wavefunction realization preserve trace one. A vector defines a rank-one density operator exactly when it has norm one. Two unit-vector projectors coincide exactly when their vectors differ by a unit phase. A density operator is idempotent exactly when it is a unit-vector projector; every such density operator comes from a unique projective ray. A density operator of norm one is rank one, by its compact Gram top eigenvector and the zero residual trace. A unit bipartite vector with rank-one reduced state has a unit tensor factorization, by the observable identity and rank-one projection.

\textbf{Proof route.} The positive rank-one trace is the squared vector norm by Parseval. Equality of unit-vector projectors forces one vector to be a scalar multiple of the other, and normalization forces unit modulus. The checked norm-one density characterization and the C*-norm identity identify idempotent density operators with rank-one projectors. Phase independence then gives an injective ray-projector map and a unique ray for every pure density operator. These supporting density results are collected in the bosonic preliminaries.

\bpdecls{\textbf{Lean.} \bpsource{GaussianTesting/Bosons/HilbertDensity.lean}{47}{\texttt{HilbertDensity.IsPositiveTraceClass}}, \bpsource{GaussianTesting/Bosons/HilbertDensity.lean}{52}{\texttt{HilbertDensity.IsDensityOperator}}, \bpsource{GaussianTesting/Bosons/HilbertDensity.lean}{117}{\texttt{HilbertDensity.isPositiveTraceClass\_basis\_iff}}, \bpsource{GaussianTesting/Bosons/HilbertDensity.lean}{127}{\texttt{HilbertDensity.isDensityOperator\_basis\_iff}}, \bpsource{GaussianTesting/Bosons/HilbertDensity.lean}{204}{\texttt{HilbertDensity.IsDensityOperator.exists\_normalized\_gram\_representation}}, \bpsource{GaussianTesting/Bosons/HilbertDensity.lean}{238}{\texttt{HilbertDensity.IsPositiveTraceClass.operatorTrace\_basis}}, \bpsource{GaussianTesting/Bosons/HilbertDensity.lean}{216}{\texttt{HilbertDensity.operatorTrace\_mul\_gramOp}}, \bpsource{GaussianTesting/Bosons/HilbertDensity.lean}{244}{\texttt{HilbertDensity.IsPositiveTraceClass.operatorTrace\_mul\_comm}}, \bpsource{GaussianTesting/Bosons/SignedSwapPurity.lean}{122}{\texttt{HilbertDensity.IsDensityOperator.embeddedTensor}}, \bpsource{GaussianTesting/Bosons/Density.lean}{131}{\texttt{isDensityOperator\_iff\_exists\_gram\_representation}}, \bpsource{GaussianTesting/Bosons/DensityPure.lean}{72}{\texttt{HilbertDensity.IsDensityOperator.exists\_rankOne\_of\_norm\_eq\_one}}, \bpsource{GaussianTesting/Bosons/PureFactorization.lean}{34}{\texttt{HilbertDensity.exists\_unit\_tmul\_of\_reducedRight\_eq\_rankOne}}, \bpsource{GaussianTesting/Bosons/PureStateDensity.lean}{33}{\texttt{HilbertDensity.positiveTrace\_rankOne}}, \bpsource{GaussianTesting/Bosons/PureStateDensity.lean}{45}{\texttt{HilbertDensity.isPositiveTraceClass\_rankOne}}, \bpsource{GaussianTesting/Bosons/PureStateDensity.lean}{57}{\texttt{HilbertDensity.isDensityOperator\_rankOne\_iff}}, \bpsource{GaussianTesting/Bosons/PureStateDensity.lean}{70}{\texttt{HilbertDensity.rankOne\_eq\_iff\_phase}}, \bpsource{GaussianTesting/Bosons/PureStateDensity.lean}{89}{\texttt{HilbertDensity.IsDensityOperator.norm\_eq\_one\_iff\_rankOne}}, \bpsource{GaussianTesting/Bosons/PureStateDensity.lean}{98}{\texttt{HilbertDensity.IsDensityOperator.idempotent\_iff\_rankOne}}, \bpsource{GaussianTesting/Bosons/PureStateDensity.lean}{117}{\texttt{HilbertDensity.rayProjector}}, \bpsource{GaussianTesting/Bosons/PureStateDensity.lean}{120}{\texttt{HilbertDensity.isDensityOperator\_rayProjector}}, \bpsource{GaussianTesting/Bosons/PureStateDensity.lean}{125}{\texttt{HilbertDensity.rayProjector\_mk}}, \bpsource{GaussianTesting/Bosons/PureStateDensity.lean}{133}{\texttt{HilbertDensity.rayProjector\_injective}}, \bpsource{GaussianTesting/Bosons/PureStateDensity.lean}{142}{\texttt{HilbertDensity.IsDensityOperator.existsUnique\_rayProjector}}.}

\textbf{Blueprint.} \bplink{prelim:bosons}{Bosonic Gaussian preliminaries}, \bplink{lem:purity}{Purity bound for a density operator}.

\section{Manuscript section 3. Robustness}
\subsection{Acceptance probability (Eq. 3.1)}\label{paper:acceptance}
\textbf{Checked.} A swap-commuting orthogonal projection on two copies accepts with its Born expectation in the two-copy product state. \bpsource{material/manuscript.tex}{415}{Manuscript, line 415.}

The completed tensor square realizes the two-copy state and the acceptance expectation for any swap-commuting orthogonal projection. Its norm is at most one. The implementations of particular tests by physical measurements are tracked under Figure 1.

\bpdecls{\textbf{Lean.} \bpsource{GaussianTesting/GradientFlow/HilbertTensor.lean}{29}{\texttt{HilbertSquare}}, \bpsource{GaussianTesting/GradientFlow/HilbertTensor.lean}{147}{\texttt{hilbertTensorSquare}}, \bpsource{GaussianTesting/GradientFlow/Abstract.lean}{80}{\texttt{TensorSquare.IsAcceptance.norm\_le\_one}}, \bpsource{GaussianTesting/GradientFlow/Defs.lean}{94}{\texttt{IsAcceptanceProjection}}, \bpsource{GaussianTesting/GradientFlow/Defs.lean}{100}{\texttt{accProb}}, \bpsource{GaussianTesting/GradientFlow/Abstract.lean}{73}{\texttt{TensorSquare.IsAcceptance}}, \bpsource{GaussianTesting/GradientFlow/Abstract.lean}{86}{\texttt{TensorSquare.accProbA}}.}

\textbf{Blueprint.} \bplink{prelim:flow}{Hilbert-space and projective preliminaries}, \bplink{def:acceptance}{Acceptance projection}.

\subsection{Certainty set and supremum fidelity (Eq. 3.2)}\label{paper:fidelity}
\textbf{Checked.} The good set is the nonempty certainty set. Fidelity with psi is the supremum of the absolute overlap over states in that set. \bpsource{material/manuscript.tex}{419}{Manuscript, line 419.}

For unit-vector representatives, certaintyStates p selects unit vectors with p = 1, and pureFidelity is the supremum of the absolute overlaps. For a nonempty class of unit vectors and a unit input it lies in [0,1], dominates every witness overlap, and is unchanged by an input phase.

\bpdecls{\textbf{Lean.} \bpsource{GaussianTesting/Fidelity.lean}{17}{\texttt{certaintyStates}}, \bpsource{GaussianTesting/Fidelity.lean}{25}{\texttt{pureFidelity}}, \bpsource{GaussianTesting/Fidelity.lean}{35}{\texttt{overlap\_le\_pureFidelity}}, \bpsource{GaussianTesting/Fidelity.lean}{51}{\texttt{sq\_overlap\_le\_pureFidelity\_sq}}, \bpsource{GaussianTesting/Fidelity.lean}{39}{\texttt{pureFidelity\_nonneg}}, \bpsource{GaussianTesting/Fidelity.lean}{44}{\texttt{pureFidelity\_le\_one}}, \bpsource{GaussianTesting/Fidelity.lean}{76}{\texttt{pureFidelity\_phase}}, \bpsource{GaussianTesting/Fidelity.lean}{58}{\texttt{pureFidelity\_sq\_le\_iff}}.}

\textbf{Blueprint.} \bplink{prelim:flow}{Hilbert-space and projective preliminaries}, \bplink{def:fidelity}{Certainty set and fidelity}, \bplink{thm:bound-fidelity-gradient-flow}{Sharp angular and linear bounds, including Eq. (3.7)}, \bplink{thm:main-fermions}{Fidelity bound for the fermionic test}, \bplink{thm:main-boson-mean-zero}{Fidelity bound for mean-zero Gaussian states}, \bplink{thm:main-boson-coherent}{Fidelity bound for coherent states}.

\subsection{Gradient flow (Eq. 3.3)}\label{paper:flow}
\textbf{Checked.} Follow the gradient of the acceptance probability from any pure initial state on nonnegative time. \bpsource{material/manuscript.tex}{429}{Manuscript, line 429.}

In every affine chart containing the current ray, its coordinates solve the ODE obtained by applying the chart derivative to the Fubini–Study gradient in Eq. (3.4) at a normalized representative. Every unit-vector solution projects to a ray solution. Smoothness and uniqueness hold for the chart-defined ray trajectories themselves.

\bpdecls{\textbf{Lean.} \bpsource{GaussianTesting/GradientFlow/ProjectiveFlow.lean}{47}{\texttt{TensorSquare.IsProjectiveFlowOn}}, \bpsource{GaussianTesting/GradientFlow/ProjectiveFlow.lean}{30}{\texttt{TensorSquare.chartGrad}}, \bpsource{GaussianTesting/GradientFlow/ProjectiveFlow.lean}{72}{\texttt{TensorSquare.IsUnitFlowOn.projective}}, \bpsource{GaussianTesting/GradientFlow/ProjectiveFlow.lean}{190}{\texttt{TensorSquare.IsProjectiveFlowOn.eqOn\_Icc}}, \bpsource{GaussianTesting/GradientFlow/ProjectiveFlow.lean}{114}{\texttt{TensorSquare.IsProjectiveFlowOn.contDiffWithinAt\_chart\_Ici}}, \bpsource{GaussianTesting/GradientFlow/ProjectiveGradientFlow.lean}{86}{\texttt{projective\_gradient\_flow}}.}

\textbf{Blueprint.} \bplink{lem:flow-existence}{Existence and uniqueness of the global gradient flow}.

\subsection{Fubini-Study distance}\label{paper:projective-angle}
\textbf{Checked.} The projective distance between pure states equals arccos of the absolute overlap. \bpsource{material/manuscript.tex}{435}{Manuscript, line 435.}

The distance on complex rays equals arccos of the absolute inner product for any unit representatives. It induces the quotient topology on the projectivization, and the metric is complete when the Hilbert space is complete. At a unit representative the tangent norm induced by this distance is the Hilbert norm on its complex-orthogonal complement.

\bpdecls{\textbf{Lean.} \bpsource{GaussianTesting/GradientFlow/FubiniStudy.lean}{52}{\texttt{rayDist}}, \bpsource{GaussianTesting/GradientFlow/FubiniStudy.lean}{56}{\texttt{rayDist\_mk}}, \bpsource{GaussianTesting/GradientFlow/FubiniStudy.lean}{189}{\texttt{rayMetricSpace}}, \bpsource{GaussianTesting/GradientFlow/FubiniStudy.lean}{197}{\texttt{rayCompleteSpace}}, \bpsource{GaussianTesting/GradientFlow/ProjectiveSpace.lean}{144}{\texttt{rayChartHomeomorph}}, \bpsource{GaussianTesting/GradientFlow/ProjectiveSpace.lean}{169}{\texttt{contDiffOn\_rayTransition}}, \bpsource{GaussianTesting/GradientFlow/Angle.lean}{59}{\texttt{projective\_angle\_le\_of\_chord\_bound}}, \bpsource{GaussianTesting/GradientFlow/RayTangentChart.lean}{56}{\texttt{centeredRayChart}}, \bpsource{GaussianTesting/GradientFlow/RayTangentMetric.lean}{101}{\texttt{tendsto\_dist\_div\_time\_rayOfTangent}}, \bpsource{GaussianTesting/GradientFlow/RayTangentMetric.lean}{134}{\texttt{tendsto\_rayTangentMetric\_of\_dist}}, \bpsource{GaussianTesting/GradientFlow/RayTangentMetric.lean}{43}{\texttt{rayTangentMetric\_eq\_polarization}}.}

\textbf{Blueprint.} \bplink{prelim:flow}{Hilbert-space and projective preliminaries}.

\subsection{Acceptance gradient (Eq. 3.4)}\label{paper:gradient}
\textbf{Checked.} The gradient is four times the contraction of P applied to the tensor square of psi, minus four times p(psi) times psi. \bpsource{material/manuscript.tex}{447}{Manuscript, line 447.}

At every ray of an arbitrary complex Hilbert space, use a unit representative psi and the centered chart sending v to [psi+v] for v complex-orthogonal to psi. In these coordinates the Fubini–Study metric pairs v and w by the real part of their inner product. The field in Eq. (3.4) is the unique horizontal vector whose pairing with every tangent vector gives the real differential of ray acceptance. Changing the phase of psi changes the field by the same phase.

\bpdecls{\textbf{Lean.} \bpsource{GaussianTesting/GradientFlow/FubiniStudyGradient.lean}{89}{\texttt{fubiniStudyGradient\_rayAcceptance}}, \bpsource{GaussianTesting/GradientFlow/FubiniStudyGradient.lean}{27}{\texttt{TensorSquare.hasGradientAt\_rayAcceptance}}, \bpsource{GaussianTesting/GradientFlow/FubiniStudyGradient.lean}{51}{\texttt{TensorSquare.isFubiniStudyGradient\_iff}}, \bpsource{GaussianTesting/GradientFlow/FubiniStudyGradient.lean}{74}{\texttt{TensorSquare.gradient\_rayAcceptance\_phase}}, \bpsource{GaussianTesting/GradientFlow/RayTangentMetric.lean}{31}{\texttt{rayTangentMetric}}, \bpsource{GaussianTesting/GradientFlow/RayTangentMetric.lean}{101}{\texttt{tendsto\_dist\_div\_time\_rayOfTangent}}, \bpsource{GaussianTesting/GradientFlow/RayTangentMetric.lean}{134}{\texttt{tendsto\_rayTangentMetric\_of\_dist}}, \bpsource{GaussianTesting/GradientFlow/RayTangentMetric.lean}{60}{\texttt{rayOfTangent\_phase}}, \bpsource{GaussianTesting/GradientFlow/RayTangentMetric.lean}{68}{\texttt{rayTangentMetric\_phase}}, \bpsource{GaussianTesting/GradientFlow/RayTangentChart.lean}{56}{\texttt{centeredRayChart}}, \bpsource{GaussianTesting/GradientFlow/RayDifferential.lean}{52}{\texttt{TensorSquare.hasFDerivAt\_accProbA\_rayChartSection}}, \bpsource{GaussianTesting/GradientFlow/Abstract.lean}{236}{\texttt{TensorSquare.hasFDerivAt\_accProbA}}, \bpsource{GaussianTesting/GradientFlow/Abstract.lean}{120}{\texttt{TensorSquare.gradAccA}}, \bpsource{GaussianTesting/GradientFlow/Phase.lean}{52}{\texttt{TensorSquare.inner\_gradAccA\_eq\_zero}}, \bpsource{GaussianTesting/GradientFlow/Phase.lean}{37}{\texttt{TensorSquare.gradAccA\_phase}}, \bpsource{GaussianTesting/GradientFlow/Abstract.lean}{309}{\texttt{TensorSquare.contDiff\_gradAccA}}.}

\textbf{Blueprint.} \bplink{lem:acceptance-gradient}{Eq. (3.4), Fubini--Study gradient of acceptance}.

\subsection{Lemma 3.1}\label{paper:r3-1}
\textbf{Checked.} Every pure initial state has a smooth unique global gradient-flow solution. \bpsource{material/manuscript.tex}{456}{Manuscript, line 456.}

For a swap-commuting orthogonal projection on the completed tensor square of any complex Hilbert space, each initial ray has a solution on all nonnegative times. It is smooth in every affine chart through the current ray, including one-sided smoothness at zero, and every chart-defined solution with the same initial ray agrees throughout that domain.

\textbf{Proof route.} Both proofs glue local solutions by uniqueness and rule out a finite maximal time by an energy estimate. Lean projects smooth unit-vector solutions to rays. Local Lipschitz uniqueness in affine charts applies to arbitrary ray solutions. On each closed subinterval, the squared Fubini–Study distance is bounded by elapsed time times the increase in acceptance. Completeness gives a limiting ray at every putative finite endpoint. The chart ODE restarts there, with matching derivatives at the joining time, so the supremum of reachable times cannot be finite.

\bpdecls{\textbf{Lean.} \bpsource{GaussianTesting/GradientFlow/ProjectiveGradientFlow.lean}{86}{\texttt{projective\_gradient\_flow}}, \bpsource{GaussianTesting/GradientFlow/ProjectiveGradientFlow.lean}{63}{\texttt{TensorSquare.exists\_unique\_smooth\_projectiveFlow}}, \bpsource{GaussianTesting/GradientFlow/ProjectiveGradientFlow.lean}{28}{\texttt{SmoothRayCurveOn}}, \bpsource{GaussianTesting/GradientFlow/ProjectiveContinuation.lean}{27}{\texttt{TensorSquare.exists\_local\_projectiveFlow}}, \bpsource{GaussianTesting/GradientFlow/ProjectiveFlow.lean}{190}{\texttt{TensorSquare.IsProjectiveFlowOn.eqOn\_Icc}}, \bpsource{GaussianTesting/GradientFlow/ProjectiveFlow.lean}{114}{\texttt{TensorSquare.IsProjectiveFlowOn.contDiffWithinAt\_chart\_Ici}}, \bpsource{GaussianTesting/GradientFlow/ProjectiveGradientFlow.lean}{38}{\texttt{TensorSquare.IsProjectiveFlowOn.continue\_at\_finite\_time}}, \bpsource{GaussianTesting/GradientFlow/MaximalFlow.lean}{21}{\texttt{exists\_global\_of\_finite\_extension}}, \bpsource{GaussianTesting/GradientFlow/Abstract.lean}{649}{\texttt{TensorSquare.IsUnitFlowOn.hasDerivWithinAt\_accProb}}, \bpsource{GaussianTesting/GradientFlow/LocalFlow.lean}{62}{\texttt{TensorSquare.IsUnitFlowOn.integral\_norm\_grad\_sq\_eq}}, \bpsource{GaussianTesting/GradientFlow/LocalFlow.lean}{78}{\texttt{TensorSquare.IsUnitFlowOn.integral\_norm\_velocity\_sq\_eq}}, \bpsource{GaussianTesting/GradientFlow/ProjectiveEnergy.lean}{35}{\texttt{TensorSquare.rayAcceptance\_mk}}, \bpsource{GaussianTesting/GradientFlow/ProjectiveEnergy.lean}{44}{\texttt{TensorSquare.IsProjectiveFlowOn.exists\_unit\_lift\_Icc}}, \bpsource{GaussianTesting/GradientFlow/ProjectiveEnergy.lean}{73}{\texttt{TensorSquare.IsProjectiveFlowOn.dist\_sq\_le\_energy}}, \bpsource{GaussianTesting/GradientFlow/ProjectiveEnergy.lean}{87}{\texttt{TensorSquare.IsProjectiveFlowOn.dist\_sq\_le\_time\_sub}}, \bpsource{GaussianTesting/GradientFlow/ProjectiveEnergy.lean}{97}{\texttt{TensorSquare.IsProjectiveFlowOn.exists\_limit\_Ico}}.}

\textbf{Blueprint.} \bplink{lem:flow-existence}{Existence and uniqueness of the global gradient flow}.

\subsection{Integrated energy on a local flow}\label{paper:local-flow-energy}
\textbf{Checked.} In the proof of Lemma 3.1, the integral of squared speed from a to b equals the increase in acceptance over that interval. \bpsource{material/manuscript.tex}{470}{Manuscript, line 470.}

On every closed subinterval of a unit-vector solution, the integral of squared velocity equals the increase in acceptance, including on singleton intervals. Every solution of the chart ODE has such a lift on each closed subinterval. Thus its squared Fubini–Study displacement is at most elapsed time times the increase in the phase-independent acceptance, and hence at most elapsed time. This gives a limit of the ray trajectory at each finite right endpoint without a spectral gap assumption.

\bpdecls{\textbf{Lean.} \bpsource{GaussianTesting/GradientFlow/Abstract.lean}{649}{\texttt{TensorSquare.IsUnitFlowOn.hasDerivWithinAt\_accProb}}, \bpsource{GaussianTesting/GradientFlow/LocalFlow.lean}{62}{\texttt{TensorSquare.IsUnitFlowOn.integral\_norm\_grad\_sq\_eq}}, \bpsource{GaussianTesting/GradientFlow/LocalFlow.lean}{78}{\texttt{TensorSquare.IsUnitFlowOn.integral\_norm\_velocity\_sq\_eq}}, \bpsource{GaussianTesting/GradientFlow/ProjectiveEnergy.lean}{35}{\texttt{TensorSquare.rayAcceptance\_mk}}, \bpsource{GaussianTesting/GradientFlow/ProjectiveEnergy.lean}{44}{\texttt{TensorSquare.IsProjectiveFlowOn.exists\_unit\_lift\_Icc}}, \bpsource{GaussianTesting/GradientFlow/ProjectiveEnergy.lean}{73}{\texttt{TensorSquare.IsProjectiveFlowOn.dist\_sq\_le\_energy}}, \bpsource{GaussianTesting/GradientFlow/ProjectiveEnergy.lean}{87}{\texttt{TensorSquare.IsProjectiveFlowOn.dist\_sq\_le\_time\_sub}}, \bpsource{GaussianTesting/GradientFlow/ProjectiveEnergy.lean}{97}{\texttt{TensorSquare.IsProjectiveFlowOn.exists\_limit\_Ico}}.}

\textbf{Blueprint.} \bplink{lem:flow-existence}{Existence and uniqueness of the global gradient flow}.

\subsection{Contraction and the critical-point equation}\label{paper:contraction}
\textbf{Checked.} Contracting the projection with one copy defines K\_psi. The gradient vanishes exactly when psi is its eigenvector with eigenvalue p(psi). \bpsource{material/manuscript.tex}{485}{Manuscript, line 485.}

For every swap-commuting orthogonal projection on the completed tensor square of an arbitrary complex Hilbert space, contractionOperator is the bounded operator J\_psi† P J\_psi, where J\_psi x = x tensor psi. It is positive and self-adjoint for every psi, and has norm at most one for unit psi. Its first two operator moments are p and q. For unit psi, the actual differential of projective acceptance in the centered chart vanishes if and only if K\_psi psi = p(psi) psi.

\textbf{Proof route.} Use the manuscript compression K\_psi = J\_psi† P J\_psi. The adjoint agrees with the Riesz contraction, so K\_psi psi is the existing contracted vector. Positivity is preserved by compression; norm bounds for insertion, projection and contraction give contractivity for unit psi. Self-adjointness gives the second moment. The already checked Fubini–Study gradient formula identifies vanishing of the projective differential with vanishing of 4(K\_psi psi - p(psi) psi), hence with the stated eigenvector equation.

\bpdecls{\textbf{Lean.} \bpsource{GaussianTesting/GradientFlow/Contraction.lean}{34}{\texttt{TensorSquare.contractionOperator}}, \bpsource{GaussianTesting/GradientFlow/Contraction.lean}{38}{\texttt{TensorSquare.contractionOperator\_apply}}, \bpsource{GaussianTesting/GradientFlow/Contraction.lean}{47}{\texttt{TensorSquare.contractionOperator\_self}}, \bpsource{GaussianTesting/GradientFlow/Contraction.lean}{54}{\texttt{TensorSquare.contractionOperator\_isPositive}}, \bpsource{GaussianTesting/GradientFlow/Contraction.lean}{61}{\texttt{TensorSquare.contractionOperator\_isSelfAdjoint}}, \bpsource{GaussianTesting/GradientFlow/Contraction.lean}{66}{\texttt{TensorSquare.norm\_contractionOperator\_le\_one}}, \bpsource{GaussianTesting/GradientFlow/Contraction.lean}{79}{\texttt{TensorSquare.contractionOperator\_first\_moment}}, \bpsource{GaussianTesting/GradientFlow/Contraction.lean}{84}{\texttt{TensorSquare.contractionOperator\_second\_moment}}, \bpsource{GaussianTesting/GradientFlow/Contraction.lean}{91}{\texttt{TensorSquare.gradAccA\_eq\_contractionOperator}}, \bpsource{GaussianTesting/GradientFlow/Contraction.lean}{107}{\texttt{TensorSquare.gradAccA\_eq\_zero\_iff\_contractionOperator}}, \bpsource{GaussianTesting/GradientFlow/ContractionCriticalPoint.lean}{26}{\texttt{TensorSquare.fderiv\_rayAcceptance\_eq\_zero\_iff\_gradAccA\_eq\_zero}}, \bpsource{GaussianTesting/GradientFlow/ContractionCriticalPoint.lean}{43}{\texttt{TensorSquare.fderiv\_rayAcceptance\_eq\_zero\_iff\_contractionOperator}}, \bpsource{GaussianTesting/GradientFlow/ContractionCriticalPoint.lean}{54}{\texttt{contraction\_criticalPoint\_iff}}, \bpsource{GaussianTesting/GradientFlow/HilbertMoments.lean}{27}{\texttt{hilbert\_contr\_coords}}, \bpsource{GaussianTesting/GradientFlow/Abstract.lean}{90}{\texttt{TensorSquare.contr}}, \bpsource{GaussianTesting/GradientFlow/Defs.lean}{104}{\texttt{Kvec}}, \bpsource{GaussianTesting/GradientFlow/Abstract.lean}{114}{\texttt{TensorSquare.KvecA}}, \bpsource{GaussianTesting/GradientFlow/Abstract.lean}{120}{\texttt{TensorSquare.gradAccA}}.}

\textbf{Blueprint.} \bplink{lem:contraction-critical}{Contraction and critical-point characterization}, \bplink{lem:gradient-identity}{Gradient norm identity}.

\subsection{Gradient identity (Eq. 3.5)}\label{paper:gradient-norm}
\textbf{Checked.} The squared gradient norm is 16(q-p\textasciicircum{}2), where q is the squared norm of the contracted vector. \bpsource{material/manuscript.tex}{491}{Manuscript, line 491.}

For every unit input the stated identity holds in the abstract TensorSquare interface and, in particular, on the completed Hilbert tensor square. The operator-moment formulation also proves that this is 16 times the variance of the positive contraction K\_psi.

\bpdecls{\textbf{Lean.} \bpsource{GaussianTesting/GradientFlow/HilbertTensor.lean}{147}{\texttt{hilbertTensorSquare}}, \bpsource{GaussianTesting/GradientFlow/Lojasiewicz.lean}{49}{\texttt{norm\_gradAcc\_sq}}, \bpsource{GaussianTesting/GradientFlow/Abstract.lean}{198}{\texttt{TensorSquare.norm\_gradAccA\_sq}}, \bpsource{GaussianTesting/GradientFlow/Contraction.lean}{98}{\texttt{TensorSquare.norm\_gradAccA\_sq\_eq\_variance}}.}

\textbf{Blueprint.} \bplink{lem:gradient-identity}{Gradient norm identity}, \bplink{lem:bound-gradient}{Moment and gradient bounds from the spectrum}.

\subsection{Three-copy average (Eq. 3.6)}\label{paper:three-copy}
\textbf{Checked.} S is the average of the three pairwise copies of P acting on a triple tensor space. \bpsource{material/manuscript.tex}{501}{Manuscript, line 501.}

The average is defined on the completed tensor cube using P tensor I, I tensor P, and the factor permutation for the remaining pair. The linear span of elementary tensors is dense in the square and cube. Their matrix elements identify the three pair actions with the coordinate formulas for any Hilbert basis.

\bpdecls{\textbf{Lean.} \bpsource{GaussianTesting/GradientFlow/HilbertGradientBound.lean}{32}{\texttt{threeCopyAverage}}, \bpsource{GaussianTesting/GradientFlow/HilbertPairOperators.lean}{70}{\texttt{threeCopyPair12}}, \bpsource{GaussianTesting/GradientFlow/HilbertPairOperators.lean}{90}{\texttt{threeCopyPair13}}, \bpsource{GaussianTesting/GradientFlow/HilbertPairOperators.lean}{75}{\texttt{threeCopyPair23}}, \bpsource{GaussianTesting/GradientFlow/HilbertPairOperators.lean}{95}{\texttt{inner\_threeCopyPair12\_tmul}}, \bpsource{GaussianTesting/GradientFlow/HilbertPairOperators.lean}{111}{\texttt{inner\_threeCopyPair13\_tmul}}, \bpsource{GaussianTesting/GradientFlow/HilbertPairOperators.lean}{105}{\texttt{inner\_threeCopyPair23\_tmul}}, \bpsource{GaussianTesting/GradientFlow/HilbertTensor.lean}{97}{\texttt{dense\_span\_hilbertTmul}}, \bpsource{GaussianTesting/GradientFlow/HilbertTensor.lean}{110}{\texttt{dense\_span\_hilbertTriple}}, \bpsource{GaussianTesting/GradientFlow/HilbertTensor.lean}{245}{\texttt{hilbertCoords}}, \bpsource{GaussianTesting/GradientFlow/HilbertGradientBound.lean}{37}{\texttt{threeCopyAverage\_eq}}, \bpsource{GaussianTesting/GradientFlow/HilbertGradientBound.lean}{43}{\texttt{threeCopyAverage\_positive\_contraction}}.}

\textbf{Blueprint.} \bplink{prelim:flow}{Hilbert-space and projective preliminaries}, \bplink{lem:three-copy-average}{Three-copy average and pair operators}, \bplink{lem:bound-gradient}{Moment and gradient bounds from the spectrum}.

\subsection{Lemma 3.2}\label{paper:r3-2}
\textbf{Checked.} On three identical copies of a unit state, the first moment of S equals p and the second equals (p + 2q)/3. \bpsource{material/manuscript.tex}{507}{Manuscript, line 507.}

For every unit vector in any complex Hilbert space, the expectations of the three-copy average and its square on the elementary tensor cube are p and (p+2q)/3. The theorem applies directly to the completed tensor cube and assumes only a swap-commuting orthogonal projection.

\textbf{Proof route.} Both proofs expand the three diagonal terms and six cross terms of the squared average. Lean does this in arbitrary-index Hilbert coordinates and uses the Riesz characterization to transport the contraction, giving both complex expectations on the original space.

\bpdecls{\textbf{Lean.} \bpsource{GaussianTesting/GradientFlow/HilbertGradientBound.lean}{49}{\texttt{threeCopyAverage\_moments}}, \bpsource{GaussianTesting/GradientFlow/HilbertMoments.lean}{55}{\texttt{hilbert\_first\_moment}}, \bpsource{GaussianTesting/GradientFlow/HilbertMoments.lean}{67}{\texttt{hilbert\_second\_moment}}, \bpsource{GaussianTesting/GradientFlow/HilbertMoments.lean}{27}{\texttt{hilbert\_contr\_coords}}, \bpsource{GaussianTesting/GradientFlow/Moments.lean}{254}{\texttt{inner\_tp3\_Sop\_tp3}}, \bpsource{GaussianTesting/GradientFlow/Moments.lean}{263}{\texttt{inner\_Sop\_Sop\_tp3}}, \bpsource{GaussianTesting/Bosons/ThreeCopy.lean}{449}{\texttt{inner\_tp3lp\_Slp}}, \bpsource{GaussianTesting/Bosons/ThreeCopy.lean}{570}{\texttt{inner\_Slp\_Slp}}.}

\textbf{Blueprint.} \bplink{lem:three-copy-moments}{Three-copy moment identities}, \bplink{lem:bound-gradient}{Moment and gradient bounds from the spectrum}.

\subsection{Three-copy fixed points and certainty}\label{paper:triple-certainty}
\textbf{Checked.} The three-copy average is a positive contraction, fixing the tensor cube of a pure state exactly at certainty. Nonempty certainty gives eigenvalue one. \bpsource{material/manuscript.tex}{533}{Manuscript, line 533.}

On any complex Hilbert space, the average fixes the tensor cube of a unit vector exactly when its two-copy acceptance is one. Its complex spectrum lies in the real interval [0,1]. If the certainty set is nonempty, one is an eigenvalue and is the largest real eigenvalue.

\bpdecls{\textbf{Lean.} \bpsource{GaussianTesting/GradientFlow/FixedPoints.lean}{34}{\texttt{threeCopyAverage\_fixed\_iff\_accProb\_eq\_one}}, \bpsource{GaussianTesting/GradientFlow/FixedPoints.lean}{70}{\texttt{threeCopyAverage\_spectrum\_subset\_unitInterval}}, \bpsource{GaussianTesting/GradientFlow/FixedPoints.lean}{56}{\texttt{threeCopyAverage\_hasEigenvalue\_one}}, \bpsource{GaussianTesting/GradientFlow/FixedPoints.lean}{79}{\texttt{threeCopyAverage\_isGreatest\_eigenvalue}}.}

\textbf{Blueprint.} \bplink{lem:triple-certainty}{Fixed tensor cubes and certainty states}.

\subsection{Footnote before Lemma 3.3: spectral gap on the symmetric cube}\label{paper:symmetric-cube-gap}
\textbf{Checked.} The spectral gap hypothesis in Lemma 3.3 need only hold on Sym³(H). \bpsource{material/manuscript.tex}{536}{Manuscript, line 536.}

On the completed Hilbert tensor cube, Sym³(H) is the closed common fixed space of the two adjacent factor swaps. Lean proves that the three-copy average commutes with these swaps and preserves this subspace. Its restriction is a positive contraction with the same first two moments on every unit tensor cube. The restricted complex spectral inclusion yields the manuscript’s moment and squared-gradient bounds, without assuming a full-cube gap or adding invariance as a hypothesis.

\textbf{Proof route.} Restrict S to the permutation-fixed Hilbert subspace, evaluate the positive spectral polynomial (S-I)(S-alphaI) on the tensor cube of psi there, substitute Lemma 3.2’s unchanged two moments, and apply Eq. (3.5). This follows the manuscript’s footnote and proof of Lemma 3.3.

\bpdecls{\textbf{Lean.} \bpsource{GaussianTesting/GradientFlow/SymmetricCube.lean}{27}{\texttt{cubeSwap12}}, \bpsource{GaussianTesting/GradientFlow/SymmetricCube.lean}{32}{\texttt{cubeSwap23}}, \bpsource{GaussianTesting/GradientFlow/SymmetricCube.lean}{74}{\texttt{symmetricCube}}, \bpsource{GaussianTesting/GradientFlow/SymmetricCube.lean}{83}{\texttt{isClosed\_symmetricCube}}, \bpsource{GaussianTesting/GradientFlow/SymmetricCube.lean}{92}{\texttt{hilbertTmul\_mem\_symmetricCube}}, \bpsource{GaussianTesting/GradientFlow/SymmetricCube.lean}{107}{\texttt{threeCopyAverage\_cubeSwap12}}, \bpsource{GaussianTesting/GradientFlow/SymmetricCube.lean}{125}{\texttt{threeCopyAverage\_cubeSwap23}}, \bpsource{GaussianTesting/GradientFlow/SymmetricCube.lean}{143}{\texttt{threeCopyAverage\_maps\_symmetricCube}}, \bpsource{GaussianTesting/GradientFlow/SymmetricCubeGradient.lean}{29}{\texttt{symmetricThreeCopyAverage}}, \bpsource{GaussianTesting/GradientFlow/SymmetricCubeGradient.lean}{52}{\texttt{symmetricTensorCube}}, \bpsource{GaussianTesting/GradientFlow/SymmetricCubeGradient.lean}{39}{\texttt{symmetricThreeCopyAverage\_positive\_contraction}}, \bpsource{GaussianTesting/GradientFlow/SymmetricCubeGradient.lean}{56}{\texttt{symmetricThreeCopyAverage\_moments}}, \bpsource{GaussianTesting/GradientFlow/SymmetricCubeGradient.lean}{67}{\texttt{lojasiewiczData\_of\_symmetricThreeCopy\_spectrum}}, \bpsource{GaussianTesting/GradientFlow/SymmetricCubeGradient.lean}{87}{\texttt{norm\_gradAcc\_sq\_ge\_of\_symmetricThreeCopy\_spectrum}}.}

\textbf{Blueprint.} \bplink{lem:bound-gradient}{Moment and gradient bounds from the spectrum}.

\subsection{Lemma 3.3}\label{paper:r3-3}
\textbf{Checked.} If the spectrum of S lies in [0,alpha] union \{1\} for 0 <= alpha <= 1, then the squared gradient norm is at least 16(p - 3 alpha/2)(1 - p). \bpsource{material/manuscript.tex}{538}{Manuscript, line 538.}

For every unit vector, the spectrum of the average on the completed Hilbert tensor cube gives q >= p - (3 alpha/2)(1-p) and the stated squared-gradient lower bound. Projection implies the required norm bound. The footnote strengthening is also checked: the same conclusions require the spectral gap only on the restriction to the closed symmetric tensor cube.

\textbf{Proof route.} Both the full-cube and symmetric-cube proofs apply positivity of the spectral polynomial to the elementary tensor cube and substitute the two identities of Lemma 3.2. For the footnote, Lean proves invariance of the closed symmetric cube under S and restricts S before applying the same argument. Rearranging gives the inequality for q, then Eq. (3.5) gives the gradient bound.

\bpdecls{\textbf{Lean.} \bpsource{GaussianTesting/GradientFlow/HilbertGradientBound.lean}{43}{\texttt{threeCopyAverage\_positive\_contraction}}, \bpsource{GaussianTesting/GradientFlow/HilbertGradientBound.lean}{75}{\texttt{lojasiewiczData\_of\_threeCopy\_spectrum}}, \bpsource{GaussianTesting/GradientFlow/HilbertGradientBound.lean}{92}{\texttt{norm\_gradAcc\_sq\_ge\_of\_threeCopy\_spectrum}}, \bpsource{GaussianTesting/GradientFlow/Lojasiewicz.lean}{76}{\texttt{qProb\_ge\_of\_gap}}, \bpsource{GaussianTesting/GradientFlow/Lojasiewicz.lean}{97}{\texttt{norm\_gradAcc\_sq\_ge}}, \bpsource{GaussianTesting/GradientFlow/Abstract.lean}{124}{\texttt{TensorSquare.LojasiewiczData}}, \bpsource{GaussianTesting/Bosons/ThreeCopy.lean}{601}{\texttt{lojasiewiczData\_of\_Slp\_gap}}, \bpsource{GaussianTesting/GradientFlow/Lojasiewicz.lean}{69}{\texttt{SpectralGapHyp}}, \bpsource{GaussianTesting/GradientFlow/Spectral.lean}{44}{\texttt{spectrum\_gap\_iff\_polynomial\_nonneg}}, \bpsource{GaussianTesting/GradientFlow/Spectral.lean}{75}{\texttt{spectrum\_gap\_iff\_quadraticForm\_nonneg}}, \bpsource{GaussianTesting/GradientFlow/SymmetricCubeGradient.lean}{67}{\texttt{lojasiewiczData\_of\_symmetricThreeCopy\_spectrum}}, \bpsource{GaussianTesting/GradientFlow/SymmetricCubeGradient.lean}{87}{\texttt{norm\_gradAcc\_sq\_ge\_of\_symmetricThreeCopy\_spectrum}}.}

\textbf{Blueprint.} \bplink{lem:bound-gradient}{Moment and gradient bounds from the spectrum}.

\subsection{Spectral gap hypothesis}\label{paper:spectral-gap}
\textbf{Checked.} The self-adjoint positive contraction S has spectrum in [0,alpha] union \{1\}. \bpsource{material/manuscript.tex}{539}{Manuscript, line 539.}

For any positive contraction on a complex Hilbert space, continuous functional calculus proves that the complex spectral containment is equivalent to positivity of (S-I)(S-alphaI), equivalently to nonnegativity of its quadratic form on all vectors. The finite predicate SpectralGapHyp uses this quadratic form.

\bpdecls{\textbf{Lean.} \bpsource{GaussianTesting/GradientFlow/Lojasiewicz.lean}{69}{\texttt{SpectralGapHyp}}, \bpsource{GaussianTesting/GradientFlow/Spectral.lean}{44}{\texttt{spectrum\_gap\_iff\_polynomial\_nonneg}}, \bpsource{GaussianTesting/GradientFlow/Spectral.lean}{75}{\texttt{spectrum\_gap\_iff\_quadraticForm\_nonneg}}.}

\textbf{Blueprint.} \bplink{lem:bound-gradient}{Moment and gradient bounds from the spectrum}, \bplink{lem:S-spectrum-coherent-states}{Spectrum of the three-copy coherent test}.

\subsection{Theorem 3.4}\label{paper:r3-4}
\textbf{Checked.} Suppose the three-copy spectrum lies in [0,alpha] union \{1\}, with 0 <= alpha < 2/3. For p > 3 alpha/2, the angle arccos(F) is at most half of arcsin(sqrt(C\_alpha (1 - p))), with C\_alpha = 1/(1 - 3 alpha/2). For every tau >= 0, p >= 1 - tau gives F\textasciicircum{}2 >= 1 - C\_alpha tau. For each input, one certainty-state witness satisfies the stated bounds with its own overlap. \bpsource{material/manuscript.tex}{561}{Manuscript, line 561.}

The two stated bounds hold for a single certainty-state witness for each unit input on any complex Hilbert space. The theorem takes the complex spectrum of the three-copy average, 0 <= alpha < 2/3, and a nonempty certainty set. The angular bound holds above threshold and the linear bound for every admissible tau. Both give the corresponding bounds on supremum fidelity. The public spectral theorem now uses the projective flow and its certainty limit to prove both conclusions.

\textbf{Proof route.} The public spectral theorem follows the manuscript projective argument. Lemma 3.1 supplies the global projective flow from the input ray. For 3 alpha/2 < p\_0 < 1, Eq. (3.7) gives strict increase and convergence of acceptance to one. The path-length estimate bounds every tail by the half-arcsine potential, making the real-time trajectory Cauchy. Completeness gives a limiting ray, and continuity of acceptance makes it a certainty ray. A normalized representative supplies the angular witness because Fubini–Study distance is the arccosine of the absolute overlap. The manuscript cosine estimate gives the linear bound for the same witness and every admissible error. At acceptance one, the input is its own witness. At or below the threshold, the linear lower bound is nonpositive and any member of the nonempty certainty set suffices. Comparing witness overlap with supremum fidelity proves both fidelity bounds.

\bpdecls{\textbf{Lean.} \bpsource{GaussianTesting/GradientFlow/HilbertSpectralGap.lean}{50}{\texttt{robustness\_of\_threeCopy\_spectrum}}, \bpsource{GaussianTesting/GradientFlow/HilbertSpectralGap.lean}{27}{\texttt{angular\_witness\_of\_threeCopy\_spectrum}}, \bpsource{GaussianTesting/GradientFlow/HilbertSpectralGap.lean}{60}{\texttt{angular\_fidelity\_of\_threeCopy\_spectrum}}, \bpsource{GaussianTesting/GradientFlow/SharpRobustness.lean}{36}{\texttt{bounds\_of\_angular\_witness}}, \bpsource{GaussianTesting/GradientFlow/HilbertSpectralGap.lean}{72}{\texttt{testing\_fidelity\_of\_threeCopy\_spectrum}}, \bpsource{GaussianTesting/GradientFlow/SharpRobustness.lean}{25}{\texttt{RobustnessBounds}}, \bpsource{GaussianTesting/GradientFlow/SharpRobustness.lean}{76}{\texttt{RobustnessBounds.angular\_fidelity}}, \bpsource{GaussianTesting/GradientFlow/SharpRobustness.lean}{69}{\texttt{RobustnessBounds.linear\_fidelity}}, \bpsource{GaussianTesting/GradientFlow/RayGradientBound.lean}{37}{\texttt{ray\_flow\_derivative\_bounds}}, \bpsource{GaussianTesting/GradientFlow/RayGradientBound.lean}{24}{\texttt{rayGradientNormSq\_ge\_of\_threeCopy\_spectrum}}, \bpsource{GaussianTesting/GradientFlow/RayEnergyIdentity.lean}{32}{\texttt{TensorSquare.rayGradientNormSq\_mk}}, \bpsource{GaussianTesting/GradientFlow/RaySpeed.lean}{43}{\texttt{tendsto\_dist\_div\_time\_of\_horizontal\_lift}}, \bpsource{GaussianTesting/GradientFlow/RayEnergyIdentity.lean}{57}{\texttt{TensorSquare.IsProjectiveFlowOn.hasDerivWithinAt\_rayAcceptance}}, \bpsource{GaussianTesting/GradientFlow/RayEnergyIdentity.lean}{68}{\texttt{TensorSquare.IsProjectiveFlowOn.rayAcceptance\_mono}}, \bpsource{GaussianTesting/GradientFlow/RayEnergyIdentity.lean}{77}{\texttt{TensorSquare.IsProjectiveFlowOn.tendsto\_dist\_div\_time}}.}

\textbf{Theorem node.} \bplink{thm:bound-fidelity-gradient-flow}{Sharp angular and linear bounds, including Eq. (3.7)}.

\textbf{Related blueprint entries.} \bplink{lem:bound-gradient}{Moment and gradient bounds from the spectrum}, \bplink{thm:gradient-flow-intro}{Linear fidelity bound from the spectral gap}, \bplink{thm:main-fermions}{Fidelity bound for the fermionic test}, \bplink{thm:main-boson-mean-zero}{Fidelity bound for mean-zero Gaussian states}, \bplink{thm:main-boson-coherent}{Fidelity bound for coherent states}.

\subsection{Energy identity along the flow (Eq. 3.7)}\label{paper:energy-identity}
\textbf{Checked.} The derivative of p along the flow is the squared speed, with a lower bound from the spectral gap. \bpsource{material/manuscript.tex}{586}{Manuscript, line 586.}

For any forward ray solution and any t >= 0, the derivative of acceptance equals both the squared Fubini–Study metric speed and the squared norm of the field in Eq. (3.4). The derivative and speed at zero are taken from the right. If the three-copy spectrum lies in [0,alpha] union \{1\} with 0 <= alpha <= 1, the derivative is at least 16(p\_t-3 alpha/2)(1-p\_t), which is at least 16(p\_0-3 alpha/2)(1-p\_t).

\bpdecls{\textbf{Lean.} \bpsource{GaussianTesting/GradientFlow/RayGradientBound.lean}{37}{\texttt{ray\_flow\_derivative\_bounds}}, \bpsource{GaussianTesting/GradientFlow/RayGradientBound.lean}{24}{\texttt{rayGradientNormSq\_ge\_of\_threeCopy\_spectrum}}, \bpsource{GaussianTesting/GradientFlow/RayEnergyIdentity.lean}{32}{\texttt{TensorSquare.rayGradientNormSq\_mk}}, \bpsource{GaussianTesting/GradientFlow/RaySpeed.lean}{43}{\texttt{tendsto\_dist\_div\_time\_of\_horizontal\_lift}}, \bpsource{GaussianTesting/GradientFlow/RayEnergyIdentity.lean}{57}{\texttt{TensorSquare.IsProjectiveFlowOn.hasDerivWithinAt\_rayAcceptance}}, \bpsource{GaussianTesting/GradientFlow/RayEnergyIdentity.lean}{68}{\texttt{TensorSquare.IsProjectiveFlowOn.rayAcceptance\_mono}}, \bpsource{GaussianTesting/GradientFlow/RayEnergyIdentity.lean}{77}{\texttt{TensorSquare.IsProjectiveFlowOn.tendsto\_dist\_div\_time}}.}

\textbf{Blueprint.} \bplink{thm:bound-fidelity-gradient-flow}{Sharp angular and linear bounds, including Eq. (3.7)}.

\subsection{Strict increase and convergence of acceptance after Eq. (3.7)}\label{paper:flow-acceptance-limit}
\textbf{Checked.} When 3 alpha/2 < p\_0 < 1, acceptance along the projective flow is strictly increasing, stays below one at every finite time, and tends to one. \bpsource{material/manuscript.tex}{595}{Manuscript, line 595.}

For every forward ray trajectory on an arbitrary complex Hilbert space, the three-copy spectral inclusion with 0 <= alpha <= 1 and 3 alpha/2 < p\_0 < 1 gives StrictMonoOn on [0,infinity), p\_t < 1 for all finite t >= 0, and convergence to one as real time tends to infinity. The rejection estimate 1-p\_t >= exp(-16t)(1-p\_0) needs no spectral assumption. The convergence result alone requires only p\_0 > 3 alpha/2. These projective estimates supply the acceptance step in the checked proof of Theorem 3.4.

\textbf{Proof route.} The upper derivative bound makes exp(16t)(1-p\_t) nondecreasing, proving that p\_t stays below one. Equation (3.7) then gives a positive derivative, hence strict increase. Bounded monotonicity supplies the supremum limit. If it were below one, Eq. (3.7) would force a uniform positive derivative and contradict p\_t <= 1. Monotonicity extends convergence along integer times to all real forward times.

\bpdecls{\textbf{Lean.} \bpsource{GaussianTesting/GradientFlow/RayAcceptanceLimit.lean}{46}{\texttt{ray\_acceptance\_strictly\_increases\_to\_one}}, \bpsource{GaussianTesting/GradientFlow/RayAcceptanceLimit.lean}{25}{\texttt{rayAcceptance\_strictMonoOn}}, \bpsource{GaussianTesting/GradientFlow/RayConvergence.lean}{25}{\texttt{rayAcceptance\_tendsto\_one}}, \bpsource{GaussianTesting/GradientFlow/RayNonattainment.lean}{24}{\texttt{TensorSquare.IsProjectiveFlowOn.one\_sub\_rayAcceptance\_ge}}, \bpsource{GaussianTesting/GradientFlow/RayNonattainment.lean}{64}{\texttt{TensorSquare.IsProjectiveFlowOn.rayAcceptance\_lt\_one}}.}

\textbf{Blueprint.} \bplink{thm:bound-fidelity-gradient-flow}{Sharp angular and linear bounds, including Eq. (3.7)}.

\subsection{Path-length estimate and convergence to certainty in Theorem 3.4}\label{paper:flow-ray-limit}
\textbf{Checked.} For 3 alpha/2 < p\_0 < 1, the manuscript path-length estimate makes every forward projective trajectory Cauchy. It converges in real time to a certainty ray, with the same sharp angular bound from each trajectory point to the limit. \bpsource{material/manuscript.tex}{601}{Manuscript, line 601.}

For every forward projective solution on an arbitrary complex Hilbert space under the three-copy spectral inclusion with 0 <= alpha <= 1 and 3 alpha/2 < p\_0 < 1, for 0 <= s <= t the integrated Fubini–Study speed on [s,t] is at most V(p\_s)-V(p\_t), where V(p) = arcsin(sqrt((1-p)/(1-3 alpha/2)))/2. The trajectory is Cauchy as real time tends to infinity and converges to a ray with acceptance one. Its distance from the ray at each s >= 0 is at most V(p\_s), and a normalized representative of the limit belongs to the certainty set. The checked proof of Theorem 3.4 uses the limiting ray to obtain its angular witness and both fidelity bounds.

\textbf{Proof route.} A horizontal lift and the partition estimate for angles bound Fubini–Study distance by integrated speed. The chain rule and fundamental theorem of calculus for V(p\_t) express the manuscript change of variables and evaluated integral; Eq. (3.7) bounds speed by minus the potential derivative. Since acceptance tends to one and V(1)=0, the tail bound tends to zero, proving the Cauchy property. Completeness of the Fubini–Study metric supplies the limiting ray. Continuity of ray acceptance identifies it as a certainty ray, and continuity of distance passes the angular bound to the limit.

\bpdecls{\textbf{Lean.} \bpsource{GaussianTesting/GradientFlow/RayPathLength.lean}{41}{\texttt{TensorSquare.IsProjectiveFlowOn.dist\_le\_integral\_speed}}, \bpsource{GaussianTesting/GradientFlow/RayLength.lean}{25}{\texttt{ray\_integral\_speed\_le\_potential\_sub}}, \bpsource{GaussianTesting/GradientFlow/RayLength.lean}{86}{\texttt{ray\_dist\_le\_potential\_sub}}, \bpsource{GaussianTesting/GradientFlow/RayTrajectoryLimit.lean}{25}{\texttt{ray\_dist\_le\_potential}}, \bpsource{GaussianTesting/GradientFlow/RayCertaintyLimit.lean}{26}{\texttt{TensorSquare.continuous\_rayAcceptance}}, \bpsource{GaussianTesting/GradientFlow/RayCertaintyLimit.lean}{41}{\texttt{TensorSquare.cauchy\_ray\_of\_dist\_le\_potential}}, \bpsource{GaussianTesting/GradientFlow/RayCertaintyLimit.lean}{64}{\texttt{TensorSquare.exists\_certainty\_limit\_of\_dist\_le\_potential}}, \bpsource{GaussianTesting/GradientFlow/RayTrajectoryLimit.lean}{38}{\texttt{ray\_trajectory\_cauchy}}, \bpsource{GaussianTesting/GradientFlow/RayTrajectoryLimit.lean}{51}{\texttt{ray\_trajectory\_tendsto\_certainty}}, \bpsource{GaussianTesting/GradientFlow/RayTrajectoryLimit.lean}{66}{\texttt{ray\_trajectory\_tendsto\_certaintyState}}.}

\textbf{Blueprint.} \bplink{thm:bound-fidelity-gradient-flow}{Sharp angular and linear bounds, including Eq. (3.7)}.

\section{Manuscript section 4. Kac walk}
\subsection{SO(3) axis-invariant subspaces}\label{paper:projections-invariant}
\textbf{Checked.} For each axis of a strongly continuous unitary SO(3) representation, project onto the vectors fixed by the corresponding circle. \bpsource{material/manuscript.tex}{651}{Manuscript, line 651.}

The three coordinate-axis subgroups and their closed fixed subspaces are defined for any unitary SO(3) representation. For a strongly continuous representation, each orthogonal projection is the normalized Hilbert-space Bochner integral of its axis orbit over a full turn. On finite-dimensional invariant subspaces, strong continuity also gives Hermitian angular-momentum generators with exactly these kernels.

\textbf{Proof route.} Periodicity of the vector orbit makes its normalized integral invariant under the axis subgroup. Pairing with every fixed vector is unchanged by averaging, identifying the integral as the orthogonal projection.

\bpdecls{\textbf{Lean.} \bpsource{GaussianTesting/Kac/SO3AxisSubgroups.lean}{42}{\texttt{SO3.axisSubgroups}}, \bpsource{GaussianTesting/Kac/UnitaryRepresentation.lean}{107}{\texttt{UnitaryRepresentation.fixedSubspace}}, \bpsource{GaussianTesting/Kac/UnitaryRepresentation.lean}{113}{\texttt{UnitaryRepresentation.fixedSubspace\_isClosed}}, \bpsource{GaussianTesting/Kac/SO3Representation.lean}{66}{\texttt{SO3.exists\_angularMomentum}}, \bpsource{GaussianTesting/Kac/AxisAverage.lean}{31}{\texttt{SO3.axisRotation}}, \bpsource{GaussianTesting/Kac/AxisAverage.lean}{75}{\texttt{SO3.axisCircleAverage}}, \bpsource{GaussianTesting/Kac/AxisAverage.lean}{133}{\texttt{SO3.axisCircleAverage\_eq\_fixedProjection}}, \bpsource{GaussianTesting/Kac/AxisAverage.lean}{160}{\texttt{SO3.axisCircleAverageCLM\_eq\_fixedProjection}}.}

\textbf{Blueprint.} \bplink{prelim:kac}{Kac preliminaries}, \bplink{thm:kac-walk}{Spectral gap for the Kac walk}, \bplink{thm:main-fermions}{Fidelity bound for the fermionic test}.

\subsection{Kac averaging operator (Eq. 4.1)}\label{paper:kac-operator}
\textbf{Checked.} The average of the three axis projections specializes to the Kac transition operator on the square-integrable functions on the sphere. \bpsource{material/manuscript.tex}{653}{Manuscript, line 653.}

SO3.kacAverage averages the three axis-invariant projections for an arbitrary unitary representation. Each projection equals uniform circle averaging. On the actual sphere with normalized Euclidean surface measure, the transition is defined by uniform axis/angle integration. Its forward-rotation formula is proved pointwise for continuous observables; the unique bounded extension to complex L²(S²) equals SO3.kacAverage sphereRepresentation. No everywhere pointwise formula for arbitrary L² representatives is asserted.

\textbf{Proof route.} Define circle averages as normalized vector-valued integrals and identify them with fixed-space projections by periodicity and orthogonality. Integrate continuous observables in their uniform norm; continuous evaluation and the embedding into L² commute with these integrals. Reflection over a full period converts inverse pullback to forward rotation. Averaging the three axes gives the transition formula, and density of continuous observables identifies its unique bounded L² extension.

\bpdecls{\textbf{Lean.} \bpsource{GaussianTesting/Kac/Defs.lean}{60}{\texttt{kacT}}, \bpsource{GaussianTesting/Kac/GeneralRepresentation.lean}{36}{\texttt{SO3.kacAverage}}, \bpsource{GaussianTesting/Kac/GroupAverage.lean}{42}{\texttt{UnitaryRepresentation.fixedProjection\_restrict}}, \bpsource{GaussianTesting/Kac/AxisAverage.lean}{133}{\texttt{SO3.axisCircleAverage\_eq\_fixedProjection}}, \bpsource{GaussianTesting/Kac/SphereKacTransition.lean}{84}{\texttt{SO3.sphereAxisContinuousAverage}}, \bpsource{GaussianTesting/Kac/SphereKacTransition.lean}{118}{\texttt{SO3.sphereAxisContinuousAverage\_apply}}, \bpsource{GaussianTesting/Kac/SphereKacTransition.lean}{139}{\texttt{SO3.sphereContinuousTransition}}, \bpsource{GaussianTesting/Kac/SphereKacTransition.lean}{144}{\texttt{SO3.sphereContinuousTransition\_apply}}, \bpsource{GaussianTesting/Kac/SphereKacTransition.lean}{80}{\texttt{SO3.sphereContinuousToLp\_denseRange}}, \bpsource{GaussianTesting/Kac/SphereKacTransition.lean}{154}{\texttt{SO3.sphereKacTransition}}, \bpsource{GaussianTesting/Kac/SphereKacTransition.lean}{161}{\texttt{SO3.sphereKacTransition\_apply}}, \bpsource{GaussianTesting/Kac/SphereKacTransition.lean}{170}{\texttt{SO3.sphereKacTransition\_eq\_kacAverage}}, \bpsource{GaussianTesting/Kac/SphereKacTransition.lean}{176}{\texttt{SO3.sphereKacTransition\_continuous}}, \bpsource{GaussianTesting/Kac/SphereKacTransition.lean}{186}{\texttt{SO3.sphereKacTransition\_unique}}.}

\textbf{Blueprint.} \bplink{prelim:kac}{Kac preliminaries}, \bplink{thm:kac-walk}{Spectral gap for the Kac walk}.

\subsection{Kac sphere measure and rotation action}\label{paper:sphere-action}
\textbf{Checked.} The natural rotation action on complex L² of the unit sphere, with normalized Euclidean surface measure, is a strongly continuous unitary representation of SO(3). \bpsource{material/manuscript.tex}{656}{Manuscript, line 656.}

The sphere is the actual unit sphere in Euclidean real three-space. Its probability measure is Mathlib’s Euclidean surface measure volume.toSphere divided by its nonzero finite total mass. Ordinary matrix rotations preserve both the surface measure and its normalization. Inverse pullback defines a complex linear isometric equivalence for each rotation, with the representation multiplication law and the almost-everywhere formula (pi(g)f)(x)=f(g\textasciicircum{}(-1)x). Every L² orbit is continuous in the norm topology. This background supports the precise spherical Kac-gap calculation; the circle-average transition identification is checked under Kac averaging operator, while the fixed-space characterization and the exact spherical gap, including spin-four attainment, are checked as part of Theorem 4.1.

\textbf{Proof route.} Normalize the finite nonzero surface measure. A rotation maps radial cones to the corresponding radial cones and preserves Euclidean volume, proving surface-measure invariance. Measure-preserving inverse pullback gives the unitary action and its group law. Joint continuity of point rotations gives continuous variation in the continuous-map topology; continuity of L² composition then proves strong continuity for every square-integrable function.

\bpdecls{\textbf{Lean.} \bpsource{GaussianTesting/Kac/SphereMeasure.lean}{29}{\texttt{SO3.SphereAmbient}}, \bpsource{GaussianTesting/Kac/SphereMeasure.lean}{32}{\texttt{SO3.Sphere}}, \bpsource{GaussianTesting/Kac/SphereMeasure.lean}{35}{\texttt{SO3.sphereSurfaceMeasure}}, \bpsource{GaussianTesting/Kac/SphereMeasure.lean}{45}{\texttt{SO3.sphereMeasure}}, \bpsource{GaussianTesting/Kac/SphereMeasure.lean}{53}{\texttt{SO3.sphereMeasure\_univ}}, \bpsource{GaussianTesting/Kac/SphereMeasure.lean}{69}{\texttt{SO3.rotationEquiv}}, \bpsource{GaussianTesting/Kac/SphereMeasure.lean}{87}{\texttt{SO3.rotationEquiv\_apply}}, \bpsource{GaussianTesting/Kac/SphereMeasure.lean}{101}{\texttt{SO3.sphereRotation}}, \bpsource{GaussianTesting/Kac/SphereMeasure.lean}{108}{\texttt{SO3.sphereRotation\_one}}, \bpsource{GaussianTesting/Kac/SphereMeasure.lean}{112}{\texttt{SO3.sphereRotation\_mul}}, \bpsource{GaussianTesting/Kac/SphereMeasure.lean}{118}{\texttt{SO3.continuous\_sphereRotation}}, \bpsource{GaussianTesting/Kac/SphereMeasure.lean}{155}{\texttt{SO3.sphereRotation\_surfaceMeasurePreserving}}, \bpsource{GaussianTesting/Kac/SphereMeasure.lean}{172}{\texttt{SO3.sphereRotation\_measurePreserving}}, \bpsource{GaussianTesting/Kac/SphereRepresentation.lean}{32}{\texttt{SO3.SphereSpace}}, \bpsource{GaussianTesting/Kac/SphereRepresentation.lean}{51}{\texttt{SO3.sphereRotationAction}}, \bpsource{GaussianTesting/Kac/SphereRepresentation.lean}{67}{\texttt{SO3.sphereRotationAction\_coeFn}}, \bpsource{GaussianTesting/Kac/SphereRepresentation.lean}{77}{\texttt{SO3.sphereRotationAction\_mul}}, \bpsource{GaussianTesting/Kac/SphereRepresentation.lean}{85}{\texttt{SO3.sphereRepresentation}}, \bpsource{GaussianTesting/Kac/SphereRepresentation.lean}{109}{\texttt{SO3.continuous\_sphereRepresentation\_orbit}}.}

\textbf{Blueprint.} \bplink{prelim:kac}{Kac preliminaries}.

\subsection{Theorem 4.1}\label{paper:r4-1}
\textbf{Checked.} The Kac walk on S\textasciicircum{}2 has exact spectral gap 5/12, and every strongly continuous SO(3) representation has spectrum in [0,7/12] union \{1\} for its axis average. \bpsource{material/manuscript.tex}{666}{Manuscript, line 666.}

Both assertions of Theorem 4.1 are checked. Every strongly continuous unitary SO(3) representation on a complete complex Hilbert space has Kac spectrum in [0,7/12] union \{1\}, without finite-dimensional or separability assumptions. On the actual sphere with normalized Euclidean surface measure, the transition is the uniform circle average and its fixed space is exactly the constants. The restriction to their orthogonal complement has norm exactly 7/12, spectrum contained in [0,7/12], and an explicit nonzero eigenvector attaining 7/12. The resulting discrete-time spectral gap is 5/12. The witness uses the three degree-four zonal harmonics, with their homogeneous harmonic polynomial certificates and the manuscript coefficient a\textsubscript{4} = 3/8; a general classification of all spherical harmonics is not required. The full manuscript integer-spin calculation is also checked: the alternating binomial-square overlap, exact nonzero Gram eigenvalues, actual largest eigenvalue, small-spin values and strict j≥6 bound. The public spectral theorem uses these exact results.

\textbf{Proof route.} Following the manuscript's representation decomposition, Haar averages of rank-one orbit operators give compact commuting operators and finite-dimensional invariant eigenspaces. An orthogonal-complement argument proves that finite-dimensional irreducibles span densely. On each finite invariant subspace, continuous axis actions have skew-adjoint generators; local Euler factorization gives their Lie commutators and their kernels are the axis-fixed spaces. The normalized weight-string basis is a coordinate realization of the manuscript's integer-spin representation. Each axis kernel is one-dimensional, and choosing phases makes all pairwise overlaps equal to a\_j = 2\textasciicircum{}(-j) sum\_k (-1)\textasciicircum{}k binom(j,k)\textasciicircum{}2. The middle coefficient of (1-X)\textasciicircum{}j(1+X)\textasciicircum{}j=(1-X\textasciicircum{}2)\textasciicircum{}j proves its odd/even closed forms. The common three-by-three Gram matrix has characteristic polynomial (z-(1+2a\_j))(z-(1-a\_j))\textasciicircum{}2. Transfer between the three kernel vectors and the rank-one projection average proves that its nonzero spectral values are exactly (1+2a\_j)/3 and (1-a\_j)/3, and their maximum is the actual largest eigenvalue. The cases j=2 and j=4 give 1/2 and 7/12; the central-binomial recurrence makes the absolute even-spin overlaps strictly decrease, yielding a strict bound below 7/12 for every j≥6. The public irreducible gap theorem now derives its quadratic-form inequality from this exact maximum. Fixed-space projections restrict correctly, and continuity extends the quadratic-form inequality to all vectors, yielding the spectral inclusion. The spherical measure, rotation action, and circle-average transition identification are checked in Kac preliminaries. A vector fixed by the average of the three orthogonal projections lies in every axis-fixed space. Coordinate-plane generation identifies this with full SO(3) invariance. Rotations act transitively on the sphere, so left invariance of normalized Haar measure makes the average of a continuous observable constant. Boundedness of vector-valued Haar averaging, density of continuous observables, and closedness of the constant subspace extend this to L². An invariant vector equals its Haar average. On each finite invariant subspace of the orthogonal complement of constants, the fixed-space characterization excludes eigenvalue one. The existing spectral inclusion then gives norm at most 7/12; density passes this bound to the full complement. For attainment, the homogeneous quartic H\_a(x)=(35 x\_a\textasciicircum{}4−30 x\_a² r²+3 r\textasciicircum{}4)/8 has vanishing polynomial Laplacian. Its restriction h\_a=P\textsubscript{4}(x\_a) is fixed by its own axis average and maps to (3/8)h\_b under a different axis average, realizing the manuscript’s j=4 matrix. The sum of the three h\_a is a 7/12 eigenfunction, nonzero at a coordinate pole and hence nonzero in L² because sphere measure has full support. Symmetry places it orthogonal to constants. Therefore the restricted norm and largest spectral value are 7/12, giving the exact gap 5/12.

\bpdecls{\textbf{Lean.} \bpsource{GaussianTesting/Kac/GeneralRepresentation.lean}{42}{\texttt{SO3.kacAverage\_spectrum\_subset}}, \bpsource{GaussianTesting/Kac/GeneralRepresentation.lean}{36}{\texttt{SO3.kacAverage}}, \bpsource{GaussianTesting/Kac/SO3Representation.lean}{66}{\texttt{SO3.exists\_angularMomentum}}, \bpsource{GaussianTesting/Kac/CompactDecomposition.lean}{249}{\texttt{UnitaryRepresentation.irreducibleSpan\_dense\_haar}}, \bpsource{GaussianTesting/Kac/CompactDecomposition.lean}{255}{\texttt{UnitaryRepresentation.finiteInvariantVectors\_dense\_haar}}, \bpsource{GaussianTesting/Kac/GroupAverage.lean}{140}{\texttt{UnitaryRepresentation.groupKacAverage\_spectrum\_of\_finite}}, \bpsource{GaussianTesting/Kac/GroupAverage.lean}{42}{\texttt{UnitaryRepresentation.fixedProjection\_restrict}}, \bpsource{GaussianTesting/Kac/Defs.lean}{44}{\texttt{AngularMomentum}}, \bpsource{GaussianTesting/Kac/KacWalk.lean}{29}{\texttt{kacT\_gap}}, \bpsource{GaussianTesting/Kac/Gap.lean}{1332}{\texttt{AngularMomentum.kacGapOn\_of\_irreducible}}, \bpsource{GaussianTesting/Kac/Gap.lean}{1351}{\texttt{AngularMomentum.kacGapOn\_all}}, \bpsource{GaussianTesting/Kac/KacFixedSpace.lean}{33}{\texttt{UnitaryRepresentation.groupKacAverage\_eq\_self\_iff}}, \bpsource{GaussianTesting/Kac/KacFixedSpace.lean}{122}{\texttt{SO3.invariant\_iff\_axisRotation}}, \bpsource{GaussianTesting/Kac/KacFixedSpace.lean}{154}{\texttt{SO3.kacAverage\_eq\_self\_iff\_invariant}}, \bpsource{GaussianTesting/Kac/SphereTransitivity.lean}{102}{\texttt{SO3.sphereRotation\_transitive}}, \bpsource{GaussianTesting/Kac/SphereFixedSpace.lean}{32}{\texttt{SO3.sphereConstant}}, \bpsource{GaussianTesting/Kac/SphereFixedSpace.lean}{36}{\texttt{SO3.sphereConstants}}, \bpsource{GaussianTesting/Kac/SphereFixedSpace.lean}{47}{\texttt{SO3.mem\_sphereConstants}}, \bpsource{GaussianTesting/Kac/SphereFixedSpace.lean}{132}{\texttt{SO3.sphereRepresentation\_invariant\_iff\_constant}}, \bpsource{GaussianTesting/Kac/SphereFixedSpace.lean}{143}{\texttt{SO3.sphereKacTransition\_eq\_self\_iff\_constant}}, \bpsource{GaussianTesting/Kac/SphereFixedSpace.lean}{149}{\texttt{SO3.sphereKacTransition\_eq\_self\_iff\_ae\_constant}}, \bpsource{GaussianTesting/Kac/SphereFixedSpace.lean}{159}{\texttt{SO3.sphereKacTransition\_fixedSpace}}, \bpsource{GaussianTesting/Kac/SphereGapBound.lean}{85}{\texttt{SO3.sphereMeanZero}}, \bpsource{GaussianTesting/Kac/SphereGapBound.lean}{103}{\texttt{SO3.sphereMeanZeroKac}}, \bpsource{GaussianTesting/Kac/SphereGapBound.lean}{125}{\texttt{SO3.sphereMeanZeroKac\_norm\_le}}, \bpsource{GaussianTesting/Kac/SphereGapBound.lean}{131}{\texttt{SO3.sphereKacTransition\_norm\_le\_on\_orthogonal}}, \bpsource{GaussianTesting/Kac/SphereGapBound.lean}{137}{\texttt{SO3.sphereKacTransition\_quadratic\_le\_on\_orthogonal}}, \bpsource{GaussianTesting/Kac/SphereSupport.lean}{36}{\texttt{SO3.sphereContinuousToLp\_injective}}, \bpsource{GaussianTesting/Kac/SphereHarmonicQuartic.lean}{34}{\texttt{SO3.zonalFourPolynomial}}, \bpsource{GaussianTesting/Kac/SphereHarmonicQuartic.lean}{45}{\texttt{SO3.zonalFourPolynomial\_isHomogeneous}}, \bpsource{GaussianTesting/Kac/SphereHarmonicQuartic.lean}{52}{\texttt{SO3.zonalFourPolynomial\_laplacian}}, \bpsource{GaussianTesting/Kac/SphereGap.lean}{28}{\texttt{SO3.sphereZonalFour\_eq\_polynomial}}, \bpsource{GaussianTesting/Kac/SphereSpinFour.lean}{99}{\texttt{SO3.sphereZonalFour}}, \bpsource{GaussianTesting/Kac/SphereSpinFour.lean}{149}{\texttt{SO3.sphereAxisContinuousAverage\_zonalFour}}, \bpsource{GaussianTesting/Kac/SphereSpinFour.lean}{166}{\texttt{SO3.sphereSpinFour}}, \bpsource{GaussianTesting/Kac/SphereSpinFour.lean}{189}{\texttt{SO3.sphereKacTransition\_spinFour}}, \bpsource{GaussianTesting/Kac/SphereSpinFour.lean}{199}{\texttt{SO3.sphereSpinFour\_toLp\_ne\_zero}}, \bpsource{GaussianTesting/Kac/SphereGap.lean}{35}{\texttt{SO3.sphereSpinFourVector}}, \bpsource{GaussianTesting/Kac/SphereGap.lean}{44}{\texttt{SO3.sphereMeanZeroKac\_spinFour}}, \bpsource{GaussianTesting/Kac/SphereGap.lean}{51}{\texttt{SO3.sphereMeanZeroKac\_norm\_eq}}, \bpsource{GaussianTesting/Kac/SphereGap.lean}{62}{\texttt{SO3.sphereMeanZeroKac\_spectrum\_subset}}, \bpsource{GaussianTesting/Kac/SphereGap.lean}{82}{\texttt{SO3.sphereMeanZeroKac\_sharp\_mem\_spectrum}}, \bpsource{GaussianTesting/Kac/SphereGap.lean}{91}{\texttt{SO3.sphereMeanZeroKac\_isGreatest\_spectrum}}, \bpsource{GaussianTesting/Kac/SphereGap.lean}{102}{\texttt{SO3.sphereKac\_spectralGap}}, \bpsource{GaussianTesting/Kac/SpinOverlap.lean}{28}{\texttt{spinOverlap}}, \bpsource{GaussianTesting/Kac/SpinOverlap.lean}{38}{\texttt{alternating\_choose\_sq\_sum}}, \bpsource{GaussianTesting/Kac/SpinOverlap.lean}{58}{\texttt{spinOverlap\_odd}}, \bpsource{GaussianTesting/Kac/SpinOverlap.lean}{69}{\texttt{spinOverlap\_even}}, \bpsource{GaussianTesting/Kac/SpinOverlap.lean}{80}{\texttt{centralOverlap\_succ}}, \bpsource{GaussianTesting/Kac/SpinOverlap.lean}{93}{\texttt{centralOverlap\_strictAnti}}, \bpsource{GaussianTesting/Kac/SpinOverlap.lean}{101}{\texttt{spinKacMaximum}}, \bpsource{GaussianTesting/Kac/SpinOverlap.lean}{104}{\texttt{spinKacMaximum\_odd}}, \bpsource{GaussianTesting/Kac/SpinOverlap.lean}{113}{\texttt{spinKacMaximum\_two}}, \bpsource{GaussianTesting/Kac/SpinOverlap.lean}{116}{\texttt{spinKacMaximum\_four}}, \bpsource{GaussianTesting/Kac/SpinOverlap.lean}{126}{\texttt{spinKacMaximum\_lt\_of\_six\_le}}, \bpsource{GaussianTesting/ThreeVectorGram.lean}{31}{\texttt{threeVectorGram\_charpoly\_eval}}, \bpsource{GaussianTesting/ThreeVectorGram.lean}{39}{\texttt{threeVectorGram\_spectrum}}, \bpsource{GaussianTesting/Kac/Gap.lean}{1122}{\texttt{AngularMomentum.nonempty\_axisKernelVectors\_of\_integer\_spin}}, \bpsource{GaussianTesting/Kac/Gap.lean}{1217}{\texttt{AngularMomentum.kacT\_mem\_spectrum\_iff\_of\_integer\_spin}}, \bpsource{GaussianTesting/Kac/Gap.lean}{1226}{\texttt{AngularMomentum.kacT\_isGreatest\_eigenvalue\_of\_integer\_spin}}, \bpsource{GaussianTesting/Kac/Gap.lean}{1233}{\texttt{AngularMomentum.kacT\_eigenvalue\_lt\_of\_six\_le}}, \bpsource{GaussianTesting/Kac/Gap.lean}{1240}{\texttt{AngularMomentum.kacT\_quadratic\_le\_of\_integer\_spin}}.}

\textbf{Theorem node.} \bplink{thm:kac-walk}{Spectral gap for the Kac walk}.

\textbf{Related blueprint entries.} \bplink{prelim:kac}{Kac preliminaries}, \bplink{thm:main-fermions}{Fidelity bound for the fermionic test}, \bplink{thm:main-boson-mean-zero}{Fidelity bound for mean-zero Gaussian states}.

\section{Manuscript section 5. Fermions}
\subsection{Majoranas and Jordan-Wigner operators}\label{paper:fermion-majoranas}
\textbf{Checked.} Self-adjoint fermionic Majoranas obey the CAR and are realized by Pauli strings on occupation vectors. \bpsource{material/manuscript.tex}{733}{Manuscript, line 733.}

Checked in the manuscript's explicit Jordan-Wigner occupation coordinates: self-adjoint Majoranas, square-to-identity, anticommutation, and reconstruction c\_X=a†+a and c\_Y=i(a†−a) from the creation and annihilation operators. These are reusable background in the single fermionic preliminaries node. The checked exterior-algebra presentation is recorded separately and is not needed for Lemma 5.1.

\bpdecls{\textbf{Lean.} \bpsource{GaussianTesting/Fermions/JordanWigner.lean}{36}{\texttt{HF}}, \bpsource{GaussianTesting/Fermions/JordanWigner.lean}{53}{\texttt{maj}}, \bpsource{GaussianTesting/Fermions/JordanWigner.lean}{134}{\texttt{maj\_isSymmetric}}, \bpsource{GaussianTesting/Fermions/JordanWigner.lean}{146}{\texttt{maj\_sq}}, \bpsource{GaussianTesting/Fermions/JordanWigner.lean}{156}{\texttt{maj\_anticomm}}, \bpsource{GaussianTesting/Fermions/CreationAnnihilation.lean}{29}{\texttt{fermionAnnihilation}}, \bpsource{GaussianTesting/Fermions/CreationAnnihilation.lean}{33}{\texttt{fermionCreation}}, \bpsource{GaussianTesting/Fermions/CreationAnnihilation.lean}{36}{\texttt{fermionAnnihilation\_apply}}, \bpsource{GaussianTesting/Fermions/CreationAnnihilation.lean}{44}{\texttt{fermionCreation\_apply}}, \bpsource{GaussianTesting/Fermions/CreationAnnihilation.lean}{121}{\texttt{fermionCreation\_adjoint}}, \bpsource{GaussianTesting/Fermions/CreationAnnihilation.lean}{90}{\texttt{fermionCreation\_add\_annihilation}}, \bpsource{GaussianTesting/Fermions/CreationAnnihilation.lean}{98}{\texttt{imaginary\_fermionCreation\_sub\_annihilation}}.}

\textbf{Blueprint.} \bplink{prelim:fermions}{Fermionic Gaussian preliminaries}.

\subsection{Fermionic parity}\label{paper:fermion-parity}
\textbf{Checked.} Particle-number parity splits Fock space into even and odd sectors and anticommutes with each Majorana. \bpsource{material/manuscript.tex}{736}{Manuscript, line 736.}

Parity is checked in occupation coordinates: diagonal particle-number signs, involution, Majorana anticommutation, Gaussian invariance, and the manuscript's ordered product formula Pi=(-i)\textasciicircum{}m c\_1 ... c\_(2m). orderedMajoranaProduct groups consecutive X/Y factors by increasing mode; associativity preserves the displayed Majorana order. The formula includes zero modes.

\textbf{Proof route.} Each consecutive X/Y Majorana pair acts by i times the corresponding occupation sign. The ordered product therefore acts by i\textasciicircum{}m times particle-number parity; multiplying by (-i)\textasciicircum{}m gives Pi.

\bpdecls{\textbf{Lean.} \bpsource{GaussianTesting/Fermions/JordanWigner.lean}{178}{\texttt{parityOp}}, \bpsource{GaussianTesting/Fermions/JordanWigner.lean}{213}{\texttt{parityOp\_sq}}, \bpsource{GaussianTesting/Fermions/JordanWigner.lean}{205}{\texttt{parityOp\_anticomm\_maj}}, \bpsource{GaussianTesting/Fermions/ParityOrbits.lean}{27}{\texttt{occupationParity}}, \bpsource{GaussianTesting/Fermions/ParityOrbits.lean}{32}{\texttt{parityOp\_single}}, \bpsource{GaussianTesting/Fermions/GaussianImplementers.lean}{92}{\texttt{gaussianUnitary\_parity}}, \bpsource{GaussianTesting/Fermions/ExteriorOccupation.lean}{151}{\texttt{fermionExteriorEquiv}}, \bpsource{GaussianTesting/Fermions/CopyCoordinates.lean}{141}{\texttt{fermionTripleEquiv\_majorana}}, \bpsource{GaussianTesting/Fermions/ParityProduct.lean}{25}{\texttt{orderedMajoranaProduct}}, \bpsource{GaussianTesting/Fermions/ParityProduct.lean}{56}{\texttt{parityOp\_eq\_majoranaProduct}}.}

\textbf{Blueprint.} \bplink{prelim:fermions}{Fermionic Gaussian preliminaries}.

\subsection{Majorana rotations and the spin action (Eq. 5.1)}\label{paper:fermion-action}
\textbf{Checked.} An implementer for R in SO(2m) rotates the Majoranas by the columns of R and is unique projectively. \bpsource{material/manuscript.tex}{742}{Manuscript, line 742.}

Required background in the fermionic preliminaries node, in the manuscript's Jordan-Wigner coordinates. Every R in SO(2m) has a unitary implementer with the exact column action on the Majoranas. Finite coordinate-plane generation and half-angle lifts construct it. Two implementers of the same R differ by a unit scalar, because the Majorana commutant is scalar. Arbitrary Gaussian implementers are exactly the plane-rotation-and-phase group and preserve parity. No multiplicative choice of scalar phases or Spin double-cover construction is assumed.

\bpdecls{\textbf{Lean.} \bpsource{GaussianTesting/Fermions/GaussianAction.lean}{27}{\texttt{FermionRotation}}, \bpsource{GaussianTesting/Fermions/GaussianAction.lean}{40}{\texttt{ImplementsMajoranaRotation}}, \bpsource{GaussianTesting/Fermions/GaussianAction.lean}{97}{\texttt{gaussianUnitaryGroup}}, \bpsource{GaussianTesting/Fermions/GaussianImplementers.lean}{46}{\texttt{exists\_gaussianImplementer}}, \bpsource{GaussianTesting/Fermions/GaussianImplementers.lean}{77}{\texttt{gaussianImplementer}}, \bpsource{GaussianTesting/Fermions/GaussianImplementers.lean}{85}{\texttt{gaussianImplementer\_majorana}}, \bpsource{GaussianTesting/Fermions/GaussianImplementers.lean}{36}{\texttt{ImplementsMajoranaRotation.unique\_phase}}, \bpsource{GaussianTesting/Fermions/GaussianImplementers.lean}{62}{\texttt{gaussianUnitaryGroup\_eq\_planeUnitaryGroup}}, \bpsource{GaussianTesting/Fermions/GaussianImplementers.lean}{92}{\texttt{gaussianUnitary\_parity}}, \bpsource{GaussianTesting/Fermions/SpecialOrthogonalGeneration.lean}{238}{\texttt{specialOrthogonal\_induction}}, \bpsource{GaussianTesting/Fermions/PlaneImplementers.lean}{63}{\texttt{givensEquiv\_implements}}, \bpsource{GaussianTesting/Fermions/MajoranaCommutant.lean}{137}{\texttt{exists\_scalar\_of\_commutes\_maj}}.}

\textbf{Blueprint.} \bplink{prelim:fermions}{Fermionic Gaussian preliminaries}.

\subsection{Fermionic second quantization and copy actions}\label{paper:fermion-second-quantization}
\textbf{Partial coverage.} On each exterior-power sector, second quantization restricts the tensor power of a one-particle unitary; copy rotations act through this construction. \bpsource{material/manuscript.tex}{749}{Manuscript, line 749.}

The real special orthogonal mode rotations needed by Theorem 5.2 are fully constructed: vacuum normalization gives an actual unitary representation, its action is identified with ExteriorAlgebra.map and with the induced map on each exterior-power sector, and it is strongly continuous. This required background belongs to the fermionic preliminaries node. The broader complex SU(m) statement at this manuscript location is not exported and is not needed by the proof. The general t-copy O(t) action and Howe commutant statement added later in Section 5 are recorded under Orthogonal copy symmetries; the needed SO(2) and SO(3) actions remain sufficient.

\bpdecls{\textbf{Lean.} \bpsource{GaussianTesting/Fermions/SecondQuantization.lean}{155}{\texttt{fermionSecondQuantization}}, \bpsource{GaussianTesting/Fermions/SecondQuantization.lean}{169}{\texttt{fermionSecondQuantization\_implements}}, \bpsource{GaussianTesting/Fermions/SecondQuantization.lean}{177}{\texttt{fermionSecondQuantization\_creation}}, \bpsource{GaussianTesting/Fermions/SecondQuantization.lean}{224}{\texttt{continuous\_fermionSecondQuantization\_orbit}}, \bpsource{GaussianTesting/Fermions/ExteriorOccupation.lean}{151}{\texttt{fermionExteriorEquiv}}, \bpsource{GaussianTesting/Fermions/ExteriorOccupation.lean}{206}{\texttt{fermionSecondQuantization\_eq\_exterior\_map}}, \bpsource{GaussianTesting/Fermions/ExteriorOccupation.lean}{218}{\texttt{fermionSecondQuantization\_exteriorPower}}, \bpsource{GaussianTesting/Fermions/ExteriorOccupation.lean}{228}{\texttt{fermionSecondQuantization\_wedge}}.}

\textbf{Blueprint.} \bplink{prelim:fermions}{Fermionic Gaussian preliminaries}.

\subsection{Even and odd fermionic Gaussian classes}\label{paper:fermion-gaussian}
\textbf{Checked.} Within each parity sector the Gaussian states form one orbit of the SO(2m) Majorana implementers; the union of the even and odd orbits equals IsGaussian. \bpsource{material/manuscript.tex}{752}{Manuscript, line 752.}

Required orbit identification in the fermionic preliminaries node. The Gaussian unitary group is independently defined by its SO(2m) Majorana action, not by IsGaussian or the certainty criterion. Its orbit from any occupation vector x is precisely the IsGaussian states with parity (-1)\textasciicircum{}|x|. Two-bit flips give transitivity on occupation vectors of the same parity; phase uniqueness identifies all implementers with the explicit generated group. The even class is the vacuum orbit, any odd occupation vector generates the odd class, and their union is IsGaussian. With zero modes the odd class is empty. No covariance, kernel or robustness classification enters this proof. The unused quadratic-ground-state equivalence is outside this checked orbit statement.

\bpdecls{\textbf{Lean.} \bpsource{GaussianTesting/Fermions/GaussianOrbits.lean}{28}{\texttt{fermionicGaussianOrbit}}, \bpsource{GaussianTesting/Fermions/GaussianOrbits.lean}{37}{\texttt{evenFermionicGaussianStates}}, \bpsource{GaussianTesting/Fermions/GaussianOrbits.lean}{41}{\texttt{oddFermionicGaussianStates}}, \bpsource{GaussianTesting/Fermions/GaussianOrbits.lean}{45}{\texttt{fermionicGaussianStates}}, \bpsource{GaussianTesting/Fermions/GaussianOrbits.lean}{50}{\texttt{mem\_fermionicGaussianOrbit\_iff}}, \bpsource{GaussianTesting/Fermions/GaussianOrbits.lean}{58}{\texttt{fermionicGaussianOrbit\_eq\_parityClass}}, \bpsource{GaussianTesting/Fermions/GaussianOrbits.lean}{68}{\texttt{evenFermionicGaussianStates\_eq\_vacuum\_orbit}}, \bpsource{GaussianTesting/Fermions/GaussianOrbits.lean}{74}{\texttt{oddFermionicGaussianStates\_eq\_orbit}}, \bpsource{GaussianTesting/Fermions/GaussianOrbits.lean}{94}{\texttt{fermionicGaussianStates\_eq\_isGaussian}}, \bpsource{GaussianTesting/Fermions/GaussianOrbits.lean}{100}{\texttt{oddFermionicGaussianStates\_zero}}, \bpsource{GaussianTesting/Fermions/ParityOrbits.lean}{73}{\texttt{exists\_planeUnitary\_basis\_of\_parity\_eq}}, \bpsource{GaussianTesting/Fermions/Gaussian.lean}{42}{\texttt{IsGaussian}}.}

\textbf{Blueprint.} \bplink{prelim:fermions}{Fermionic Gaussian preliminaries}.

\subsection{Fermionic occupation-sign operators}\label{paper:fermion-number}
\textbf{Checked.} The Majorana pair i c\_(2k-1) c\_(2k) is 2N\_k-I, diagonal with eigenvalue 2x\_k-1 on occupation vectors. \bpsource{material/manuscript.tex}{756}{Manuscript, line 756.}

Checked in the manuscript's Jordan-Wigner coordinates: fermionNumber=a†a, its occupation eigenvalues, and i c\_X c\_Y=2N−I. nOp denotes the signed observable 2N−I, not N itself. The expectation-saturation argument belongs to Lemma 5.1. No exterior-Fock identification is needed for that argument.

\bpdecls{\textbf{Lean.} \bpsource{GaussianTesting/Fermions/CreationAnnihilation.lean}{69}{\texttt{fermionNumber}}, \bpsource{GaussianTesting/Fermions/CreationAnnihilation.lean}{72}{\texttt{fermionNumber\_apply}}, \bpsource{GaussianTesting/Fermions/CreationAnnihilation.lean}{81}{\texttt{fermionNumber\_single}}, \bpsource{GaussianTesting/Fermions/CreationAnnihilation.lean}{107}{\texttt{majorana\_pair\_eq\_two\_fermionNumber\_sub\_id}}, \bpsource{GaussianTesting/Fermions/CovarianceBasic.lean}{97}{\texttt{nOp}}, \bpsource{GaussianTesting/Fermions/CovarianceBasic.lean}{100}{\texttt{nOp\_apply}}.}

\textbf{Blueprint.} \bplink{prelim:fermions}{Fermionic Gaussian preliminaries}.

\subsection{Fermionic two-copy operator and its kernel projector}\label{paper:lambda-f}
\textbf{Checked.} Lambda is the sum of the two-copy Majorana products, equivalently twice the indicated sum of ladder-operator products. P projects onto ker Lambda. \bpsource{material/manuscript.tex}{778}{Manuscript, line 778.}

The fermionic test-construction node contains the Majorana and creation/annihilation expressions for Lambda, its symmetry and swap commutation, and its orthogonal kernel acceptance projection in the manuscript's occupation coordinates. Perfect acceptance is equivalent to the kernel condition for unit inputs.

\bpdecls{\textbf{Lean.} \bpsource{GaussianTesting/Fermions/Lambda.lean}{138}{\texttt{lambdaF}}, \bpsource{GaussianTesting/Fermions/CreationAnnihilation.lean}{151}{\texttt{lambdaF\_eq\_ladder}}, \bpsource{GaussianTesting/Fermions/Lambda.lean}{164}{\texttt{PF}}, \bpsource{GaussianTesting/Fermions/Lambda.lean}{142}{\texttt{lambdaF\_isSymmetric}}, \bpsource{GaussianTesting/Fermions/Lambda.lean}{152}{\texttt{lambdaF\_swap\_comm}}, \bpsource{GaussianTesting/Fermions/Lambda.lean}{168}{\texttt{PF\_isAcceptanceProjection}}, \bpsource{GaussianTesting/Fermions/Lambda.lean}{191}{\texttt{accProb\_PF\_eq\_one\_iff}}.}

\textbf{Blueprint.} \bplink{def:PF}{Fermionic test projection}.

\subsection{Simultaneous Gaussian conjugation of Lambda (Eq. 5.2)}\label{paper:fermion-invariance}
\textbf{Checked.} Applying the same fermionic Gaussian unitary to both copies leaves Lambda invariant. \bpsource{material/manuscript.tex}{786}{Manuscript, line 786.}

Equation (5.2) is proved for every implementer of any SO(2m) rotation, on the full two-copy Hilbert space. The proof expands the Majorana action and uses orthogonality of the rotation. The kernel projector commutes with the action, and acceptance is invariant. These test-specific calculations belong to the test-construction node.

\bpdecls{\textbf{Lean.} \bpsource{GaussianTesting/Fermions/GaussianTestInvariance.lean}{27}{\texttt{fermionPairAction}}, \bpsource{GaussianTesting/Fermions/GaussianTestInvariance.lean}{77}{\texttt{ImplementsMajoranaRotation.lambdaF\_comm}}, \bpsource{GaussianTesting/Fermions/GaussianTestInvariance.lean}{112}{\texttt{ImplementsMajoranaRotation.lambdaF\_conjugate}}, \bpsource{GaussianTesting/Fermions/GaussianTestInvariance.lean}{119}{\texttt{ImplementsMajoranaRotation.PF\_comm}}, \bpsource{GaussianTesting/Fermions/GaussianTestInvariance.lean}{142}{\texttt{ImplementsMajoranaRotation.accProb\_PF}}.}

\textbf{Blueprint.} \bplink{def:PF}{Fermionic test projection}.

\subsection{Lemma 5.1}\label{paper:r5-1}
\textbf{Checked.} For any pure fermionic input, with no parity restriction, Gaussianity is equivalent to annihilation of its tensor square by Lambda and to perfect acceptance. \bpsource{material/manuscript.tex}{793}{Manuscript, line 793.}

Both equivalences hold for the independently defined union of the even and odd SO(2m) Gaussian orbits, for every unit input, including zero modes and without any input parity assumption. The proof follows the manuscript: Gaussian invariance and basis-state annihilation; the covariance norm identity; general SO antisymmetric block normal form and Gaussian covariance transport; block-coefficient bounds and saturation to a joint occupation eigenstate. No open obligations remain for this lemma in the manuscript's explicit occupation presentation.

\textbf{Proof route.} The manuscript's normal-form proof is used directly: saturate the covariance norm, block-diagonalize the arbitrary real antisymmetric matrix by an SO rotation, implement it by a Gaussian unitary, and deduce that every expectation of 2N\_k−I is ±1. The rotated state is an occupation basis vector up to phase.

\bpdecls{\textbf{Lean.} \bpsource{GaussianTesting/Fermions/GaussianCertainty.lean}{65}{\texttt{fermionicGaussian\_iff\_lambdaF}}, \bpsource{GaussianTesting/Fermions/GaussianCertainty.lean}{70}{\texttt{fermionicGaussian\_iff\_accProb}}, \bpsource{GaussianTesting/Fermions/GaussianTestInvariance.lean}{149}{\texttt{fermionicGaussian\_lambdaF\_eq\_zero}}, \bpsource{GaussianTesting/Fermions/GaussianCertainty.lean}{40}{\texttt{fermionicGaussian\_of\_lambdaF\_eq\_zero}}, \bpsource{GaussianTesting/Fermions/GaussianCertainty.lean}{30}{\texttt{sum\_sq\_covMat\_of\_lambdaF\_eq\_zero}}, \bpsource{GaussianTesting/Fermions/CovarianceOccupation.lean}{109}{\texttt{exists\_phase\_basis\_of\_covMat\_skewBlockMatrix}}, \bpsource{GaussianTesting/Fermions/CovarianceOccupation.lean}{71}{\texttt{abs\_covMat\_pair\_le\_one}}, \bpsource{GaussianTesting/Fermions/CovarianceOccupation.lean}{81}{\texttt{covMat\_pair\_signs\_of\_sum\_eq}}, \bpsource{GaussianTesting/Fermions/CovarianceBasic.lean}{152}{\texttt{nOp\_eigen\_of\_expectation}}, \bpsource{GaussianTesting/Fermions/CovarianceBasic.lean}{179}{\texttt{joint\_eigen\_basis}}.}

\textbf{Blueprint.} \bplink{lem:exact-fermions}{Fermionic states accepted with certainty}, \bplink{thm:main-fermions}{Fidelity bound for the fermionic test}, \bplink{cor:fermion-testing}{Testing fermionic Gaussian states}.

\subsection{Covariance saturation identity (Eq. 5.3)}\label{paper:fermion-covariance}
\textbf{Checked.} The squared norm of Lambda(psi tensor psi) is 2m minus the Hilbert-Schmidt norm squared of the covariance matrix. \bpsource{material/manuscript.tex}{806}{Manuscript, line 806.}

Equation (5.3) and its saturation consequences are proof steps of Lemma 5.1. The squared-entry sum is exactly the Hilbert-Schmidt norm squared. Gaussian covariance congruence and orthogonal norm invariance transport the saturated norm to canonical blocks; the block coefficients then force an occupation basis state. These applications stay in the lemma node, not as assumptions in the preliminaries.

\bpdecls{\textbf{Lean.} \bpsource{GaussianTesting/Fermions/CovarianceBasic.lean}{23}{\texttt{majMoment}}, \bpsource{GaussianTesting/Fermions/CovarianceBasic.lean}{26}{\texttt{covMat}}, \bpsource{GaussianTesting/Fermions/CovarianceBasic.lean}{69}{\texttt{norm\_lambdaF\_tp2}}, \bpsource{GaussianTesting/Fermions/GaussianCertainty.lean}{30}{\texttt{sum\_sq\_covMat\_of\_lambdaF\_eq\_zero}}, \bpsource{GaussianTesting/Fermions/CovarianceAction.lean}{58}{\texttt{ImplementsMajoranaRotation.covMat\_symm}}, \bpsource{GaussianTesting/Fermions/CovarianceAction.lean}{106}{\texttt{ImplementsMajoranaRotation.sum\_sq\_covMat\_symm}}, \bpsource{GaussianTesting/Fermions/CovarianceOccupation.lean}{41}{\texttt{covMat\_pair\_sum\_eq\_of\_skewBlockMatrix}}, \bpsource{GaussianTesting/Fermions/CovarianceOccupation.lean}{109}{\texttt{exists\_phase\_basis\_of\_covMat\_skewBlockMatrix}}.}

\textbf{Blueprint.} \bplink{lem:exact-fermions}{Fermionic states accepted with certainty}.

\subsection{Real antisymmetric block normal form}\label{paper:fermion-skew-normal-form}
\textbf{Checked.} Every real antisymmetric matrix of even dimension has paired two-dimensional skew blocks after an SO change of basis. \bpsource{material/manuscript.tex}{810}{Manuscript, line 810.}

Standard linear algebra used by Lemma 5.1 and collected in the fermionic preliminaries node. The theorem applies to arbitrary real antisymmetric matrices, with no covariance saturation or state-class assumptions. The proof uses the real spectral theorem for the symmetric operator −A², invariant orthogonal complements, and an orientation correction; it includes zero dimension.

\bpdecls{\textbf{Lean.} \bpsource{GaussianTesting/Fermions/SkewNormalFormOrientation.lean}{28}{\texttt{skewBlockMatrix}}, \bpsource{GaussianTesting/Fermions/SkewNormalForm.lean}{210}{\texttt{exists\_specialOrthogonal\_skewNormalForm}}, \bpsource{GaussianTesting/Fermions/CovarianceAction.lean}{86}{\texttt{sum\_sq\_specialOrthogonal\_congruence}}.}

\textbf{Blueprint.} \bplink{prelim:fermions}{Fermionic Gaussian preliminaries}.

\subsection{Fermionic kernel is spanned by Gaussian squares}\label{paper:fermion-kernel-span}
\textbf{Checked.} The cited Gaussian tensor squares span the full fermionic kernel, so P has support in the symmetric subspace. \bpsource{material/manuscript.tex}{823}{Manuscript, line 823.}

For every mode count, the full kernel of lambdaF is the linear span of tensor squares from the independently defined union of both Gaussian parity orbits. Its projection PF is supported on the swap-symmetric subspace on arbitrary two-copy inputs. The proof uses Lemma 5.1, the joint Bell eigenlines of the Majorana pair involutions, a vacuum-square overlap, and signed coordinate permutations. Both parities are essential for invariance under individual Majoranas. The manuscript specifies twirling over the full orthogonal group, which includes both parity classes. Lean uses the previously agreed direct invariant-span argument, not a separately constructed Haar average. This assertion is recorded with the introductory span and optimality claim, not as an additional obligation of Lemma 5.1.

\bpdecls{\textbf{Lean.} \bpsource{GaussianTesting/Fermions/GaussianSquareSpan.lean}{36}{\texttt{fermionicGaussianSquareSpan}}, \bpsource{GaussianTesting/Fermions/GaussianSquareSpan.lean}{39}{\texttt{fermionicGaussianSquareSpan\_le\_kernel}}, \bpsource{GaussianTesting/Fermions/GaussianSquareSpan.lean}{426}{\texttt{fermionicGaussianSquareSpan\_eq\_kernel}}, \bpsource{GaussianTesting/Fermions/GaussianSquareSpan.lean}{446}{\texttt{swap2\_eq\_self\_of\_mem\_fermionicGaussianSquareSpan}}, \bpsource{GaussianTesting/Fermions/GaussianSquareSpan.lean}{458}{\texttt{swap2\_PF}}.}

\textbf{Blueprint.} \bplink{prelim:fermions}{Fermionic Gaussian preliminaries}, \bplink{target:span-optimality}{Gaussian-square span and optimality}.

\subsection{Two-copy parity strings (Eq. 5.4)}\label{paper:fermion-two-copy-majoranas}
\textbf{Checked.} The ordinary tensor coordinates acquire parity operators on earlier copies to recover Majoranas satisfying the CAR on the enlarged system. \bpsource{material/manuscript.tex}{830}{Manuscript, line 830.}

The actual occupation unitary fermionPairEquiv identifies HF(2*m) with the two-copy Hilbert space. With modes ordered by copy then mode, it transports the Jordan-Wigner Majoranas to exactly c\_j tensor I and Pi tensor c\_j. The transported operators are Hermitian, square to identity, and distinct ones anticommute. The three-copy counterpart used in Theorem 5.2 remains checked independently.

\textbf{Proof route.} Split a 2m-mode occupation vector into its first and second m-mode blocks. The Jordan-Wigner prefix factors into parity of preceding copies and the within-copy prefix; bit flips respect the splitting. Transport through the resulting occupation unitary proves the two explicit formulas, Hermiticity and the CAR.

\bpdecls{\textbf{Lean.} \bpsource{GaussianTesting/Fermions/PairCoordinates.lean}{34}{\texttt{fermionPairOccupationEquiv}}, \bpsource{GaussianTesting/Fermions/PairCoordinates.lean}{54}{\texttt{fermionPairEquiv}}, \bpsource{GaussianTesting/Fermions/PairCoordinates.lean}{115}{\texttt{maj2\textsubscript{1}}}, \bpsource{GaussianTesting/Fermions/PairCoordinates.lean}{118}{\texttt{maj2\textsubscript{2}}}, \bpsource{GaussianTesting/Fermions/PairCoordinates.lean}{122}{\texttt{fermionPairMajorana}}, \bpsource{GaussianTesting/Fermions/PairCoordinates.lean}{134}{\texttt{fermionPairEquiv\_majorana}}, \bpsource{GaussianTesting/Fermions/PairCoordinates.lean}{183}{\texttt{fermionPairMajorana\_isSymmetric}}, \bpsource{GaussianTesting/Fermions/PairCoordinates.lean}{153}{\texttt{fermionPairMajorana\_sq}}, \bpsource{GaussianTesting/Fermions/PairCoordinates.lean}{163}{\texttt{fermionPairMajorana\_anticomm}}.}

\textbf{Blueprint.} \bplink{prelim:fermions}{Fermionic Gaussian preliminaries}.

\subsection{Two-copy Majoranas and the rotation generator (Eq. 5.5)}\label{paper:fermion-rotation}
\textbf{Checked.} The parity-twisted Majoranas give L = -i (Pi tensor I) Lambda. L has the same kernel as Lambda and generates copy rotations, up to angular normalization. \bpsource{material/manuscript.tex}{837}{Manuscript, line 837.}

The standalone two-copy statement is fully checked. fermionPairL is the manuscript's i sum\_j c\_(j,1)c\_(j,2), equals -i(Pi tensor I)Lambda, and has exactly ker Lambda. fermionPairRepresentation is the actual strongly continuous exterior second quantization of R tensor I\_m for R in SO(2), transported by the occupation unitary. Its usual-angle derivative is i L/2. Full SO(2) invariance is equivalent to Lambda v=0, so PF is exactly the invariant-space projection. The three-copy axis generators and Kac equality remain proof steps of Theorem 5.2. The unused complex SU(m) action is not added.

\textbf{Proof route.} Parity anticommutation proves L=-i(Pi tensor I)Lambda; parity's involution gives equal kernels. The actual SO(2) second-quantized action rotates the two Majorana families and fixes the vacuum. Differentiating covariance determines the generator through its commutators; the scalar commutant and fixed vacuum remove the scalar ambiguity. The existing finite-dimensional one-parameter fixed-space theorem identifies the generator kernel. Coordinate-plane generation extends fixedness at every angle to full SO(2), and equality of subspaces identifies the orthogonal projection.

\bpdecls{\textbf{Lean.} \bpsource{GaussianTesting/Fermions/PairGenerator.lean}{28}{\texttt{fermionPairL}}, \bpsource{GaussianTesting/Fermions/PairGenerator.lean}{37}{\texttt{fermionPairL\_eq}}, \bpsource{GaussianTesting/Fermions/PairGenerator.lean}{50}{\texttt{ker\_fermionPairL}}, \bpsource{GaussianTesting/Fermions/PairGenerator.lean}{67}{\texttt{fermionPairGenerator}}, \bpsource{GaussianTesting/Fermions/PairGenerator.lean}{71}{\texttt{ker\_fermionPairGenerator}}, \bpsource{GaussianTesting/Fermions/PairRepresentation.lean}{29}{\texttt{fermionPairModeRotation}}, \bpsource{GaussianTesting/Fermions/PairRepresentation.lean}{82}{\texttt{fermionPairRepresentation}}, \bpsource{GaussianTesting/Fermions/PairRepresentation.lean}{93}{\texttt{continuous\_fermionPairRepresentation\_orbit}}, \bpsource{GaussianTesting/Fermions/PairRepresentation.lean}{112}{\texttt{fermionPairRepresentation\_majorana}}, \bpsource{GaussianTesting/Fermions/PairRotation.lean}{28}{\texttt{fermionPairCircle}}, \bpsource{GaussianTesting/Fermions/PairRotation.lean}{88}{\texttt{fermionPairRepresentation\_hasDerivAt}}, \bpsource{GaussianTesting/Fermions/PairRotation.lean}{120}{\texttt{fermionPairCircle\_fixed\_iff}}, \bpsource{GaussianTesting/Fermions/PairRotation.lean}{131}{\texttt{fermionPairRepresentation\_fixed\_iff}}, \bpsource{GaussianTesting/Fermions/PairRotation.lean}{156}{\texttt{fermionPairRepresentation\_fixedSubspace}}, \bpsource{GaussianTesting/Fermions/PairRotation.lean}{166}{\texttt{PF\_eq\_pairRotation\_fixedProjection}}.}

\textbf{Blueprint.} \bplink{prelim:fermions}{Fermionic Gaussian preliminaries}, \bplink{thm:main-fermions}{Fidelity bound for the fermionic test}.

\subsection{Theorem 5.2}\label{paper:r5-2}
\textbf{Checked.} Acceptance p >= 1 - tau implies squared fidelity at least 1 - 8 tau to the fermionic Gaussian class. \bpsource{material/manuscript.tex}{856}{Manuscript, line 856.}

For every mode count and unit input, acceptance at least 1-tau gives a witness in the independent even/odd Gaussian orbit union, squared overlap at least 1-8 tau, and the same bound for squared supremal fidelity. The actual strongly continuous second-quantized SO(3) action has the prescribed parity-string Majoranas. Its axis-fixed spaces are exactly the pair-test kernels, giving S=T\_pi. Theorem 4.1 supplies the complex spectral inclusion; the completed-tensor/projective-flow proof of Theorem 3.4 supplies the witness; Lemma 5.1 identifies its class. Nonnegativity of tau follows from the acceptance hypothesis. No parity restriction, positive-mode hypothesis, or extra spectral assumption is supplied.

\textbf{Proof route.} The Lean proof follows the manuscript: construct the second-quantized SO(3) action on three copies; identify the coordinate-axis fixed projections with P12, P13 and P23 using the parity-string generators; obtain S=T\_pi and apply Theorem 4.1; transport finite coordinates to completed tensor products and apply the projective Theorem 3.4; use Lemma 5.1 for Gaussianity and Eq. (3.2) for class fidelity.

\bpdecls{\textbf{Lean.} \bpsource{GaussianTesting/Fermions/Robustness.lean}{39}{\texttt{fermionic\_fidelity\_witness}}, \bpsource{GaussianTesting/Fermions/Robustness.lean}{63}{\texttt{fermionic\_fidelity}}, \bpsource{GaussianTesting/Fermions/Robustness.lean}{27}{\texttt{fermionicGaussianStates\_eq\_certaintyStates}}, \bpsource{GaussianTesting/Fermions/CopyRotationAxis.lean}{40}{\texttt{fermion\_copyRepresentation\_axis\_derivatives}}, \bpsource{GaussianTesting/Fermions/CopyKac.lean}{26}{\texttt{fermion\_copy\_axis\_fixed}}, \bpsource{GaussianTesting/Fermions/CopyKac.lean}{38}{\texttt{fermion\_Sop\_eq\_kacAverage}}, \bpsource{GaussianTesting/Fermions/CopyKac.lean}{47}{\texttt{fermionic\_spectrum\_subset}}, \bpsource{GaussianTesting/GradientFlow/FiniteHilbert.lean}{200}{\texttt{finite\_robustness\_of\_spectrum}}.}

\textbf{Theorem node.} \bplink{thm:main-fermions}{Fidelity bound for the fermionic test}.

\textbf{Related blueprint entries.} \bplink{cor:fermion-testing}{Testing fermionic Gaussian states}.

\subsection{Second-quantized SO(3) copy action}\label{paper:fermion-copy-action}
\textbf{Checked.} The action Gamma\_wedge(R tensor I\_m) is a strongly continuous unitary SO(3) representation on three fermionic copies. \bpsource{material/manuscript.tex}{874}{Manuscript, line 874.}

The required real copy action is fully constructed in the fermionic preliminaries. A vacuum-fixed implementer realizes second quantization on every exterior-power sector. The mode embedding is exactly R tensor I\_m, and the copy-major occupation unitary transports the Majoranas to the manuscript parity strings. The test-specific derivatives, fixed projections, and Kac equality belong to Theorem 5.2.

\bpdecls{\textbf{Lean.} \bpsource{GaussianTesting/Fermions/CopyModeRotation.lean}{29}{\texttt{fermionCopyModeRotation}}, \bpsource{GaussianTesting/Fermions/CopyRepresentation.lean}{26}{\texttt{fermionCopyRepresentation}}, \bpsource{GaussianTesting/Fermions/CopyRepresentation.lean}{37}{\texttt{continuous\_fermionCopyRepresentation\_orbit}}, \bpsource{GaussianTesting/Fermions/CopyRepresentation.lean}{44}{\texttt{fermionCopyRepresentation\_vacuum}}, \bpsource{GaussianTesting/Fermions/CopyRepresentation.lean}{56}{\texttt{fermionCopyRepresentation\_majorana}}, \bpsource{GaussianTesting/Fermions/ExteriorOccupation.lean}{218}{\texttt{fermionSecondQuantization\_exteriorPower}}.}

\textbf{Blueprint.} \bplink{prelim:fermions}{Fermionic Gaussian preliminaries}.

\subsection{Corollary 5.3}\label{paper:r5-3}
\textbf{Checked.} For epsilon, delta in (0,1), repeated measurements on fresh pairs have perfect completeness and rejection probability at least 1-delta when squared Gaussian fidelity is at most 1-epsilon², with mode-independent O(epsilon\textasciicircum{}(-2) log(1/delta)) copy complexity. \bpsource{material/manuscript.tex}{888}{Manuscript, line 888.}

The protocol uses the independent fermionic Gaussian orbit class and the Born probability of PF on each fresh pair. The product outcome law has exact acceptance p\textasciicircum{}n and rejection 1-p\textasciicircum{}n. It has perfect completeness and soundness for the manuscript’s squared-fidelity promise, with n=ceil((8/epsilon²)*log(1/delta)) and exactly 2n input copies. The uniform bound is 16/epsilon²*log(1/delta)+2; for delta <= exp(-1), the copy count is at most 32/epsilon²*log(1/delta). The Bell-readout identity is checked separately in Figure 1(a).

\textbf{Proof route.} Follow the manuscript: Lemma 5.1 gives p=1 for Gaussian states; Theorem 5.2 applied with tau=1-p gives p<=1-epsilon²/8 under the class-fidelity promise. Independent repetition on fresh pairs gives p\textasciicircum{}n, and the exponential estimate with the chosen ceiling bounds this by delta. The complement event gives rejection at least 1-delta. Count two input copies per round and retain the ceiling in the uniform estimate.

\bpdecls{\textbf{Lean.} \bpsource{GaussianTesting/Fermions/Testing.lean}{63}{\texttt{fermionic\_testing}}, \bpsource{GaussianTesting/Fermions/Testing.lean}{37}{\texttt{fermionicTestProtocol}}, \bpsource{GaussianTesting/Fermions/Testing.lean}{28}{\texttt{fermionic\_acceptance\_le\_of\_far}}, \bpsource{GaussianTesting/Fermions/Testing.lean}{46}{\texttt{fermionicTestProtocol\_acceptance}}, \bpsource{GaussianTesting/Fermions/Testing.lean}{53}{\texttt{fermionicTestProtocol\_rejection}}, \bpsource{GaussianTesting/Fermions/Testing.lean}{43}{\texttt{fermionicTestCopies}}, \bpsource{GaussianTesting/Fermions/Testing.lean}{97}{\texttt{fermionicTestCopies\_le\_rate}}, \bpsource{GaussianTesting/RepeatedTest.lean}{95}{\texttt{testingRounds}}, \bpsource{GaussianTesting/Fermions/GaussianCertainty.lean}{70}{\texttt{fermionicGaussian\_iff\_accProb}}.}

\textbf{Blueprint.} \bplink{cor:fermion-testing}{Testing fermionic Gaussian states}, \bplink{target:measurement-endpoints}{Testing Gaussianity with constant sample complexity}.

\section{Manuscript section 6. Bosons}
\subsection{Bosonic Fock space and number basis}\label{paper:boson-fock}
\textbf{Checked.} The normalized oscillator number vectors form a complete orthonormal basis of L2(R\textasciicircum{}m), indexed by occupation tuples; their coefficients give a unitary identification with the number-state Fock model. \bpsource{material/manuscript.tex}{904}{Manuscript, line 904.}

Ordered creation powers have the manuscript normalization and satisfy the sqrt(n+1)/sqrt(n) ladder formulas. Gaussian integration by parts proves adjointness, the canonical commutator gives orthogonality and factorial norms, and Fourier uniqueness proves density of their polynomial Gaussian span. schrodingerNumberBasis is a complete HilbertBasis, fockSchrodingerEquiv is the actual occupation-to-wavefunction unitary, and the vacuum corresponds to zero occupation. For two copies the finite and closed generator graphs, kernels, and accepting projections correspond exactly. The paired unitary maps completed product tensors to products of the single-copy wavefunctions. The three-copy identifications are tracked separately where needed by the tests.

\bpdecls{\textbf{Lean.} \bpsource{GaussianTesting/Bosons/Fock.lean}{32}{\texttt{FockB}}, \bpsource{GaussianTesting/Bosons/Fock.lean}{35}{\texttt{degB}}, \bpsource{GaussianTesting/Bosons/Fock.lean}{61}{\texttt{BlockSp}}, \bpsource{GaussianTesting/Bosons/Lambda.lean}{53}{\texttt{aF}}, \bpsource{GaussianTesting/Bosons/Lambda.lean}{58}{\texttt{adF}}, \bpsource{GaussianTesting/Bosons/SchrodingerTensor.lean}{34}{\texttt{WaveSpace}}, \bpsource{GaussianTesting/Bosons/SchrodingerRotation.lean}{33}{\texttt{SchrodingerSpace}}, \bpsource{GaussianTesting/Bosons/SchrodingerCore.lean}{147}{\texttt{fockNumberPolynomial}}, \bpsource{GaussianTesting/Bosons/SchrodingerCore.lean}{154}{\texttt{finiteFockPolynomialCore}}, \bpsource{GaussianTesting/Bosons/SchrodingerDerivative.lean}{84}{\texttt{polynomialGaussianFun\_annihilation}}, \bpsource{GaussianTesting/Bosons/SchrodingerDerivative.lean}{96}{\texttt{polynomialGaussianFun\_creation}}, \bpsource{GaussianTesting/Bosons/NumberPolynomial.lean}{53}{\texttt{polynomialAnnihilation\_creation}}, \bpsource{GaussianTesting/Bosons/NumberPolynomial.lean}{200}{\texttt{finiteFockPolynomialCore\_eq\_top}}, \bpsource{GaussianTesting/Bosons/PolynomialGaussianDensity.lean}{183}{\texttt{denseRange\_polynomialGaussian}}, \bpsource{GaussianTesting/Bosons/SchrodingerAdjoint.lean}{183}{\texttt{inner\_polynomialGaussian\_creation}}, \bpsource{GaussianTesting/Bosons/SchrodingerAdjoint.lean}{200}{\texttt{inner\_polynomialGaussian\_annihilation}}, \bpsource{GaussianTesting/Bosons/NumberCoreDense.lean}{38}{\texttt{schrodingerFiniteLambda\_domain\_dense}}, \bpsource{GaussianTesting/Bosons/NumberCoreDense.lean}{44}{\texttt{schrodingerLambda\_domain\_dense}}, \bpsource{GaussianTesting/Bosons/SchrodingerNumberBasis.lean}{145}{\texttt{schrodingerNumberBasis}}, \bpsource{GaussianTesting/Bosons/SchrodingerNumberBasis.lean}{106}{\texttt{orthonormal\_schrodingerNumberVector}}, \bpsource{GaussianTesting/Bosons/SchrodingerNumberBasis.lean}{140}{\texttt{dense\_span\_schrodingerNumberVector}}, \bpsource{GaussianTesting/Bosons/SchrodingerNumberBasis.lean}{165}{\texttt{fockSchrodingerEquiv}}, \bpsource{GaussianTesting/Bosons/SchrodingerNumberBasis.lean}{186}{\texttt{fockSchrodingerEquiv\_vacuum}}, \bpsource{GaussianTesting/Bosons/NumberLadder.lean}{98}{\texttt{polynomialCreation\_schrodingerNumberPolynomial}}, \bpsource{GaussianTesting/Bosons/NumberLadder.lean}{108}{\texttt{polynomialAnnihilation\_schrodingerNumberPolynomial}}, \bpsource{GaussianTesting/Bosons/SchrodingerFockPair.lean}{56}{\texttt{fockPairSchrodingerEquiv}}, \bpsource{GaussianTesting/Bosons/SchrodingerFockPair.lean}{256}{\texttt{fockPairGraphEquiv\_closedGraph}}, \bpsource{GaussianTesting/Bosons/SchrodingerFockPair.lean}{305}{\texttt{fockPairSchrodingerEquiv\_P0B}}, \bpsource{GaussianTesting/Bosons/NumberTensor.lean}{219}{\texttt{fockPairSchrodingerEquiv\_tpPair}}.}

\textbf{Blueprint.} \bplink{prelim:bosons}{Bosonic Gaussian preliminaries}, \bplink{def:P0B}{Bosonic kernel-test projection}.

\subsection{Weyl operators in the Schrödinger representation}\label{paper:weyl-action}
\textbf{Checked.} The explicit translation and phase multiplier define unitary Weyl displacements on L2(R\textasciicircum{}m), with the Weyl composition law, adjoints, and strong continuity. \bpsource{material/manuscript.tex}{921}{Manuscript, line 921.}

weylOperator uses the manuscript's explicit wavefunction action as its definition on actual Lebesgue L2, for every number of modes. Its almost-everywhere action, unitarity, zero displacement, composition phase exp(-i sigma(r,s)/2), opposite-parameter inverse and adjoint, and joint parameter/wavefunction continuity are checked. The manuscript retains its equivalent unbounded-exponential notation as background; that identification is not used or proved here. This completed component is documented inside the combined bosonic preliminaries node, which is checked for the background required by the Section 6 proofs. In the manuscript coordinates r=(x,p), the wavefunction variable is z; the local translation parameter q in Lean denotes this x.

\bpdecls{\textbf{Lean.} \bpsource{GaussianTesting/Bosons/Weyl.lean}{162}{\texttt{weylOperator}}, \bpsource{GaussianTesting/Bosons/Weyl.lean}{167}{\texttt{weylOperator\_coeFn}}, \bpsource{GaussianTesting/Bosons/Weyl.lean}{196}{\texttt{weylOperator\_zero}}, \bpsource{GaussianTesting/Bosons/Weyl.lean}{214}{\texttt{weylOperator\_mul}}, \bpsource{GaussianTesting/Bosons/Weyl.lean}{253}{\texttt{weylOperator\_symm\_neg}}, \bpsource{GaussianTesting/Bosons/Weyl.lean}{260}{\texttt{weylOperator\_adjoint}}, \bpsource{GaussianTesting/Bosons/Weyl.lean}{267}{\texttt{weylOperator\_norm}}, \bpsource{GaussianTesting/Bosons/Weyl.lean}{270}{\texttt{weylOperator\_inner}}, \bpsource{GaussianTesting/Bosons/WeylContinuity.lean}{128}{\texttt{continuous\_weylOperator}}, \bpsource{GaussianTesting/Bosons/WeylContinuity.lean}{144}{\texttt{continuous\_weylOperator\_apply}}.}

\textbf{Blueprint.} \bplink{prelim:bosons}{Bosonic Gaussian preliminaries}.

\subsection{Characteristic functions of bosonic states}\label{paper:weyl}
\textbf{Partial coverage.} The characteristic function is Tr(rho D\_r), equals one at zero, is continuous, has modulus at most one, conjugates under parameter negation, and determines the state. \bpsource{material/manuscript.tex}{922}{Manuscript, line 922.}

The characteristic function uses the explicit Weyl action and product order Tr(rho D\_r). Basis independence, Gram expansion, conjugation, and continuity hold for positive trace-class operators; normalization and the modulus bound hold for density operators. Pure-state expectations equal rank-one traces. Displacements give symplectic linear phases, tensor products multiply, and the vacuum and displaced-vacuum formulas are checked. Pure-state injectivity is proved by L1 Fourier uniqueness, correlation kernels, and Fubini, without regularity or moment assumptions. The pure Gaussian characteristic-function reconstruction and vacuum-orbit characterization required by Section 6 are checked. Injectivity for arbitrary mixed states is unused by these proofs and remains unformalized.

\bpdecls{\textbf{Lean.} \bpsource{GaussianTesting/Bosons/CharacteristicFunction.lean}{95}{\texttt{waveCharacteristic}}, \bpsource{GaussianTesting/Bosons/CharacteristicFunction.lean}{102}{\texttt{waveCharacteristic\_gramOp}}, \bpsource{GaussianTesting/Bosons/CharacteristicFunction.lean}{109}{\texttt{waveCharacteristic\_basis}}, \bpsource{GaussianTesting/Bosons/CharacteristicFunction.lean}{116}{\texttt{waveCharacteristic\_zero}}, \bpsource{GaussianTesting/Bosons/CharacteristicFunction.lean}{129}{\texttt{norm\_waveCharacteristic\_le\_one}}, \bpsource{GaussianTesting/Bosons/CharacteristicFunction.lean}{136}{\texttt{waveCharacteristic\_neg}}, \bpsource{GaussianTesting/Bosons/CharacteristicFunction.lean}{147}{\texttt{continuous\_waveCharacteristic}}, \bpsource{GaussianTesting/Bosons/PureCharacteristic.lean}{31}{\texttt{waveCharacteristic\_rankOne}}, \bpsource{GaussianTesting/Bosons/PureCharacteristic.lean}{85}{\texttt{pureWaveCharacteristic\_displacement}}, \bpsource{GaussianTesting/Bosons/WeylRotation.lean}{282}{\texttt{pureTensorCharacteristic\_tmul}}, \bpsource{GaussianTesting/Bosons/WeylRotation.lean}{301}{\texttt{pureTensorCharacteristic\_balanced\_product}}, \bpsource{GaussianTesting/Bosons/CharacteristicInjectivity.lean}{103}{\texttt{rankOne\_eq\_of\_pureWaveCharacteristic\_eq}}.}

\textbf{Blueprint.} \bplink{prelim:bosons}{Bosonic Gaussian preliminaries}.

\subsection{Symplectic Gaussian unitaries}\label{paper:boson-symplectic}
\textbf{Checked.} Every real symplectic matrix has a unitary implementation on actual L2 with the manuscript inverse Weyl covariance; identical copies commute with copy rotations. \bpsource{material/manuscript.tex}{928}{Manuscript, line 928.}

Quadratic phase multiplication, Jacobian-normalized invertible position changes, and the unitary Fourier transform implement symmetric shears and diagonal blocks. A proved symplectic block factorization constructs an implementation for every S. Lean proves U\_S adjoint D\_r U\_S = D\_(S inverse r), characteristic-function covariance, and exact commutation of U\_S tensor U\_S with every copy rotation. These are the action properties required by Section 6. Projective/metaplectic structure and quadrature/exponential interpretations are not prerequisites and are not asserted by these declarations. The newly explicit general U(m) second-quantization description is also unused background; the real orthogonal copy actions needed here are independently constructed.

\bpdecls{\textbf{Lean.} \bpsource{GaussianTesting/Bosons/SymplecticShear.lean}{82}{\texttt{waveQuadraticPhase\_weyl}}, \bpsource{GaussianTesting/Bosons/LinearChange.lean}{207}{\texttt{waveLinearChange\_weyl}}, \bpsource{GaussianTesting/Bosons/FourierAction.lean}{197}{\texttt{waveFourier\_conjugate}}, \bpsource{GaussianTesting/Bosons/FourierTensor.lean}{247}{\texttt{schrodingerTensorRotation\_fourier}}, \bpsource{GaussianTesting/Bosons/SymplecticDecomposition.lean}{161}{\texttt{symplectic\_decomposition}}, \bpsource{GaussianTesting/Bosons/SymplecticAction.lean}{101}{\texttt{symplectic\_hasWaveImplementation}}, \bpsource{GaussianTesting/Bosons/SymplecticAction.lean}{111}{\texttt{waveSymplecticUnitary}}, \bpsource{GaussianTesting/Bosons/SymplecticAction.lean}{125}{\texttt{waveSymplecticUnitary\_conjugate\_weyl}}, \bpsource{GaussianTesting/Bosons/SymplecticAction.lean}{146}{\texttt{pureWaveCharacteristic\_symplectic}}, \bpsource{GaussianTesting/Bosons/SymplecticAction.lean}{139}{\texttt{waveSymplecticUnitary\_commutesWithCopyRotations}}.}

\textbf{Blueprint.} \bplink{prelim:bosons}{Bosonic Gaussian preliminaries}.

\subsection{Physical bosonic Gaussian classes}\label{paper:boson-gaussian}
\textbf{Partial coverage.} The paper describes pure Gaussians through quadratic ground states, Gaussian wave functions, displaced symplectic vacuum orbits, or Gaussian characteristic functions. Centered Gaussians and coherents are specified separately. \bpsource{material/manuscript.tex}{939}{Manuscript, line 939.}

The physical state sets use the manuscript independent vacuum-orbit descriptions: meanZeroGaussianWaveStates, gaussianWaveStates, and coherentWaveStates. Gaussian characteristic functions of actual unit wavefunctions are proved to give the centered/displaced symplectic vacuum orbits. Nondegeneracy follows from the Weyl action; Fourier inversion identifies correlation kernels, a continuous Gaussian representative, and its symplectic vacuum realization. No moment or nonsingular-covariance hypothesis is added. These are the characterizations required by Section 6. Equivalence to quadratic-Hamiltonian ground states and explicit unbounded mean-operator descriptions is unused and not asserted.

\bpdecls{\textbf{Lean.} \bpsource{GaussianTesting/Bosons/GaussianWaveStates.lean}{31}{\texttt{meanZeroGaussianWaveStates}}, \bpsource{GaussianTesting/Bosons/GaussianWaveStates.lean}{37}{\texttt{gaussianWaveStates}}, \bpsource{GaussianTesting/Bosons/GaussianWaveStates.lean}{43}{\texttt{coherentWaveStates}}, \bpsource{GaussianTesting/Bosons/GaussianCharacteristic.lean}{100}{\texttt{norm\_pureWaveCharacteristic\_lt\_one}}, \bpsource{GaussianTesting/Bosons/GaussianCharacteristic.lean}{131}{\texttt{gaussianCharacteristic\_neg\_posDef}}, \bpsource{GaussianTesting/Bosons/WaveGaussianOrbit.lean}{158}{\texttt{exists\_symplectic\_vacuum\_of\_quadratic\_wavefunction}}, \bpsource{GaussianTesting/Bosons/GaussianCharacteristicReconstruction.lean}{283}{\texttt{exists\_symplectic\_vacuum\_of\_pureCharacteristic\_quadratic}}, \bpsource{GaussianTesting/Bosons/GaussianCharacteristicReconstruction.lean}{303}{\texttt{exists\_displaced\_symplectic\_vacuum\_of\_gaussian\_characteristic}}.}

\textbf{Blueprint.} \bplink{prelim:bosons}{Bosonic Gaussian preliminaries}.

\subsection{Bosonic kernel projector and graph domain}\label{paper:lambda-b}
\textbf{Checked.} Lambda is defined on the graph-norm closure of the finite Fock number-vector span. Its closed kernel has a bounded orthogonal projection P0 defined on every input. \bpsource{material/manuscript.tex}{955}{Manuscript, line 955.}

Both generators are the graph closure of the specified finite number-vector expression. The actual occupation-to-wavefunction unitary maps the finite graph and its closure exactly; hence it maps their domains and kernels, and intertwines P0B with schrodingerP0. The initial domain is a dense core. The physical closed kernel equals the subspace fixed by all copy rotations.

\bpdecls{\textbf{Lean.} \bpsource{GaussianTesting/Bosons/ClosedLambda.lean}{128}{\texttt{closedLambdaB\_graph}}, \bpsource{GaussianTesting/Bosons/ClosedLambda.lean}{162}{\texttt{closedLambdaB\_eq\_closure}}, \bpsource{GaussianTesting/Bosons/ClosedLambda.lean}{169}{\texttt{closedLambdaB\_hasCore}}, \bpsource{GaussianTesting/Bosons/ClosedLambda.lean}{267}{\texttt{P0B\_eq\_closedLambdaB\_kernelProjection}}, \bpsource{GaussianTesting/Bosons/SchrodingerLambda.lean}{64}{\texttt{schrodingerFiniteLambda\_mem\_domain}}, \bpsource{GaussianTesting/Bosons/SchrodingerLambda.lean}{37}{\texttt{schrodingerLambda\_graph}}, \bpsource{GaussianTesting/Bosons/SchrodingerLambda.lean}{81}{\texttt{schrodingerLambda\_eq\_closure}}, \bpsource{GaussianTesting/Bosons/SchrodingerLambda.lean}{88}{\texttt{schrodingerLambda\_hasCore}}, \bpsource{GaussianTesting/Bosons/SchrodingerLambda.lean}{123}{\texttt{schrodingerP0}}, \bpsource{GaussianTesting/Bosons/SchrodingerLambda.lean}{136}{\texttt{schrodingerRotation\_P0}}, \bpsource{GaussianTesting/Bosons/SchrodingerFockPair.lean}{227}{\texttt{fockPairGraphEquiv\_finiteGraph}}, \bpsource{GaussianTesting/Bosons/SchrodingerFockPair.lean}{256}{\texttt{fockPairGraphEquiv\_closedGraph}}, \bpsource{GaussianTesting/Bosons/SchrodingerFockPair.lean}{268}{\texttt{fockPairSchrodingerEquiv\_mem\_graph\_iff}}, \bpsource{GaussianTesting/Bosons/SchrodingerFockPair.lean}{283}{\texttt{fockPairSchrodingerEquiv\_mem\_kernel\_iff}}, \bpsource{GaussianTesting/Bosons/SchrodingerFockPair.lean}{305}{\texttt{fockPairSchrodingerEquiv\_P0B}}, \bpsource{GaussianTesting/Bosons/NumberProjection.lean}{56}{\texttt{waveNumberCoordinates\_conj\_P0}}, \bpsource{GaussianTesting/Bosons/PhysicalMeanZero.lean}{89}{\texttt{schrodingerMeanZeroAcceptance}}.}

\textbf{Blueprint.} \bplink{def:P0B}{Bosonic kernel-test projection}, \bplink{lem:purity}{Purity bound for a density operator}.

\subsection{Rotation of two copies (Eq. 6.1)}\label{paper:unbounded-rotation}
\textbf{Checked.} The strongly continuous copy-rotation group has the manuscript wavefunction action and finite-core derivative -i Lambda; its fixed subspace equals the prescribed closed generator kernel. Orthogonal copy matrices act by inverse pullback. \bpsource{material/manuscript.tex}{957}{Manuscript, line 957.}

The actual L2 pullback satisfies the manuscript action, group law and strong continuity. On the prescribed finite core its derivative is -i Lambda, and the integrated identity extends to the graph closure. Every kernel vector is rotation-fixed; conversely, periodic averaging preserves the finite-dimensional polynomial orbit span, has zero generator, and extends by density and graph closure. Thus the fixed space equals the closed kernel. The quarter-turn is the signed swap. The O(t)/SO(t) copy representations and their two-copy restriction are checked. Exponential notation is unused background, not an outstanding prerequisite for the Section 6 proofs. Copy coordinates are ordered as Fin t × Fin m, so the checked inverse pullback is exactly the manuscript R inverse tensor I\_m action.

\bpdecls{\textbf{Lean.} \bpsource{GaussianTesting/Bosons/SchrodingerRotation.lean}{161}{\texttt{schrodingerRotation\_coeFn}}, \bpsource{GaussianTesting/Bosons/SchrodingerRotation.lean}{177}{\texttt{schrodingerRotation\_add}}, \bpsource{GaussianTesting/Bosons/SchrodingerRotation.lean}{193}{\texttt{continuous\_schrodingerRotation}}, \bpsource{GaussianTesting/Bosons/SchrodingerCore.lean}{119}{\texttt{hasDerivAt\_schrodingerRotation\_polynomialGaussian}}, \bpsource{GaussianTesting/Bosons/RotationGraphClosure.lean}{68}{\texttt{RotationGraphClosure.integrated\_identity}}, \bpsource{GaussianTesting/Bosons/SchrodingerLambda.lean}{94}{\texttt{schrodingerRotation\_fixed\_of\_mem\_kernel}}, \bpsource{GaussianTesting/Bosons/SchrodingerTensor.lean}{244}{\texttt{schrodingerRotation\_quarter\_waveTensor}}, \bpsource{GaussianTesting/Bosons/SchrodingerCopies.lean}{182}{\texttt{schrodingerCopiesRepresentation}}, \bpsource{GaussianTesting/Bosons/SchrodingerCopies.lean}{160}{\texttt{schrodingerCopiesAction\_coeFn}}, \bpsource{GaussianTesting/Bosons/SchrodingerCopies.lean}{172}{\texttt{schrodingerCopiesAction\_mul}}, \bpsource{GaussianTesting/Bosons/SchrodingerCopies.lean}{249}{\texttt{schrodingerSpecialOrthogonalRepresentation}}, \bpsource{GaussianTesting/Bosons/SchrodingerCopies.lean}{262}{\texttt{continuous\_schrodingerSpecialOrthogonalRepresentation\_apply}}, \bpsource{GaussianTesting/Bosons/SchrodingerCopies.lean}{290}{\texttt{schrodingerSpecialOrthogonalRepresentation\_planeRotation}}, \bpsource{GaussianTesting/Bosons/RotationFixedSpace.lean}{160}{\texttt{schrodingerRotationAverage\_mem\_kernel}}, \bpsource{GaussianTesting/Bosons/RotationFixedSpace.lean}{176}{\texttt{schrodingerLambda\_mem\_kernel\_iff\_rotation\_fixed}}.}

\textbf{Blueprint.} \bplink{prelim:bosons}{Bosonic Gaussian preliminaries}.

\subsection{Gaussian conjugation of the bosonic generator (Eq. 6.2)}\label{paper:boson-invariance}
\textbf{Partial coverage.} The doubled symplectic Gaussian unitary preserves Lambda. \bpsource{material/manuscript.tex}{973}{Manuscript, line 973.}

For every actual symplectic Gaussian unitary, identical copies commute with every physical copy rotation. The fixed-space/kernel equivalence then proves preservation of the prescribed closed kernel and commutation with its orthogonal projection P0, without a supplied domain hypothesis. This is the invariance required by the manuscript propositions. Full equality of the unbounded operators on their domains is not asserted here and is not needed for these kernel-test proofs. The explicitly stated consequence p0(U\_S psi)=p0(psi) is exported for every wavefunction, without a normalization or moment hypothesis.

\bpdecls{\textbf{Lean.} \bpsource{GaussianTesting/Bosons/GaussianKernel.lean}{64}{\texttt{CommutesWithCopyRotations.mem\_kernel}}, \bpsource{GaussianTesting/Bosons/GaussianKernel.lean}{93}{\texttt{CommutesWithCopyRotations.tensorP0}}, \bpsource{GaussianTesting/Bosons/GaussianKernel.lean}{104}{\texttt{waveSymplecticUnitary\_tensorP0}}, \bpsource{GaussianTesting/Bosons/PhysicalMeanZero.lean}{94}{\texttt{CommutesWithCopyRotations.meanZeroAcceptance}}, \bpsource{GaussianTesting/Bosons/PhysicalMeanZero.lean}{114}{\texttt{schrodingerMeanZeroAcceptance\_symplectic}}.}

\textbf{Blueprint.} \bplink{def:P0B}{Bosonic kernel-test projection}, \bplink{prp:mean-zero-gaussian}{Mean-zero states accepted with certainty}.

\subsection{Proposition 6.1}\label{paper:r6-1}
\textbf{Checked.} For pure inputs, physical mean-zero Gaussianity is equivalent to Lambda(psi tensor psi)=0 on its domain and to p0=1. \bpsource{material/manuscript.tex}{980}{Manuscript, line 980.}

bosonicMeanZeroGaussian\_characterization proves all three manuscript items equivalent for every unit wavefunction on actual L2(R\textasciicircum{}m): membership in the independently defined symplectic-vacuum orbit, membership of its tensor square in the prescribed graph-closed Lambda kernel (including its domain), and physical acceptance p0=1. No moment, covariance nondegeneracy, basis completeness, or operator-intertwining hypothesis is supplied. The Gaussian set uses the manuscript vacuum-orbit characterization; unused ground-state and explicit mean-operator descriptions are not assumed equivalences. The forward implication uses the vacuum finite kernel and symplectic invariance. The converse follows the manuscript balanced characteristic equation, positivity, logarithm and polarization, then the checked background Gaussian reconstruction.

\textbf{Proof route.} The physical proof follows the manuscript. The constant vacuum polynomial has zero finite generator, and symplectic kernel invariance proves the forward orbit implication. In the converse, kernel membership gives rotation invariance; Weyl covariance yields the balanced characteristic equation; continuity and iteration give positivity; log chi obeys the parallelogram law and polarization yields a real quadratic exponent. The proved Gaussian-characteristic reconstruction identifies the independent symplectic vacuum orbit, deriving nondegeneracy internally. Projection fixation is equivalent to acceptance one and includes the prescribed closed-operator domain.

\bpdecls{\textbf{Lean.} \bpsource{GaussianTesting/Bosons/PhysicalMeanZero.lean}{178}{\texttt{bosonicMeanZeroGaussian\_characterization}}, \bpsource{GaussianTesting/Bosons/PhysicalMeanZero.lean}{168}{\texttt{schrodingerMeanZeroAcceptance\_eq\_one\_iff}}, \bpsource{GaussianTesting/Bosons/PhysicalMeanZero.lean}{159}{\texttt{mem\_meanZeroGaussianWaveStates\_iff\_kernel}}, \bpsource{GaussianTesting/Bosons/PhysicalMeanZero.lean}{137}{\texttt{schrodingerMeanZeroAcceptance\_eq\_one\_iff\_kernel}}, \bpsource{GaussianTesting/Bosons/KernelCore.lean}{57}{\texttt{P0B\_fixed\_iff\_blocks}}, \bpsource{GaussianTesting/Bosons/MeanZeroCharacteristic.lean}{45}{\texttt{pureWaveCharacteristic\_balanced\_of\_kernel}}, \bpsource{GaussianTesting/Bosons/CharacteristicEquation.lean}{80}{\texttt{characteristic\_rotation\_positive}}, \bpsource{GaussianTesting/Bosons/CharacteristicEquation.lean}{116}{\texttt{exists\_bilinear\_of\_continuous\_parallelogram}}, \bpsource{GaussianTesting/Bosons/MeanZeroCharacteristic.lean}{62}{\texttt{pureWaveCharacteristic\_kernel\_eq\_exp\_quadratic}}, \bpsource{GaussianTesting/Bosons/PhysicalMeanZero.lean}{40}{\texttt{waveTensor\_vacuum\_mem\_kernel}}, \bpsource{GaussianTesting/Bosons/PhysicalMeanZero.lean}{75}{\texttt{schrodingerTensorP0\_symplectic\_vacuum}}.}

\textbf{Blueprint.} \bplink{prp:mean-zero-gaussian}{Mean-zero states accepted with certainty}, \bplink{thm:main-boson-mean-zero}{Fidelity bound for mean-zero Gaussian states}.

\subsection{Characteristic-function rotation identity (Eq. 6.3)}\label{paper:chi-rotation}
\textbf{Checked.} Two-copy rotation invariance implies a functional equation for chi whose continuous positive solutions have quadratic log chi. \bpsource{material/manuscript.tex}{995}{Manuscript, line 995.}

The equation is derived from the actual closed-kernel projection and Weyl covariance on the completed tensor square. For normalized continuous characteristic functions with conjugation symmetry and the usual bound, Lean proves strict positivity, a quadratic logarithm, and a nonpositive real quadratic exponent. This checked intermediate argument belongs to the completed Proposition 6.1 node, including its independent Gaussian-class identification and forward implication.

\bpdecls{\textbf{Lean.} \bpsource{GaussianTesting/Bosons/MeanZeroCharacteristic.lean}{45}{\texttt{pureWaveCharacteristic\_balanced\_of\_kernel}}, \bpsource{GaussianTesting/Bosons/CharacteristicEquation.lean}{56}{\texttt{characteristic\_rotation\_ne\_zero}}, \bpsource{GaussianTesting/Bosons/CharacteristicEquation.lean}{80}{\texttt{characteristic\_rotation\_positive}}, \bpsource{GaussianTesting/Bosons/CharacteristicEquation.lean}{160}{\texttt{characteristic\_rotation\_log\_parallelogram}}, \bpsource{GaussianTesting/Bosons/CharacteristicEquation.lean}{116}{\texttt{exists\_bilinear\_of\_continuous\_parallelogram}}, \bpsource{GaussianTesting/Bosons/CharacteristicEquation.lean}{210}{\texttt{characteristic\_rotation\_eq\_exp\_nonpos\_quadratic}}.}

\textbf{Blueprint.} \bplink{prp:mean-zero-gaussian}{Mean-zero states accepted with certainty}.

\subsection{Lemma 6.2}\label{paper:r6-2}
\textbf{Checked.} For every density operator rho, two-copy acceptance at least 1-tau implies that its operator norm is at least 1-2*tau. \bpsource{material/manuscript.tex}{1009}{Manuscript, line 1009.}

purity\_schrodinger proves the full statement for every positive trace-class density operator on actual L2(R\textasciicircum{}m), with no supplied Gram family, moment assumption, or intertwining hypothesis. The two-copy state is the canonical product-wavefunction realization of rho tensor rho, and P0 is the orthogonal projection onto the kernel of the specified finite Fock generator's graph closure. The complex operator trace is basis independent. The acceptance trace is proved real with real part in [0,1].

\textbf{Proof route.} The proof follows the manuscript. The finite-core differential identity passes to the graph closure, so U(pi/2) fixes the range of P0. Hence twice acceptance minus one is at most the real signed-swap trace. The actual product-wavefunction formula identifies that complex trace with Tr(J rho\textasciicircum{}2). Cauchy-Schwarz and unitarity give |Tr(J rho\textasciicircum{}2)| <= Tr(rho\textasciicircum{}2) <= norm(rho). Positive-square-root Gram representations and absolute double-sum interchange justify every trace identity for arbitrary density operators.

\bpdecls{\textbf{Lean.} \bpsource{GaussianTesting/Bosons/SchrodingerPurity.lean}{96}{\texttt{purity\_schrodinger}}, \bpsource{GaussianTesting/Bosons/SchrodingerPurity.lean}{80}{\texttt{schrodinger\_signedSwap\_trace}}, \bpsource{GaussianTesting/Bosons/SchrodingerPurity.lean}{43}{\texttt{schrodingerDensityTensor\_product}}, \bpsource{GaussianTesting/Bosons/SchrodingerLambda.lean}{136}{\texttt{schrodingerRotation\_P0}}, \bpsource{GaussianTesting/Bosons/HilbertDensity.lean}{204}{\texttt{HilbertDensity.IsDensityOperator.exists\_normalized\_gram\_representation}}, \bpsource{GaussianTesting/Bosons/SignedSwapPurity.lean}{196}{\texttt{HilbertDensity.signedSwap\_trace\_norm\_le\_purity}}, \bpsource{GaussianTesting/Bosons/SignedSwapPurity.lean}{143}{\texttt{HilbertDensity.IsDensityOperator.purity\_le\_norm}}, \bpsource{GaussianTesting/Bosons/SchrodingerPurity.lean}{58}{\texttt{schrodinger\_acceptance\_real}}, \bpsource{GaussianTesting/Bosons/SchrodingerPurity.lean}{69}{\texttt{schrodinger\_acceptance\_mem\_Icc}}.}

\textbf{Blueprint.} \bplink{prelim:bosons}{Bosonic Gaussian preliminaries}, \bplink{lem:purity}{Purity bound for a density operator}.

\subsection{Beam splitter rotates displacements (Eq. 6.4)}\label{paper:displacement-rotation}
\textbf{Checked.} The positive balanced splitter rotates a pair of Weyl displacements to difference and sum parameters. \bpsource{material/manuscript.tex}{1038}{Manuscript, line 1038.}

Coordinate pullback proves covariance for every angle on actual two-copy L2. The explicit two-copy Weyl action is identified with the completed tensor product of the one-copy Weyl operators. The balanced-angle conjugation theorem is exactly the manuscript formula on every completed-tensor vector, with difference/sum parameters and no assumed representation bridge.

\bpdecls{\textbf{Lean.} \bpsource{GaussianTesting/Bosons/WeylRotation.lean}{140}{\texttt{schrodingerWeyl\_waveTensor}}, \bpsource{GaussianTesting/Bosons/WeylRotation.lean}{210}{\texttt{schrodingerRotation\_weyl}}, \bpsource{GaussianTesting/Bosons/WeylRotation.lean}{268}{\texttt{schrodingerTensorRotation\_balanced\_weyl}}.}

\textbf{Blueprint.} \bplink{prelim:bosons}{Bosonic Gaussian preliminaries}, \bplink{prp:general-certainty}{Gaussian and coherent certainty}.

\subsection{Four-copy accepting projection}\label{paper:four-copy}
\textbf{Checked.} Undo the balanced splitter on each identical pair, then apply P0 on registers two and four. \bpsource{material/manuscript.tex}{1043}{Manuscript, line 1043.}

The accepting projection is constructed on four completed copies of actual one-copy wavefunctions, using the physical inverse balanced rotation and the graph-closed kernel projection on registers 2 and 4. Symmetry and idempotence are proved. Its Born expectation is the physical reduced-state trace. The physical proofs use their direct construction and partial-trace identity.

\bpdecls{\textbf{Lean.} \bpsource{GaussianTesting/Bosons/SchrodingerFourCopy.lean}{70}{\texttt{schrodingerFourCopyProjection}}, \bpsource{GaussianTesting/Bosons/SchrodingerFourCopy.lean}{74}{\texttt{schrodingerFourCopyProjection\_symmetric}}, \bpsource{GaussianTesting/Bosons/SchrodingerFourCopy.lean}{81}{\texttt{schrodingerFourCopyProjection\_idempotent}}, \bpsource{GaussianTesting/Bosons/SchrodingerFourCopy.lean}{90}{\texttt{schrodingerFourCopyAcceptance}}, \bpsource{GaussianTesting/Bosons/SchrodingerFourCopy.lean}{155}{\texttt{schrodingerFourCopyProjection\_expectation\_eq\_reducedTrace}}, \bpsource{GaussianTesting/Bosons/SchrodingerFourCopy.lean}{166}{\texttt{schrodingerFourCopyAcceptance\_eq\_reducedTrace}}.}

\textbf{Blueprint.} \bplink{lem:four-copy}{Four-copy projection and reduced-state expectation}, \bplink{target:fig1-general}{Figure 1(d), four-copy measurement}.

\subsection{Reduced mixed-state expectation (Eq. 6.5)}\label{paper:boson-reduced-state}
\textbf{Checked.} The four-copy acceptance probability equals the two-copy kernel test on two copies of the reduced state after undoing the splitter. \bpsource{material/manuscript.tex}{1047}{Manuscript, line 1047.}

The actual inverse-rotation output has a basis-independent reduced operator, with the defining partial-trace identity for every bounded observable. Unit inputs give positive trace-class reduced states of trace one. Regrouping two copies yields the tensor square of that reduced operator, so the physical four-copy Born probability equals Tr(P0 rho tensor rho) and lies in [0,1].

\bpdecls{\textbf{Lean.} \bpsource{GaussianTesting/Bosons/PartialTrace.lean}{137}{\texttt{HilbertDensity.operatorTrace\_mul\_reducedRight}}, \bpsource{GaussianTesting/Bosons/PartialTrace.lean}{154}{\texttt{HilbertDensity.reducedRight\_basis}}, \bpsource{GaussianTesting/Bosons/PartialTraceTensor.lean}{56}{\texttt{HilbertDensity.reducedRight\_reshuffle\_square}}, \bpsource{GaussianTesting/Bosons/SchrodingerFourCopy.lean}{106}{\texttt{schrodingerReducedState\_isDensityOperator}}, \bpsource{GaussianTesting/Bosons/SchrodingerFourCopy.lean}{113}{\texttt{schrodingerReducedState\_basis}}, \bpsource{GaussianTesting/Bosons/SchrodingerFourCopy.lean}{155}{\texttt{schrodingerFourCopyProjection\_expectation\_eq\_reducedTrace}}, \bpsource{GaussianTesting/Bosons/SchrodingerFourCopy.lean}{166}{\texttt{schrodingerFourCopyAcceptance\_eq\_reducedTrace}}, \bpsource{GaussianTesting/Bosons/SchrodingerFourCopy.lean}{175}{\texttt{schrodingerFourCopyAcceptance\_mem\_Icc}}.}

\textbf{Blueprint.} \bplink{lem:purity}{Purity bound for a density operator}, \bplink{lem:four-copy}{Four-copy projection and reduced-state expectation}.

\subsection{Coherent test and acceptance probability (Eq. 6.6)}\label{paper:coherent-q}
\textbf{Checked.} The coherent test conjugates the second-register vacuum projection by the balanced splitter. \bpsource{material/manuscript.tex}{1060}{Manuscript, line 1060.}

The physical operator Q is constructed from the actual balanced coordinate rotation and the rank-one projection onto the normalized Gaussian vacuum in L2(R\textasciicircum{}m). Symmetry, idempotence, swap commutation, reality, and acceptance in [0,1] are proved. Its expectation is the vacuum trace of the basis-independent reduced state. QB\_eq\_conj proves the occupation-coordinate identity and waveNumberCoordinates\_conj\_Q identifies the two projections. The particle-counting readout is checked in Figure 1(c).

\bpdecls{\textbf{Lean.} \bpsource{GaussianTesting/Bosons/SchrodingerCoherent.lean}{38}{\texttt{schrodingerQ}}, \bpsource{GaussianTesting/Bosons/SchrodingerVacuum.lean}{95}{\texttt{waveVacuum\_norm}}, \bpsource{GaussianTesting/Bosons/SchrodingerVacuum.lean}{53}{\texttt{waveVacuum\_annihilation}}, \bpsource{GaussianTesting/Bosons/SchrodingerVacuum.lean}{106}{\texttt{waveParity\_vacuum}}, \bpsource{GaussianTesting/Bosons/SchrodingerSwap.lean}{81}{\texttt{schrodingerTensorRotation\_balanced\_swap}}, \bpsource{GaussianTesting/Bosons/SchrodingerCoherent.lean}{110}{\texttt{schrodingerQ\_isAcceptance}}, \bpsource{GaussianTesting/Bosons/SchrodingerCoherent.lean}{123}{\texttt{schrodingerQ\_expectation\_real}}, \bpsource{GaussianTesting/Bosons/SchrodingerCoherent.lean}{129}{\texttt{schrodingerCoherentAcceptance\_mem\_Icc}}, \bpsource{GaussianTesting/Bosons/SchrodingerCoherent.lean}{136}{\texttt{schrodingerQ\_expectation\_eq\_reducedTrace}}, \bpsource{GaussianTesting/Bosons/BeamSplitter.lean}{851}{\texttt{QB\_eq\_conj}}, \bpsource{GaussianTesting/Bosons/SchrodingerCoherent.lean}{161}{\texttt{schrodingerCoherentAcceptance\_eq\_vacuumExpectation}}.}

\textbf{Blueprint.} \bplink{lem:S-spectrum-coherent-states}{Spectrum of the three-copy coherent test}, \bplink{lem:coherent-projection}{Coherent accepting projection}, \bplink{thm:main-boson}{Fidelity bound for general bosonic Gaussian states}, \bplink{target:fig1-coherent}{Figure 1(c), vacuum measurement}.

\subsection{Proposition 6.3}\label{paper:r6-3}
\textbf{Checked.} Gaussianity is equivalent to perfect four-copy acceptance, and coherence is equivalent to perfect coherent-test acceptance. \bpsource{material/manuscript.tex}{1067}{Manuscript, line 1067.}

bosonicGaussian\_and\_coherent\_characterization proves both manuscript equivalences for every unit wavefunction on actual L2(R\textasciicircum{}m): membership in the independently defined displaced symplectic-vacuum Gaussian orbit iff physical four-copy acceptance p=1, and membership in the displaced-vacuum coherent orbit iff q=1. The Gaussian converse uses the actual reduced-state trace, Lemma 6.2, norm-one purity, product factorization, Proposition 6.1, symplectic transport to the vacuum, and the physical coherent equivalence. The forward implication uses copy-rotation invariance of centered Gaussians and displacement covariance. Only unit norm is assumed; no moment, factorization, or covariance hypothesis is supplied.

\textbf{Proof route.} The physical proof follows the manuscript. For a Gaussian input, the inverse balanced splitter moves its identical displacements into register one, leaving a centered Gaussian in register two, so P0 accepts the two retained factors. Conversely p=1 gives Tr(P0 rho tensor rho)=1; Lemma 6.2 gives norm rho=1, hence rho is rank one. Pure-reduction factorization and Proposition 6.1 identify its unit right factor as centered Gaussian. Applying the inverse symplectic unitary on both copies makes this factor the vacuum; rotation commutation gives coherent certainty, and Weyl covariance transports the displaced vacuum back to a general Gaussian input. The coherent equivalence uses the manuscript normalized characteristic equation, continuous linear phase and injectivity.

\bpdecls{\textbf{Lean.} \bpsource{GaussianTesting/Bosons/PhysicalGaussian.lean}{41}{\texttt{bosonicGaussian\_and\_coherent\_characterization}}, \bpsource{GaussianTesting/Bosons/PhysicalGaussian.lean}{29}{\texttt{schrodingerFourCopyAcceptance\_eq\_one\_iff}}, \bpsource{GaussianTesting/Bosons/GaussianPurityReduction.lean}{64}{\texttt{schrodingerReducedState\_rankOne\_of\_acceptance\_one}}, \bpsource{GaussianTesting/Bosons/GaussianPurityReduction.lean}{83}{\texttt{gaussian\_factorization\_of\_fourCopyAcceptance\_one}}, \bpsource{GaussianTesting/Bosons/GaussianSymplecticReduction.lean}{62}{\texttt{mem\_gaussianWaveStates\_of\_meanZero\_factorization}}, \bpsource{GaussianTesting/Bosons/GaussianForward.lean}{50}{\texttt{displaced\_meanZeroGaussian\_factorization}}, \bpsource{GaussianTesting/Bosons/GaussianForward.lean}{82}{\texttt{schrodingerFourCopyAcceptance\_of\_mem\_gaussianWaveStates}}, \bpsource{GaussianTesting/Bosons/PhysicalCoherent.lean}{119}{\texttt{schrodingerCoherentAcceptance\_eq\_one\_iff}}, \bpsource{GaussianTesting/Bosons/PhysicalCoherent.lean}{63}{\texttt{coherent\_factorization\_of\_acceptance\_one}}, \bpsource{GaussianTesting/Bosons/CoherentCharacteristic.lean}{91}{\texttt{coherent\_factorization\_normalized\_character}}, \bpsource{GaussianTesting/Bosons/CoherentCharacteristic.lean}{140}{\texttt{coherent\_factorization\_eq\_displaced\_vacuum}}, \bpsource{GaussianTesting/Bosons/PhysicalCoherent.lean}{93}{\texttt{displaced\_vacuum\_factorization}}.}

\textbf{Blueprint.} \bplink{thm:main-boson-coherent}{Fidelity bound for coherent states}, \bplink{prp:general-certainty}{Gaussian and coherent certainty}.

\subsection{Vacuum and coherent characteristic functions}\label{paper:vacuum-characteristic}
\textbf{Checked.} The vacuum characteristic function is eta(x,p) = exp(-(norm(x)\textasciicircum{}2 + norm(p)\textasciicircum{}2)/4). Displaced vacua multiply it by a symplectic linear phase, and eta satisfies the balanced-rotation product identity. \bpsource{material/manuscript.tex}{1099}{Manuscript, line 1099.}

The formula follows directly from the explicit normalized Gaussian wavefunction, the Weyl action, midpoint translation, and the Gaussian Fourier integral. The vacuum characteristic function never vanishes. Its displaced form follows from checked Weyl conjugation. These are background identities used in Proposition 6.3; the normalized functional-equation argument itself belongs to that proposition. No number-basis completeness or independent Gaussian-class equivalence is assumed.

\bpdecls{\textbf{Lean.} \bpsource{GaussianTesting/Bosons/VacuumCharacteristic.lean}{75}{\texttt{inner\_waveVacuum\_weylOperator}}, \bpsource{GaussianTesting/Bosons/VacuumCharacteristic.lean}{129}{\texttt{waveCharacteristic\_vacuum}}, \bpsource{GaussianTesting/Bosons/VacuumCharacteristic.lean}{136}{\texttt{vacuumCharacteristic\_ne\_zero}}, \bpsource{GaussianTesting/Bosons/VacuumCharacteristic.lean}{151}{\texttt{pureWaveCharacteristic\_displaced\_vacuum}}, \bpsource{GaussianTesting/Bosons/VacuumCharacteristic.lean}{164}{\texttt{vacuumCharacteristic\_balanced}}.}

\textbf{Blueprint.} \bplink{prelim:bosons}{Bosonic Gaussian preliminaries}.

\subsection{Normalized characteristic functions for coherence (Eq. 6.7)}\label{paper:normalized-chi}
\textbf{Checked.} A vacuum output gives a functional equation for the characteristic function divided by the vacuum characteristic function. Its solutions identify a displacement of the vacuum. \bpsource{material/manuscript.tex}{1105}{Manuscript, line 1105.}

For the actual inverse-splitter factorization, the characteristic function divided by the nonzero vacuum characteristic function is a continuous unit-circle character. Lean proves it is a symplectic linear phase and uses pure-state injectivity to identify a displaced vacuum. Together with the physical projection/factorization equivalence this proves the full coherent half of Proposition 6.3. These proof steps are recorded in that proposition node.

\bpdecls{\textbf{Lean.} \bpsource{GaussianTesting/Bosons/CoherentCharacteristic.lean}{76}{\texttt{coherent\_factorization\_normalized\_equation}}, \bpsource{GaussianTesting/Bosons/CoherentCharacteristic.lean}{91}{\texttt{coherent\_factorization\_normalized\_character}}, \bpsource{GaussianTesting/Bosons/CoherentCharacteristicEquation.lean}{87}{\texttt{continuous\_character\_eq\_weyl\_phase}}, \bpsource{GaussianTesting/Bosons/CoherentCharacteristic.lean}{120}{\texttt{coherent\_factorization\_characteristic\_linear}}, \bpsource{GaussianTesting/Bosons/CoherentCharacteristic.lean}{140}{\texttt{coherent\_factorization\_eq\_displaced\_vacuum}}, \bpsource{GaussianTesting/Bosons/PhysicalCoherent.lean}{119}{\texttt{schrodingerCoherentAcceptance\_eq\_one\_iff}}.}

\textbf{Blueprint.} \bplink{prp:general-certainty}{Gaussian and coherent certainty}.

\subsection{Theorem 6.4}\label{paper:r6-4}
\textbf{Checked.} If p0 >= 1 - tau, there is a mean-zero Gaussian witness with squared overlap at least 1 - 8 tau. The squared supremal fidelity is at least this squared overlap. \bpsource{material/manuscript.tex}{1121}{Manuscript, line 1121.}

For every mode count and unit physical wavefunction, the actual mean-zero test's acceptance bound gives a centered Gaussian witness with squared overlap at least 1−8 tau and at most squared supremal fidelity with meanZeroGaussianWaveStates. The physical completed tensor cube carries the strongly continuous SO(3) action; its three axis-fixed projections are exactly the pair tests. The arbitrary-representation Kac spectral theorem and Theorem 3.4 complete the manuscript proof. No additional spectral or covariance hypothesis is supplied.

\textbf{Proof route.} The proof follows the manuscript on actual wavefunctions. The canonical L2 tensor-cube equivalence intertwines the three pair rotations with the SO(3) coordinate-axis rotations. Period averaging and the checked closed-generator kernel identity identify each pair projection with its axis-fixed projection. The average is exactly the Kac operator; Theorem 4.1 gives its spectrum in [0,7/12] union \{1\}. Theorem 3.4 gives constant eight, and Proposition 6.1 identifies its certainty witness with the independent centered Gaussian vacuum orbit. The earlier particle-number-block argument is not a dependency of this physical spectral proof.

\bpdecls{\textbf{Lean.} \bpsource{GaussianTesting/Bosons/PhysicalMeanZeroRobustness.lean}{72}{\texttt{schrodinger\_meanZero\_fidelity}}, \bpsource{GaussianTesting/Bosons/PhysicalMeanZeroRobustness.lean}{44}{\texttt{schrodinger\_meanZero\_fidelity\_witness}}, \bpsource{GaussianTesting/Bosons/PhysicalMeanZeroRobustness.lean}{32}{\texttt{meanZeroGaussianWaveStates\_eq\_certaintyStates}}, \bpsource{GaussianTesting/Bosons/SchrodingerCube.lean}{125}{\texttt{waveCubeEquiv}}, \bpsource{GaussianTesting/Bosons/SchrodingerKac.lean}{57}{\texttt{schrodingerCubeRepresentation}}, \bpsource{GaussianTesting/Bosons/SchrodingerKac.lean}{131}{\texttt{schrodingerCube\_fixedProjection\_X}}, \bpsource{GaussianTesting/Bosons/SchrodingerKac.lean}{247}{\texttt{schrodingerCube\_fixedProjection\_Y}}, \bpsource{GaussianTesting/Bosons/SchrodingerKac.lean}{228}{\texttt{schrodingerCube\_fixedProjection\_Z}}, \bpsource{GaussianTesting/Bosons/SchrodingerKac.lean}{271}{\texttt{schrodinger\_threeCopyAverage\_eq\_kacAverage}}, \bpsource{GaussianTesting/Bosons/SchrodingerKac.lean}{286}{\texttt{schrodinger\_meanZero\_spectrum\_subset}}, \bpsource{GaussianTesting/Bosons/Main.lean}{1189}{\texttt{main\_boson\_mean\_zero}}, \bpsource{GaussianTesting/Bosons/Main.lean}{1096}{\texttt{lojasiewiczData\_P0B}}.}

\textbf{Theorem node.} \bplink{thm:main-boson-mean-zero}{Fidelity bound for mean-zero Gaussian states}.

\subsection{Theorem 6.5}\label{paper:r6-5}
\textbf{Checked.} If q >= 1 - tau, squared fidelity to the coherent class is at least 1 - 4 tau. \bpsource{material/manuscript.tex}{1144}{Manuscript, line 1144.}

For every mode count and unit physical wavefunction, schrodinger\_coherent\_fidelity proves that q(psi) >= 1 - tau implies squared supremal fidelity with the independent displaced-vacuum class coherentWaveStates is at least 1 - 4 tau. schrodinger\_coherent\_fidelity\_witness supplies a coherent comparison state with squared overlap between this lower bound and that fidelity. Only unit norm and the acceptance bound are supplied; tau >= 0 follows from q <= 1. The test, tensor cube, spectrum, and certainty class are the actual physical ones. No extra spectral, moment, or covariance hypothesis is assumed.

\textbf{Proof route.} The proof follows the manuscript directly. Lemma 6.6 gives the spectrum of the physical three-pair average with alpha=1/2. The physical Q is a symmetric idempotent commuting with swap, and the vacuum supplies a certainty state. Theorem 3.4 applies with Cgap(1/2)=4. Proposition 6.3 identifies its normalized certainty set with the independently defined displaced-vacuum coherent class. The same witness gives the squared-overlap bound and the squared-fidelity conclusion from Eq. (3.2). The older coefficient-model robustness theorem is separate supporting material and is not a dependency of this physical application.

\bpdecls{\textbf{Lean.} \bpsource{GaussianTesting/Bosons/PhysicalCoherentRobustness.lean}{78}{\texttt{schrodinger\_coherent\_fidelity}}, \bpsource{GaussianTesting/Bosons/PhysicalCoherentRobustness.lean}{48}{\texttt{schrodinger\_coherent\_fidelity\_witness}}, \bpsource{GaussianTesting/Bosons/PhysicalCoherentRobustness.lean}{35}{\texttt{coherentWaveStates\_eq\_certaintyStates}}, \bpsource{GaussianTesting/Bosons/Coherent.lean}{1008}{\texttt{main\_boson\_coherent}}, \bpsource{GaussianTesting/Bosons/Coherent.lean}{932}{\texttt{lojasiewiczData\_QB}}.}

\textbf{Theorem node.} \bplink{thm:main-boson-coherent}{Fidelity bound for coherent states}.

\subsection{Lemma 6.6}\label{paper:r6-6}
\textbf{Checked.} The spectrum of the average of the three pairwise coherent projections on the full three-copy space lies in [0,1/2] union \{1\}. \bpsource{material/manuscript.tex}{1151}{Manuscript, line 1151.}

The spectral inclusion is checked on the full completed tensor cube of actual L2(R\textasciicircum{}m), for every mode count. schrodingerCoherentAverage is the explicit average of the three physical pair projections built from Q = U(I tensor vacuum projection)U adjoint. The proved number basis, tensor identification, balanced rotation and vacuum projection identities conjugate it to the coefficient operator; its spectrum is therefore exactly the spectrum of Slp QB. The exact three-direction Gram spectrum is checked on every finite number block. For complement occupation k, its two branches are (1+2(-1/2)\textasciicircum{}k)/3 and (1-(-1/2)\textasciicircum{}k)/3. Their maxima are 1 at k=0, 1/2 at k=1,2, and strictly below 1/2 for k≥3. The public inclusion follows from these exact eigenvalues and the spectral-positivity argument. No basis completeness or operator-intertwining assumption is supplied.

\textbf{Proof route.} Lean uses vVec=(1,1,1) for the uniform direction V and wVec a for sqrt(6) times the manuscript normalized w\_ab in W=V-perp. The occupation frame orders the center first and complement second, exactly as the revised manuscript; k is the complement occupation. The manuscript separates the center-of-mass direction and its orthogonal complement. The normalized rotated occupation vectors eVecQ n a (p,q) realize these subspaces in coefficient coordinates. Distinct occupation labels are orthogonal, and within a label the three Gram overlaps are s=(-1/2)\textasciicircum{}|q|. Tq3 is exactly the sum of the associated rank-one projections divided by three. The common Gram matrix has characteristic polynomial (z-(1+2s))(z-(1-s))\textasciicircum{}2; nonzero spectral values of each finite number block are exactly the two corresponding branches divided by three for some occupation label. Complement occupation k=0 gives branches 1,0; k=1 gives 0,1/2; k=2 gives 1/2,1/4; the maximum is strictly less than 1/2 for k≥3. The three-projection polynomial inequality is now proved by classifying each self-adjoint eigenvector with this exact Gram calculation. Orthogonal number-block summation extends it to the infinite Fock space, and the spectral-positivity equivalence gives the inclusion. The proved number/Fock unitary, product-tensor identity, balanced splitters and vacuum projection transport the exact pair average to the physical completed tensor cube; unitary conjugation preserves the entire spectrum.

\bpdecls{\textbf{Lean.} \bpsource{GaussianTesting/Bosons/ThreeProjectionGap.lean}{40}{\texttt{three\_projection\_gap\_of\_gram}}, \bpsource{GaussianTesting/Bosons/CoherentBlocks.lean}{460}{\texttt{blockGapQ}}, \bpsource{GaussianTesting/Bosons/Coherent.lean}{910}{\texttt{coherent\_polynomial\_gap}}, \bpsource{GaussianTesting/Bosons/ThreeCopy.lean}{379}{\texttt{Slp\_nonneg}}, \bpsource{GaussianTesting/Bosons/ThreeCopy.lean}{349}{\texttt{Slp\_norm\_le\_one}}, \bpsource{GaussianTesting/Bosons/CoherentSpectrum.lean}{23}{\texttt{coherent\_spectrum\_subset}}, \bpsource{GaussianTesting/Bosons/NumberCoherent.lean}{109}{\texttt{schrodingerCoherentAverage}}, \bpsource{GaussianTesting/Bosons/NumberCoherent.lean}{102}{\texttt{waveNumberCoordinates\_conj\_Q}}, \bpsource{GaussianTesting/Bosons/NumberCoherent.lean}{131}{\texttt{schrodingerCoherentAverage\_spectrum}}, \bpsource{GaussianTesting/Bosons/NumberCoherent.lean}{144}{\texttt{schrodinger\_coherent\_spectrum\_subset}}, \bpsource{GaussianTesting/ThreeVectorGram.lean}{31}{\texttt{threeVectorGram\_charpoly\_eval}}, \bpsource{GaussianTesting/ThreeVectorGram.lean}{39}{\texttt{threeVectorGram\_spectrum}}, \bpsource{GaussianTesting/ThreeVectorGram.lean}{148}{\texttt{three\_family\_average\_mem\_spectrum\_iff}}, \bpsource{GaussianTesting/Bosons/OccupationGram.lean}{28}{\texttt{coherentSectorConstantEigenvalue}}, \bpsource{GaussianTesting/Bosons/OccupationGram.lean}{33}{\texttt{coherentSectorSumZeroEigenvalue}}, \bpsource{GaussianTesting/Bosons/OccupationGram.lean}{37}{\texttt{coherentSectorMaxEigenvalue}}, \bpsource{GaussianTesting/Bosons/OccupationGram.lean}{41}{\texttt{coherentSectorGram\_spectrum}}, \bpsource{GaussianTesting/Bosons/OccupationGram.lean}{55}{\texttt{coherentSectorEigenvalues\_zero}}, \bpsource{GaussianTesting/Bosons/OccupationGram.lean}{60}{\texttt{coherentSectorEigenvalues\_one}}, \bpsource{GaussianTesting/Bosons/OccupationGram.lean}{65}{\texttt{coherentSectorEigenvalues\_two}}, \bpsource{GaussianTesting/Bosons/OccupationGram.lean}{69}{\texttt{coherentSectorMaxEigenvalue\_zero}}, \bpsource{GaussianTesting/Bosons/OccupationGram.lean}{73}{\texttt{coherentSectorMaxEigenvalue\_one}}, \bpsource{GaussianTesting/Bosons/OccupationGram.lean}{77}{\texttt{coherentSectorMaxEigenvalue\_two}}, \bpsource{GaussianTesting/Bosons/OccupationGram.lean}{82}{\texttt{coherentSectorMaxEigenvalue\_lt\_half}}, \bpsource{GaussianTesting/Bosons/OccupationGram.lean}{96}{\texttt{coherentSectorMaxEigenvalue\_le\_half}}, \bpsource{GaussianTesting/Bosons/CoherentBlocks.lean}{440}{\texttt{Tq3\_eq\_family\_average}}, \bpsource{GaussianTesting/Bosons/CoherentBlocks.lean}{449}{\texttt{blockSpectrumQ\_nonzero}}.}

\textbf{Blueprint.} \bplink{lem:S-spectrum-coherent-states}{Spectrum of the three-copy coherent test}, \bplink{thm:main-boson-coherent}{Fidelity bound for coherent states}.

\subsection{Three-copy Fock decomposition}\label{paper:three-copy-fock}
\textbf{Partial coverage.} Split the copy index into the uniform center-of-mass direction and its orthogonal complement, then decompose by the complement's particle number. \bpsource{material/manuscript.tex}{1166}{Manuscript, line 1166.}

The rotated-mode occupation bases and Gram identities implement the particle-block estimates used in Lemma 6.6. The physical HilbertCube average is now exactly conjugate to that coefficient average using the proved number basis and test operators. The displayed abstract symmetric-Fock direct-sum isomorphism is not separately packaged; it supplies no unproved premise in the physical spectral theorem and is not an additional background target. Lean uses vVec=(1,1,1) for the uniform direction V and wVec a for sqrt(6) times the manuscript normalized w\_ab in W=V-perp. The occupation frame orders the center first and complement second, exactly as the revised manuscript; k is the complement occupation.

\bpdecls{\textbf{Lean.} \bpsource{GaussianTesting/Bosons/CoherentBlocks.lean}{460}{\texttt{blockGapQ}}, \bpsource{GaussianTesting/Bosons/ThreeProjectionGap.lean}{40}{\texttt{three\_projection\_gap\_of\_gram}}, \bpsource{GaussianTesting/Bosons/NumberCoherent.lean}{131}{\texttt{schrodingerCoherentAverage\_spectrum}}.}

\textbf{Blueprint.} \bplink{lem:S-spectrum-coherent-states}{Spectrum of the three-copy coherent test}.

\subsection{Theorem 6.7}\label{paper:r6-7}
\textbf{Checked.} If the four-copy test accepts with p >= 1 - tau, squared fidelity to pure Gaussians is at least 1 - 48 tau. \bpsource{material/manuscript.tex}{1202}{Manuscript, line 1202.}

For every mode count and unit physical wavefunction, four-copy acceptance at least 1−tau implies squared supremal fidelity at least 1−48 tau with the independently defined gaussianWaveStates. A physical Gaussian witness also satisfies the same squared-overlap bound. The top eigenstate, positive residual estimate, mean-zero robustness, symplectic transport and coherent robustness are all proved. Only unit norm and the acceptance bound are assumed; the small-error split is internal. The full reduced-operator covariance sigma=U\_S rho U\_S adjoint is exported, not just its diagonal matrix coefficients. Eqs. (6.8) and (6.9) are indexed separately as proof steps within this same theorem node.

\textbf{Proof route.} The physical proof follows the manuscript. For tau≤1/48, Lemma 6.2 gives a reduced-state top eigenvalue lambda≥1−2 tau. Positivity of the residual tensor terms gives p0(xi)≥1−tau/lambda². Theorem 6.4 supplies a centered Gaussian approximant; its inverse symplectic unitary transports the reduced expectation to vacuum. Commutation with the balanced splitter makes this coherent acceptance at least lambda−8 tau/lambda≥1−12 tau. Theorem 6.5 gives overlap at least 1−48 tau with a coherent state, and inverse transport gives a physical Gaussian witness. Larger errors are handled by nonnegativity.

\bpdecls{\textbf{Lean.} \bpsource{GaussianTesting/Bosons/PhysicalGeneralRobustness.lean}{194}{\texttt{schrodinger\_gaussian\_fidelity}}, \bpsource{GaussianTesting/Bosons/PhysicalGeneralRobustness.lean}{162}{\texttt{schrodinger\_gaussian\_fidelity\_witness}}, \bpsource{GaussianTesting/Bosons/PhysicalGeneralDensity.lean}{235}{\texttt{schrodingerReducedState\_top\_eigenstate}}, \bpsource{GaussianTesting/Bosons/PhysicalGeneralDensity.lean}{181}{\texttt{HilbertDensity.IsDensityOperator.eigenstate\_rejection\_le}}, \bpsource{GaussianTesting/Bosons/PhysicalGeneralDensity.lean}{199}{\texttt{HilbertDensity.IsDensityOperator.eigenstate\_weight\_ge\_of\_overlap}}, \bpsource{GaussianTesting/Bosons/PhysicalGeneralRobustness.lean}{61}{\texttt{schrodingerReducedState\_diagonal\_expectation}}, \bpsource{GaussianTesting/Bosons/PhysicalGeneralRobustness.lean}{99}{\texttt{schrodingerCoherentAcceptance\_inverse\_symplectic}}, \bpsource{GaussianTesting/Bosons/PhysicalGeneralRobustness.lean}{137}{\texttt{schrodinger\_general\_fidelity\_witness\_of\_reduced\_weight}}, \bpsource{GaussianTesting/Bosons/PhysicalGeneralRobustness.lean}{81}{\texttt{schrodingerReducedState\_covariance}}.}

\textbf{Theorem node.} \bplink{thm:main-boson}{Fidelity bound for general bosonic Gaussian states}.

\subsection{Reduced-state top eigenvalue (Eq. 6.8)}\label{paper:boson-top-eigenvalue}
\textbf{Checked.} Acceptance at least 1−tau implies that the largest reduced-state eigenvalue lambda is at least 1−2 tau. \bpsource{material/manuscript.tex}{1214}{Manuscript, line 1214.}

For a unit input in the tau≤1/48 branch of Theorem 6.7, the following bound is established. Theorem 6.7 applies Lemma 6.2 to the actual reduced density operator. The public top-eigenstate theorem provides a unit eigenvector, its positive eigenvalue, the lower bound, and the rejection estimate used next. This is a proof step of Theorem 6.7.

\bpdecls{\textbf{Lean.} \bpsource{GaussianTesting/Bosons/PhysicalGeneralDensity.lean}{235}{\texttt{schrodingerReducedState\_top\_eigenstate}}.}

\textbf{Blueprint.} \bplink{thm:main-boson}{Fidelity bound for general bosonic Gaussian states}.

\subsection{Centered Gaussian approximation of the top eigenstate (Eq. 6.9)}\label{paper:boson-centered-approximation}
\textbf{Checked.} The reduced-state top eigenvector has squared overlap at least 1−8 tau/lambda² with a centered Gaussian state. \bpsource{material/manuscript.tex}{1223}{Manuscript, line 1223.}

For a unit input in the tau≤1/48 branch of Theorem 6.7, the following bound is established. Theorem 6.7 applies the checked Theorem 6.4 witness to the top eigenstate after the positive-residual estimate gives p0(xi)≥1−tau/lambda². The centered Gaussian belongs to the independent symplectic-vacuum orbit, so its overlap is the displayed vacuum matrix coefficient. This application is part of the existing general-fidelity witness proof.

\bpdecls{\textbf{Lean.} \bpsource{GaussianTesting/Bosons/PhysicalMeanZeroRobustness.lean}{44}{\texttt{schrodinger\_meanZero\_fidelity\_witness}}, \bpsource{GaussianTesting/Bosons/PhysicalGeneralRobustness.lean}{137}{\texttt{schrodinger\_general\_fidelity\_witness\_of\_reduced\_weight}}, \bpsource{GaussianTesting/Bosons/PhysicalGeneralRobustness.lean}{162}{\texttt{schrodinger\_gaussian\_fidelity\_witness}}.}

\textbf{Blueprint.} \bplink{thm:main-boson}{Fidelity bound for general bosonic Gaussian states}.

\subsection{Corollary 6.8}\label{paper:r6-8}
\textbf{Checked.} The general, mean-zero and coherent tests have perfect completeness and O(epsilon\textasciicircum{}(-2) log(1/delta)) copies for trace-distance threshold epsilon, with respectively four, two and two copies per round. \bpsource{material/manuscript.tex}{1254}{Manuscript, line 1254.}

The explicit independent product Bernoulli protocol uses the three physical Born acceptance probabilities and accepts iff every round accepts. For normalized wavefunctions and epsilon,delta in (0,1), it has perfect completeness and rejection probability at least 1−delta whenever physical squared class fidelity is at most 1−epsilon². It uses exactly k*ceil(C/epsilon² log(1/delta)) copies, with (k,C)=(4,48),(2,8),(2,4). The ceiling bound includes +k uniformly; for delta≤exp(−1) the checked bound is 2kC/epsilon² log(1/delta), giving the stated asymptotic rate. The physical readouts are checked in Figure 1(b–d).

\textbf{Proof route.} Following the manuscript, physical robustness gives one-round rejection at least epsilon²/C, while the certainty propositions give perfect completeness. Independent fresh-copy outcomes have the product Bernoulli law, so all-round acceptance is exactly p\textasciicircum{}n. The exponential inequality (1−epsilon²/C)\textasciicircum{}n≤exp(−n epsilon²/C) gives the soundness bound. Taking the natural ceiling fixes the integer number of rounds and gives the exact total copy count and explicit asymptotic estimate.

\bpdecls{\textbf{Lean.} \bpsource{GaussianTesting/Bosons/PhysicalTesting.lean}{109}{\texttt{bosonic\_physical\_testing}}, \bpsource{GaussianTesting/Bosons/PhysicalTesting.lean}{100}{\texttt{bosonicTestProtocol}}, \bpsource{GaussianTesting/Bosons/PhysicalTesting.lean}{34}{\texttt{BosonicTestKind.copiesPerRound}}, \bpsource{GaussianTesting/Bosons/PhysicalTesting.lean}{90}{\texttt{BosonicTestKind.acceptance\_le\_of\_far}}, \bpsource{GaussianTesting/RepeatedTest.lean}{43}{\texttt{repeatedTest\_acceptance}}, \bpsource{GaussianTesting/RepeatedTest.lean}{59}{\texttt{repeatedTest\_rejection}}, \bpsource{GaussianTesting/RepeatedTest.lean}{97}{\texttt{testingRounds\_sufficient}}, \bpsource{GaussianTesting/RepeatedTest.lean}{101}{\texttt{testingCopies\_le}}, \bpsource{GaussianTesting/RepeatedTest.lean}{114}{\texttt{testingCopies\_le\_rate}}, \bpsource{GaussianTesting/Bosons/Main.lean}{1211}{\texttt{boson\_mean\_zero\_testing}}, \bpsource{GaussianTesting/Bosons/Coherent.lean}{1030}{\texttt{boson\_coherent\_testing}}.}

\textbf{Blueprint.} \bplink{cor:boson-testing}{Testing bosonic Gaussian and coherent states}, \bplink{target:measurement-endpoints}{Testing Gaussianity with constant sample complexity}.
